\documentclass[11pt]{article}

\usepackage[T1]{fontenc}
\usepackage[utf8]{inputenc}
\usepackage[english]{babel}
\usepackage[margin=1in]{geometry}
\usepackage{lmodern}
\usepackage{microtype}
\usepackage{mathtools}
\usepackage{amsmath,amssymb,amsthm}
\usepackage{array}
\usepackage{booktabs}
\usepackage{tabularx}
\usepackage{longtable}
\usepackage{float}
\usepackage{graphicx}
\usepackage{pdflscape}
\usepackage{caption}
\usepackage{enumitem}
\usepackage{fancyhdr}
\usepackage[numbers]{natbib}
\usepackage{doi}
\usepackage{xurl}
\usepackage{hyperref}
\usepackage{cleveref}
\usepackage{orcidlink}
\usepackage{titling}



\newcommand{\PaperOne}{Paper~1}
\newcommand{\PaperOneTitle}{Hiddenness and the Structural Law of Emergence}
\newcommand{\PaperTwo}{Paper~2}

\newcommand{\Tsat}{T^{!}_{\mathrm{SAT}}}
\newcommand{\Saud}{S_{\mathrm{aud}}^{\mathrm{SAT}}}
\newcommand{\Hsat}{H^{!}_{\mathrm{SAT}}}
\newcommand{\HsatN}[1]{H^{!}_{\mathrm{SAT}}(#1)}

\newcommand{\Ipoly}{I_{\mathrm{poly}}}
\newcommand{\EQbang}{E^{!}_{Q}}
\newcommand{\EMbang}{E^{!}_{M}}

\newcommand{\Qsat}{Q^{!}_{\mathrm{SAT}}}
\newcommand{\Msat}{M^{!}_{\mathrm{SAT}}}
\newcommand{\ComparisonMap}{\pi^{!}}

\newcommand{\OmegaSAT}{\Omega_{\mathrm{SAT}}}
\newcommand{\VisBang}{\operatorname{Vis}^{!}}

\newcommand{\SATbit}{\operatorname{SATbit}}
\newcommand{\omegaSATsel}{\omega_{\mathrm{SAT}}}   

\newcommand{\CurrentizeDec}{\operatorname{Currentize}^{\mathrm{dec}}_{Q}}
\newcommand{\CurrentizeWit}{\operatorname{Currentize}^{\mathrm{wit}}_{Q}}

\newcommand{\CSLSATHidden}{\Gamma_{\textsc{csl-sat-hidden}}}

\newcommand{\BIJudgment}[4]{#1\,;\,#2\,;\,#3\,\vdash\,#4}

\newcommand{\PSB}{P^{\mathrm{SB}}}
\newcommand{\NPSB}{NP^{\mathrm{SB}}}

\newcommand{\GateOne}{G_{1}}                \newcommand{\GateTwo}{G_{2}}                \newcommand{\GateThree}{G_{3}}              \newcommand{\GateFour}{G_{4}}               \newcommand{\GateFive}{G_{5}}               \newcommand{\GateSix}{G_{6}}                \newcommand{\GateSeven}{G_{7}}              

\newcommand{\FATCD}{\mathsf{FATCD}}
\newcommand{\BirdInt}{\mathbf{BirdInt}^{\mathrm{aud}}_{\mathrm{fin}}}

\newcommand{\RecognitionGrade}{\emph{recognition grade}}
\newcommand{\TheoremGrade}{\emph{theorem grade}}

\newcommand{\CSL}{\textsc{csl}}
\newcommand{\CSLcore}{Candidate Structural Law}
\newcommand{\VDiff}{V-Differential classification}
 
\newtheorem{theorem}{Theorem}[section]
\newtheorem{lemma}[theorem]{Lemma}
\newtheorem{proposition}[theorem]{Proposition}
\newtheorem{corollary}[theorem]{Corollary}
\theoremstyle{definition}
\newtheorem{definition}[theorem]{Definition}
\newtheorem{remark}[theorem]{Remark}
\floatstyle{ruled}
\newfloat{algorithm}{tbp}{loa}
\floatname{algorithm}{Algorithm}

\urlstyle{same}
\captionsetup{
  font=small,
  labelfont=bf,
  labelsep=period
}
\captionsetup[table]{skip=6pt,position=top}
\captionsetup[figure]{skip=6pt}
\captionsetup[algorithm]{skip=6pt,position=top}
\crefname{algorithm}{algorithm}{algorithms}
\Crefname{algorithm}{Algorithm}{Algorithms}
\crefname{theorem}{theorem}{theorems}
\Crefname{theorem}{Theorem}{Theorems}
\crefname{lemma}{lemma}{lemmas}
\Crefname{lemma}{Lemma}{Lemmas}
\crefname{proposition}{proposition}{propositions}
\Crefname{proposition}{Proposition}{Propositions}
\crefname{corollary}{corollary}{corollaries}
\Crefname{corollary}{Corollary}{Corollaries}
\crefname{definition}{definition}{definitions}
\Crefname{definition}{Definition}{Definitions}
\crefname{remark}{remark}{remarks}
\Crefname{remark}{Remark}{Remarks}
\setlist[enumerate]{leftmargin=*, itemsep=2pt, topsep=4pt}
\setlength{\droptitle}{-3.2em}
\pretitle{\begin{center}\LARGE\bfseries}
\posttitle{\par\end{center}\vspace{0.5em}}
\preauthor{\begin{center}\normalsize}
\postauthor{\par\end{center}\vspace{-0.4em}}
\predate{\begin{center}\small}
\postdate{\par\end{center}\vspace{0.6em}}
\setlength{\emergencystretch}{2em}
\clubpenalty=10000
\widowpenalty=10000
\displaywidowpenalty=10000
\interfootnotelinepenalty=10000
\allowdisplaybreaks[2]
\fancypagestyle{firstpage}{\fancyhf{}
  \fancyfoot[C]{\scriptsize Automorph Inc., Wilmington, DE, USA \quad \texttt{ioannis@automorph.io} \quad \textcopyright\ 2026 \quad CC-BY 4.0}
  \renewcommand{\headrulewidth}{0pt}
  \renewcommand{\footrulewidth}{0pt}
}

\title{P $\neq$ NP under Closure of the Saturated SAT Layer}
\author{Ioannis Tsiokos\,\orcidlink{0009-0009-7659-5964}}
\date{\today}
\hypersetup{
  pdftitle={P != NP under Closure of the Saturated SAT Layer},
  pdfauthor={Ioannis Tsiokos},
  hidelinks
}

\begin{document}

\maketitle
\thispagestyle{firstpage}

\begin{abstract}
This paper proves a conditional separation between \(P\) and \(NP\)
inside Six Birds Theory, a typed framework for closure formation in
mathematical structures.  For SAT, the framework supplies a saturated
SAT layer \(\Tsat\) --- a typed carrier of polynomial-time SAT
computations together with their audit data --- observed through a
lawful polynomial-time instrument \(\Ipoly\) that declares which
observations on the layer count as \emph{current}.  Two translation
theorems connect this setup to standard SAT semantics.  Translation
Theorem A identifies \(\mathrm{SAT}\in P\) with decision
currentization, the assertion that the SAT decision bit is itself a
current observation on \(\Tsat\).  Translation Theorem B identifies
decision currentization with currentization of the canonical
witness producer \(\omega_{\mathrm{SAT}}^{\mathrm{sel}}\) of the
autonomous lift \(\OmegaSAT\) --- the lexicographically-first
satisfying assignment on satisfiable instances and \(\bot\) otherwise
--- using SAT self-reducibility and the closure of the
current-observable family under polynomial-time post-processing,
finite products, and composition.  The structural hypothesis
\(\CSLSATHidden\) of \citet{TsiokosHiddenness2026} asserts that this
canonical witness producer is not available as a lawful current
observable on \(\Tsat\).  Under that
hypothesis, the four-step contradiction chain gives
\(\mathrm{SAT}\notin P\); combined with \(\mathrm{SAT}\in NP\) from
the standard verifier package, the framework's typed judgment
discipline accepts the separation \(\PSB\neq\NPSB\) at recognition
grade --- a status reserved for claims that follow from a named
structural hypothesis rather than from framework primitives alone.
Read outside the framework, this is the conditional theorem: closure
of the formed SAT layer under \(\Ipoly\) implies \(P\neq NP\).  Read
in the audit-closure register of \citet{Tsiokos2026AOR}, the same
closure records tag \(\Tsat\) as a carrier-attached cascade-stable
AOR package under \(\CSLSATHidden\) --- a structural-witness
packaging of the same conditional separation.  The hypothesis \(\CSLSATHidden\) further
localizes to a precise residual: the canonical lexicographic
SAT-branching readout is not a lawful current observable on the
standard fixed-presentation polynomial current carrier.  A Class-B
Irremovability theorem then rules out one route to removing the
hypothesis: no basis-generated Class-B carrier-feature method can
remove the residual, except by leaving the basis or importing a
target-equivalent lower bound.  The
barrier is analogous in role to relativization, natural proofs, and
algebrization, but methodologically distinct and internal to Six
Birds.  The paper does not claim an unconditional standard-ZFC separation of
\(P\) and \(NP\), derivability of \(\CSLSATHidden\) from framework
primitives, circumvention of the classical barriers in their native
senses, absolute irremovability of \(\CSLSATHidden\), a barrier
against methods using carrier structure outside the standard basis,
or any AOR-membership claim outside the stated formed-SAT,
standard-basis scope.
\end{abstract}

\medskip
\noindent\textbf{Keywords.} Six Birds Theory; saturated observable family; decision currentization; witness currentization; SAT self-reducibility; recognition grade; BirdInt judgment; conditional P versus NP; Class-B irremovability barrier; fixed-presentation polynomial carrier; fourth methodological barrier; audit closure.



\section{Introduction}
\label{sec:intro}

Does \(P\) equal \(NP\)?  The classical question asks whether every
language whose yes-instances have efficiently checkable witnesses also
has an efficient decision procedure.  It is a question about standard
language classes, and any proposal that touches it must say carefully
which part of the standard problem is being addressed.  Decades of
work have established structural barriers --- natural proofs
\citep{RazborovRudich1997NaturalProofs}, algebrization
\citep{AaronsonWigderson2009Algebrization}, and relativization
\citep{BakerGillSolovay1975Relativization} --- that constrain any
direct proof of separation.  The present paper does not attempt to overcome these
barriers in their native settings.  It develops a conditional theorem
in a typed structural calculus, and states exactly what that theorem
says when read from outside the calculus.  Alongside the three
classical barriers, it also proves a fourth, framework-internal barrier:
a Six Birds Class-B Irremovability theorem ruling out basis-generated
carrier-feature methods on the standard polynomial carrier.

The first result has two readings.  For a reader already working inside
Six Birds Theory, the paper proves a framework-relative judgment: under
the standard closure assumption of Foundations I, and conditional on a
specified structural recognition source, the internal predicates
\(\PSB\) and \(\NPSB\) separate at \RecognitionGrade{}
\citep{TsiokosFoundationsI2026}.  A \(\BirdInt\) judgment is the typed
audit-and-judgment shape supplied by Foundations III: it records the
source being used, the layer on which the source acts, the instrument
through which observation is made, and the proposition being accepted
\citep{TsiokosFoundationsIII2026}.  The grade matters.  At recognition
grade the acceptance depends on a structural source supplied to the
construction rather than derived within it.

For a complexity-theory reader, the same first result has a direct
reading.  The paper constructs a SAT-specific closure object and proves
that, if this formed SAT closure under the lawful polynomial instrument
carries the imported hiddenness source, then \(P\neq NP\).  The
antecedent is the structural recognition source \(\CSLSATHidden\),
supplied by \citet{TsiokosHiddenness2026}.  That paper develops the
\CSLcore{} and the \VDiff{} of substrates; the present paper uses the
SAT-specific source that those structures name, but it does not restate
the Candidate Structural Law or reprove the classification.  In this
reading, the paper's headline is the conditional theorem ``formed SAT
closure under \(\Ipoly\) implies \(P\neq NP\).''

The second result localizes the recognition source to the exact
fixed-branch residual
\(\mathrm{Br}^{\mathrm{lex}}_{\mathrm{SAT}}\notin \EQbang\).  In
ordinary language, the claim is that the canonical lexicographic
SAT-branching readout is not a lawful current observable on the
fixed-presentation carrier.  The localization identifies precisely
which residual the recognition source contributes to the contradiction
chain.

The third result is the Six Birds Class-B Irremovability barrier
theorem.  A Class-B law is a carrier-feature method: it tries to remove
the recognition source using the intrinsic structure of the carrier
rather than by changing the target.  A Class-B law is \emph{basis-generated}
when its feature is built from the native fixed-presentation basis by
the allowed closure operations.  The theorem says that no such law on
the standard fixed-presentation polynomial current carrier can remove
\(\CSLSATHidden\).  This explains why the conditional theorem lands at
recognition grade rather than upgrading further within Class-B methods.
The barrier is analogous in role to relativization, natural proofs, and
algebrization, but methodologically distinct from all three.

The construction begins by forming the saturated SAT layer
\(\Tsat\).  A saturated SAT layer is an aggregate bookkeeping object:
it contains an audited shell, a history carrier, current and
predictive observable families, a comparison map, and audit data.  The
audited shell \(\Saud\) is the finite-stage state record for the SAT
construction: it records the declared polynomial-time machine class,
the active SAT instances, the autonomous lift state, the visibility
lens, quotient packaging, audit ledger, and polynomial budget.  The
history carrier \(\Hsat\) is the all-input, machine-indexed collection
of SAT histories generated over these shells.  In ordinary language,
\(\Hsat\) is the place where every lawful polynomial-time SAT-handling
transducer is indexed, replayed, and audited.  The construction
instantiates the Foundations-II primitive-role schema \(P_{1}\)--\(P_{6}\):
substrate, lift, instrument, observable family, comparison map, and
audit are all declared as typed components rather than left as
implicit background choices \citep{TsiokosFoundationsII2026}.

Observation on this layer is controlled by the instrument \(\Ipoly\),
the lawful polynomial-time current instrument.  A current observable is
what a polynomial-time test may read at the present layer, using only
the encoded SAT instance and the declared \(Q\)-visible package data.
A predictive observable is allowed to speak about witness-side or
continuation-side behavior.  The saturated current-observable family
\(\EQbang\) admits every lawful deterministic polynomial-time
transducer over encoded SAT instances, provided it carries the audit
records required by the lawfulness discipline: a machine index, a time
certificate, an all-input replay convention, visibility bookkeeping,
and explicit no-smuggling records.  If a polynomial-time SAT decider
exists, it is admitted precisely because it is lawful computation; the
SAT bit is not inserted into the observable data by the framework.

From the same carrier the paper forms two quotients.  The current
quotient \(\Qsat\) identifies histories that all current observables
in \(\EQbang\) read the same way.  The predictive quotient \(\Msat\)
identifies histories only after the predictive observable family has
also been accounted for.  A comparison map descends from the identity
on histories, because predictive equality refines current equality.
The autonomous lift \(\OmegaSAT\) is the deterministic policy that
generates the SAT histories, their lift states, and their continuation
catalogs.  The \(\FATCD\) discipline of Foundations II supplies the
finite audited closure description, while \(\BirdInt\) supplies the
finite audited interaction calculus in which the final judgment is
recorded \citep{TsiokosFoundationsII2026,TsiokosFoundationsIII2026}.

The two translation theorems are the bridge between this typed
construction and the usual language-class statement.  Translation A
states
\[
  \mathrm{SAT}\in P
  \quad\Longleftrightarrow\quad
  \CurrentizeDec(\mathrm{SAT}).
\]
Decision currentization is the framework predicate saying that some
lawful current observable in \(\EQbang\) returns the SAT decision bit
on every encoded SAT instance.  The forward direction admits an
ordinary polynomial-time SAT decider into \(\EQbang\) by observable
saturation.  The reverse direction reads the lawfulness record of such
an observable and extracts an ordinary polynomial-time decider.

Translation B states
\[
  \CurrentizeDec(\mathrm{SAT})
  \quad\Longleftrightarrow\quad
  \CurrentizeWit(\OmegaSAT).
\]
Witness currentization is the corresponding predicate for producing a
witness: the observable returns bottom on unsatisfiable instances and
a satisfying assignment on satisfiable ones.  In the forward
direction, a SAT decider yields a witness producer by the standard
variable-by-variable self-reduction: call the decider on restrictions
of the formula and assemble a satisfying assignment when the original
instance is satisfiable.  The closure properties of \(\EQbang\) admit
the resulting wrapper, because polynomial-time post-processing,
finite products of fixed-stage observables, and polynomial-time
composition preserve the lawfulness records.  In the reverse
direction, a witness producer is post-processed to the decision bit by
testing whether its output is bottom.

The contradiction chain is then short.  Assume \(\mathrm{SAT}\in P\).
By Translation A, decision currentization holds.  By Translation B,
witness currentization holds for \(\OmegaSAT\).  The structural
recognition source \(\CSLSATHidden\), imported from
\citet{TsiokosHiddenness2026}, says that witness currentization fails
on the formed SAT layer under \(\Ipoly\).  These
statements contradict each other, so \(\mathrm{SAT}\notin P\).  Since
the standard verifier package gives \(\mathrm{SAT}\in NP\)
\citep{Cook1971TheoremProving}, the framework-relative conclusion is
that \(\PSB\neq\NPSB\) is accepted at \RecognitionGrade{} under the
standard Six Birds closure assumption.

The barrier theorem of \Cref{sec:class-b-barrier} closes one natural
route for upgrading this contradiction chain beyond recognition grade.

The source \(\CSLSATHidden\) is supplied, not derived.
\citet{TsiokosHiddenness2026} documents that the SAT-specific instance
of hiddenness is not obtained from the framework primitives by the
natural derivation routes; it is structural source content imported
into the present construction.  The \VDiff{} of that paper places PvNP
among the trace-only-exposing substrates: the available exposure
concerns trace state rather than witness content.  That placement is
consistent with the role played here by \(\CSLSATHidden\), but it is
not itself a derivation of the source.
Readers who do not adopt the framework's closure assumptions can read
the theorem conditionally, with the recognition source as the named
antecedent.

This is the scope of the paper.  It does not prove an unconditional
standard-ZFC separation of \(P\) from \(NP\).  Its headline judgment is
framework-relative and recognition-grade; outside the framework, what
is established is the conditional theorem just described.  Nor does
the paper claim to bypass natural proofs, algebrization, or
relativization in their native settings.  The construction works at a
different layer: typed observable saturation over a SAT-specific
history carrier.  Finally, the paper does not claim that
\(\CSLSATHidden\) follows from the framework primitives.  It is a
structural recognition source supplied by \citet{TsiokosHiddenness2026}.
The Class-B barrier theorem
also has its own scope: it does not prove absolute irremovability of the
recognition source, rule out future Six Birds laws using genuinely new
carrier structure outside the fixed-presentation basis, or apply to
enriched carriers such as objective SAT, branchwise carriers, or
presentation-quotiented carriers.  \Cref{sec:pvnp-aor-instance} then
reads the same closure records in the audit-closure register of
\citet{Tsiokos2026AOR}, exhibiting \(\Tsat\) as an audit-closed carrier
under \(\CSLSATHidden\) without adding new content.
 

\section{Saturated SAT layer}
\label{sec:saturated-closure}

We now build the saturated SAT layer that the rest of the paper
studies.  The construction has three parts.  First we define the
audited SAT shell, the finite-stage state record on which the
instrument acts.  Then we define the SAT history carrier, the
all-input, machine-indexed collection of histories generated over
those shells.  Finally we package the observable families, quotients,
comparison map, and audit closure into the saturated layer
\(\Tsat\).  The six-mechanism reading of the same construction, in
the Foundations-II primitive-role schema \(P_{1}\)--\(P_{6}\), is
recorded in \Cref{app:primitive-roles}; the present section gives the
construction itself and leaves that schema bookkeeping to the
appendix.

The word \emph{substrate} is used here in a plain mathematical sense:
it is the underlying typed state space on which instances, traces, and
audit data are formed.  In the present section the substrate is
specialized to SAT.  The raw inputs are encoded SAT instances, the
machine class consists of deterministic polynomial-time transducers
with declared certificates, and the records are arranged so that later
current and predictive observations have a common carrier on which to
act.

\begin{definition}[Audited SAT shell]
\label{def:pvnp:audited-shell}
The \emph{audited SAT shell} \(\Saud\) is the finite-stage typed
state record for the SAT construction.  A point of \(\Saud\) is a
tuple
\[
  s=(K_{\mathrm{poly}},\operatorname{Inst},\sigma_{L},\ell,
       \operatorname{pkg},\operatorname{Audit},\beta),
\]
where:
\begin{enumerate}
  \item \(K_{\mathrm{poly}}\) is the declared polynomial-time machine
  class, consisting of deterministic polynomial-time transducers
  together with time certificates;
  \item \(\operatorname{Inst}\) is the active SAT instance database,
  an indexed family of encoded SAT formulas under the current size
  filter;
  \item \(\sigma_{L}\) is the autonomous lift state, including phase
  and internal trace data;
  \item \(\ell\) is the current visibility lens;
  \item \(\operatorname{pkg}\) is the active quotient and
  continuation packaging;
  \item \(\operatorname{Audit}\) is the audit ledger;
  \item \(\beta\) is the polynomial budget and replay ledger.
\end{enumerate}
\footnote{The audited shell is realized in Lean by the audited SAT
shell declaration listed in \Cref{app:formalization}.}
\end{definition}

An encoded SAT instance here means the standard finite string
encoding of a propositional formula in conjunctive normal form
\citep{Cook1971TheoremProving,GareyJohnson1979Intractability},
together with its variable count and size data.  The polynomial-time
machine class is likewise meant in the standard deterministic sense
\citep{AroraBarak2009Complexity}: each admitted transducer acts on
encoded inputs and carries a declared polynomial budget in the input
length.  The visibility lens records
which projections are available to current observation; the quotient
packaging and audit ledger are defined later, where the current and
predictive quotients and the audit gates are introduced.

This seven-field shell is the SAT instance of the primitive-role
organization in Foundations II: machine update, admissibility,
timescale, lens selection, packaging, and audit are not implicit
ambient choices but named components of the typed object
\citep{TsiokosFoundationsII2026}.  The autonomous lift state is only
used here as a field of the shell; the lift itself is defined in
\Cref{sec:autonomous-lift}.  Similarly, the visibility lens anticipates
the lawful current instrument of \Cref{sec:observable-saturation}.

At a finite stage the shell restricts both sides of the construction:
instances are bounded by size, and machines are drawn from a declared
finite prefix of the indexed polynomial-time class together with their
time certificates.  This restriction is bookkeeping, not a change in
the problem.  It is the finite-audit way to speak about the usual SAT
input universe while keeping every object inspectable.  No solver
success, witness recovery, or decision correctness is built into the
shell.  The shell says what data are present before the lawfulness
discipline and the currentization predicates ask what any admitted
machine actually computes.

\begin{definition}[SAT history carrier]
\label{def:pvnp:history-carrier}
The \emph{SAT history carrier} \(\Hsat\) is the all-input,
machine-indexed collection of SAT histories generated over the audited
SAT shells.  At finite stage \(N\), a history in \(\HsatN{N}\) is a
finite typed trace
\[
  h=(\varphi,e,p_{e},I_{e},\operatorname{trace}_{Q},
     \sigma_{L},\operatorname{Cont}_{M},\operatorname{Audit}_{h}),
\]
where \(\varphi\) is an encoded SAT instance of size at most \(N\),
\(e\) is a deterministic polynomial-time transducer index, \(p_{e}\)
is its time certificate, \(I_{e}\) is the package interface,
\(\operatorname{trace}_{Q}\) is the current trace produced by the
autonomous run, \(\sigma_{L}\) is the autonomous lift state,
\(\operatorname{Cont}_{M}\) is the predictive continuation record, and
\(\operatorname{Audit}_{h}\) is the history-local audit record.
\footnote{The SAT history carrier is realized in Lean by the SAT
history carrier declaration listed in \Cref{app:formalization}, with
the finite-stage form exposed as a free argument; see
\Cref{app:formalization} for the narrowing-disclosure discipline.}
\end{definition}

The all-input condition means that histories are indexed over every
lawful pair consisting of an admitted machine and an admitted encoded
instance.  The machine-indexed condition means that the carrier
remembers which transducer produced the trace; histories are not
identified merely because their visible outputs agree.  This is
essential for the lawfulness discipline in the next section, where
machine identity, time certificate, input convention, output
convention, and replay record all remain part of the observable
audit.

A history therefore has more structure than the output of a machine.
It includes the input formula, the identity of the transducer, the
certificate under which the run is admitted, the visible trace
available at the current layer, the lift state generated during the
run, the predictive continuation record, and the local audit ledger.
These fields are deliberately separated.  A current observable may be
allowed to read the current trace and the declared \(Q\)-visible
package data, while a predictive observable may later refer to
continuation or witness-side records.  The separation of fields is
what lets the later quotients compare ``same for current observation''
against ``same for predictive observation'' without changing the
underlying histories.

The machine indexing is also important for the all-input semantics.
The carrier does not represent a single chosen algorithm for SAT.  It
represents the family of histories obtained by running every admitted
deterministic polynomial-time transducer on every admitted encoded SAT
instance, with the declared replay and audit data attached.  A
putative SAT decider, if one exists, appears as one indexed member of
the lawful machine class, subject to the same record requirements as
all other members.

Mathematically, the intended carrier is the directed colimit of the
finite-stage carriers \(\HsatN{N}\), as \(N\) ranges over the active
size and machine-prefix bounds.  The body uses the directed-colimit
reading; every construction needed for the translation theorems is
already visible at the finite stage on which the audit is performed.

\begin{definition}[Saturated SAT layer]
\label{def:pvnp:saturated-closure}
The \emph{saturated SAT layer} \(\Tsat\) is the aggregate typed
bookkeeping object
\[
  \Tsat
  =
  (\Saud,\Hsat,\EQbang,\EMbang,\Qsat,\Msat,\ComparisonMap,
    \operatorname{Audit}^{!}).
\]
Its components are the audited shell \(\Saud\), the SAT history
carrier \(\Hsat\), the saturated current and predictive observable
families \(\EQbang\) and \(\EMbang\), the saturated current and
predictive quotients \(\Qsat\) and \(\Msat\), the comparison map
\(\ComparisonMap\) between the two quotients, and the closed audit
record \(\operatorname{Audit}^{!}\) under the lawful polynomial-time
instrument \(\Ipoly\).
\footnote{The saturated SAT layer is realized in Lean by the saturated
SAT layer declaration listed in \Cref{app:formalization}; the Lean
record carries the typed core
\((\Hsat,\EQbang,\EMbang,\ComparisonMap,\operatorname{Audit}^{!})\)
plus closure-operator and admissibility/audit fields, while the
audited shell \(\Saud\) and saturated quotients \(\Qsat,\Msat\) live in
companion records.  The narrowing-disclosure discipline is recorded in
\Cref{app:formalization}.}
\end{definition}

This definition is intentionally by aggregation.  The saturated layer
does not hide an additional oracle or a new computational resource; it
bundles the typed components on which the later arguments depend.  The
observable families \(\EQbang\) and \(\EMbang\) are defined in
\Cref{sec:observable-saturation}, the quotients and comparison map in
\Cref{sec:quotients}, and the audit gates in \Cref{app:gate-audit}.
The layer \(\Tsat\) is the object over which the currentization
predicates and the two translation theorems are stated.

The saturation in this definition is saturation of the declared
observable and audit structure, not a claim that SAT has already been
decided.  The current family \(\EQbang\) will include all lawful
polynomial-time current observables once lawfulness has been defined;
the predictive family \(\EMbang\) will record the corresponding
continuation-sensitive observables.  The quotients \(\Qsat\) and
\(\Msat\) are then quotients of the same carrier \(\Hsat\).  This
common-carrier feature is what makes the comparison map meaningful:
the map compares what is forgotten by current observation against what
is retained by predictive observation, rather than comparing two
unrelated spaces.

One should also distinguish the layer from the recognition source used
later.  The present section only forms the typed object.  It does not
assert hiddenness, non-currentization, or any separation of complexity
classes.  Those claims enter later through the observable-saturation
definitions, the two translation theorems, and finally the structural
recognition source.  The role of \(\Tsat\) is to give all of those
later statements a fixed, audited SAT carrier.

The construction is therefore finitary and audit-friendly.  Every
component is a typed record with declared content: encoded instances,
machine indices, time certificates, traces, quotient packages,
visibility data, and audit ledgers.  Nothing in this section invokes
inputs beyond the framework primitives and the standard SAT encoding.
The audit and judgment shape used later is supplied by the finite
audited interaction calculus \(\BirdInt\) of Foundations III
\citep{TsiokosFoundationsIII2026}.  The next section gives
\(\EQbang\) its mathematical content by specifying which
polynomial-time transducers count as lawful current observables.
 

\section{Observable saturation}
\label{sec:observable-saturation}

The previous section formed the saturated SAT layer \(\Tsat\).  We now
specify what current observation means on that layer.  The lawful
polynomial-time current instrument \(\Ipoly\) fixes the machine class,
budget, input convention, visibility lens, and audit interface for
current observations.  The lawfulness discipline then says which
candidate observables are admissible.  Finally, the saturated
current-observable family \(\EQbang\) is defined to be exactly the
family of admissible lawful current observables.  The closure theorem
for \(\EQbang\), proved next in this section, records closure under
polynomial-time post-processing, finite products, and polynomial-time
composition.

A \emph{current observable} in this paper is first a candidate partial
finite-typed function from histories in \(\Hsat\) to an output
alphabet such as decision bits, witness strings, replay hashes, or
package data, realized by a declared polynomial-time transducer.  A
transducer is a deterministic Turing machine, or equivalent RAM
machine, that produces output rather than only accepting or rejecting.
Lawfulness will force the observable to be defined on every admitted
encoded SAT instance at the current stage.  Current observables may
read only the \(Q\)-visible projection of a history.  The visibility
lens is the projection map specifying which fields of the history are
exposed to the transducer.  A predictive observable is the same kind
of finite typed test, but with continuation-side or witness-side data
also made available.  This current-versus-predictive distinction is
the load-bearing distinction for the translation theorems.

\begin{definition}[Lawful polynomial-time current instrument]
\label{def:pvnp:i-poly}
The \emph{lawful polynomial-time current instrument} \(\Ipoly\) is the
typed instrument on \(\Tsat\) consisting of:
\begin{enumerate}
  \item the machine class admitted by the audited shell \(\Saud\),
  namely the deterministic polynomial-time transducers drawn from
  \(K_{\mathrm{poly}}\);
  \item the polynomial time budget under which those transducers are
  admitted;
  \item the input-encoding convention, namely the standard encoded SAT
  instances fixed in \Cref{sec:saturated-closure};
  \item the visibility lens selecting the \(Q\)-visible projection of
  histories in \(\Hsat\);
  \item the audit interface containing the declared no-smuggling
  commitments that any admitted transducer must satisfy.
\end{enumerate}
\end{definition}

Thus \(\Ipoly\) is not a single SAT algorithm.  It is the current
instrument under which polynomial-time machines are allowed to act on
the SAT history carrier.  Its budget component says what it means for
a run to be polynomial-time at the current stage.  Its visibility
component says what input data the run may see.  Its audit component
records the commitments later checked by the gate audit.  The
instrument is the point at which the finite-audit discipline of the
Foundations-II construction enters the current-observation side of the
SAT layer \citep{TsiokosFoundationsII2026}.

\begin{definition}[Lawfulness discipline]
\label{def:pvnp:lawfulness-discipline}
A candidate current observable on \(\Hsat\) is \emph{lawful} when it
satisfies the following five clauses.
\begin{enumerate}
  \item \emph{Machine indexing.}  The observable is realized by a
  declared deterministic polynomial-time transducer drawn from
  \(K_{\mathrm{poly}}\), with an associated machine index.
  \item \emph{Time certificate.}  The transducer carries a polynomial
  time certificate in the input length, witnessed by an explicit
  polynomial.
  \item \emph{All-input replay convention.}  The observable is defined
  on every encoded SAT instance admitted at the current stage, with a
  replay convention for the corresponding all-input transcript.
  \item \emph{Visibility bookkeeping.}  The observable reads only the
  declared \(Q\)-visible projection of the history, as selected by the
  visibility lens of \(\Ipoly\).
  \item \emph{No-smuggling records.}  The observable carries explicit
  audit records satisfying the gate predicates declared in the audit
  interface.
\end{enumerate}
\end{definition}

The five clauses are independent requirements, and an observable must
satisfy all five to be admitted as current.  Machine indexing prevents
an anonymous computation from entering the carrier.  The time
certificate records the polynomial bound rather than leaving it as an
informal assertion.  The all-input replay convention rules out a
machine that is only specified on a favorable subset of instances.
Visibility bookkeeping enforces the current layer by requiring the
transducer to read only the \(Q\)-visible data.  The no-smuggling
records are the typed commitments that the observable has not called a
SAT oracle, witness oracle, external verifier oracle, hidden SAT bit,
or other forbidden source \citep{AroraBarak2009Complexity}; the audit
predicates themselves are listed in \Cref{app:gate-audit}.

This definition also carries the saturation correction used throughout
the paper.  If a polynomial-time SAT decider exists, it is admissible
as a current observable precisely because it is lawful computation
under these five clauses.  Conversely, if a current observable is
admitted,
then its lawfulness record contains the transducer, its index, its
time certificate, its input and output conventions, its replay data,
and its visibility record.

\begin{definition}[Saturated current-observable family]
\label{def:pvnp:e-q-saturation}
The \emph{saturated current-observable family} \(\EQbang\) is the
family of all lawful polynomial-time current observables on the SAT
history carrier:
\[
  \EQbang
  =
  \{\, f:\Hsat\to Y
      \mid
      f \text{ is a lawful polynomial-time current observable} \,\},
\]
where \(Y\) ranges over the finite typed output alphabets admitted by
the instrument.
\end{definition}

Membership in \(\EQbang\) is exactly lawfulness.  In other words,
\(f\in\EQbang\) if and only if \(f\) satisfies the five-part
lawfulness discipline just stated.  There is no additional filter
after lawfulness: saturation means that every lawful polynomial-time
transducer is admitted as a current observable.

For a history \(h\in\Hsat\) whose SAT instance is \(\varphi\), the
event associated to a lawful transducer \(T\) records the machine
identity, the computed value \(T(\varphi)\), the replay data for the
run, the current-level tag, and the instrument \(\Ipoly\).  The exact
codomain is not important here; what matters is that the value is
typed, replayable, and audited.  The observable is current because it
is evaluated only through the \(Q\)-visible projection of the history.

This saturation is essential for the translation theorems.  If a
polynomial-time SAT decider exists, then after its machine index, time
certificate, all-input replay convention, visibility bookkeeping, and
audit records are declared, it belongs to \(\EQbang\) by saturation.
Conversely, any member of \(\EQbang\) is realized by such a
deterministic polynomial-time transducer, because membership is
lawfulness.  This two-way relationship is the bridge between the
framework-internal decision-currentization predicate and standard
membership in \(P\), formalized later as Translation A.

We now record the three closure properties needed downstream.  The
family \(\EQbang\) is closed under polynomial-time post-processing,
finite products of fixed-stage observables, and polynomial-time
composition \citep{Goldreich2008ConceptualComplexity,AroraBarak2009Complexity}.  The theorem is stated as a single three-clause closure
statement.  The proof is by exhibiting the corresponding lawful
transducer wrappers and checking that the time, visibility, replay,
and audit records remain lawful.

\begin{theorem}[Closure properties of \(\EQbang\)]
\label{thm:pvnp:e-q-closure-properties}
The saturated current-observable family \(\EQbang\) has the following
closure properties.
\begin{enumerate}
  \item \emph{Post-processing.}  If \(f\in\EQbang\) has codomain
  \(Y\), and \(g:Y\to Z\) is a deterministic polynomial-time function
  with a declared polynomial budget in the input length, then
  \(g\circ f\in\EQbang\).

  \item \emph{Finite products.}  If
  \(f_{1},\ldots,f_{k}\in\EQbang\) have codomains
  \(Y_{1},\ldots,Y_{k}\), and the family is fixed-stage in the audit
  sense, at a common stage \(N\), then the product observable
  \[
    h\longmapsto (f_{1}(h),\ldots,f_{k}(h))
  \]
  belongs to \(\EQbang\), with codomain
  \(Y_{1}\times\cdots\times Y_{k}\).

  \item \emph{Composition.}  Suppose \(f\in\EQbang\) has codomain
  \(Y\), the values \(f(h)\) are encoded SAT instances or other
  admitted input encodings, and \(g\in\EQbang\) is a current
  observable defined on histories of the resulting instances.  Then
  the suitable polynomial-time composition \(g\star f\) belongs to
  \(\EQbang\).  This composition is through the lawful transducer
  interfaces, not through identity on the history carrier: the
  realizing transducer is the composition of the two underlying
  transducers, with combined time, visibility, replay, and audit
  records.
\end{enumerate}
\end{theorem}

\begin{proof}
Membership in \(\EQbang\) is membership in the lawful
polynomial-time current-observable family.  It is therefore enough, in
each case, to exhibit a realizing transducer and check the five
lawfulness clauses of \Cref{def:pvnp:lawfulness-discipline}.

For post-processing, let \(f\in\EQbang\) be realized by a lawful
transducer \(T_f\) with polynomial time certificate \(p_f\), and let
\(g:Y\to Z\) be realized by a deterministic polynomial-time
post-processing transducer with budget \(p_g\).  Define the wrapper
\(W\) for \(g\circ f\) as follows: on an admitted SAT history \(h\),
run \(T_f\) on the \(Q\)-visible input allowed by \(\Ipoly\), obtain
the output \(T_f(h)\in Y\), and then run \(g\) on that output.  The
machine-indexing clause holds because \(W\) is the declared composite
of the index for \(T_f\) with the index for \(g\).  If the input
length is \(n\), then \(T_f\) runs in time \(p_f(n)\), its output has
size polynomially bounded by \(p_f(n)\), and the second stage runs in
time at most \(p_g(p_f(n))\) after replacing \(p_f\) by a polynomial
output-size bound if needed.  Hence
\[
  p_W(n)=p_f(n)+p_g(p_f(n))
\]
is a polynomial time certificate for \(W\).  The all-input replay
convention is preserved because \(T_f\) is defined on every admitted
SAT instance and \(g\) is defined on every output in its declared
codomain.  Visibility bookkeeping is unchanged: \(W\) reads no part of
the history beyond what \(T_f\) reads, and \(T_f\) reads only the
\(Q\)-visible projection.  Finally, the no-smuggling records propagate
because \(g\) is an ordinary polynomial-time post-processing
computation and introduces no oracle, hidden witness source, or
external verifier.  Thus \(g\circ f\) is lawful and belongs to
\(\EQbang\).

For finite products, let \(f_1,\ldots,f_k\in\EQbang\) be fixed at a
common audit stage \(N\), and let \(T_{f_i}\) be lawful realizing
transducers with polynomial certificates \(p_{f_i}\).  The product
wrapper \(W^{\mathrm{prod}}\) runs the \(T_{f_i}\) sequentially on the
same admitted input history, or equivalently by a standard parallel
simulation, and returns
\[
  (f_1(h),\ldots,f_k(h)).
\]
Its machine index is the finite tuple of the individual machine
indices together with the fixed product wrapper.  In the sequential
presentation the time certificate is
\[
  p_{\mathrm{prod}}(n)=\sum_{i=1}^{k}p_{f_i}(n),
\]
which is polynomial because \(k\) is fixed.  The all-input replay
condition holds because each \(f_i\) is defined on every admitted
instance at stage \(N\).  The visibility record is the finite union of
the individual visibility records; since each component reads only the
\(Q\)-visible projection, the product wrapper also reads only that
projection.  The audit and no-smuggling records are the product of the
individual records, with no new computational source introduced by
forming the tuple.  Hence the finite product observable is lawful and
belongs to \(\EQbang\).

For composition, let \(f\in\EQbang\) be realized by \(T_f\), and
suppose that \(f(h)\) is an admitted input encoding, for instance an
encoded SAT instance \(\psi\).  Let \(g\in\EQbang\) be realized by a
lawful transducer \(T_g\) on histories of such admitted instances.
The composition wrapper \(W^{\mathrm{comp}}\) first runs \(T_f\) on
the input history \(h_\varphi\) and obtains \(\psi\).  It then uses
the all-input replay convention of the carrier to package the
corresponding admitted history \(h_\psi\), and runs \(T_g\) on
\(h_\psi\).  The machine index is the declared composite of the
indices for \(T_f\), the replay-packaging step, and \(T_g\).  If
\(|h_\varphi|=n\), then \(T_f\) has polynomial runtime and polynomially
bounded output size, and the replay package \(h_\psi\) has polynomial
size in \(n\).  Applying the polynomial budget for \(T_g\) to that
polynomially bounded input size gives a polynomial combined budget.
The all-input convention holds because \(T_f\) outputs only admitted
encodings in this clause and \(T_g\) is defined on histories of all
admitted resulting instances.  Visibility and replay records compose
through the lawful transducer interfaces: \(T_f\) sees only the
\(Q\)-visible part of the original history, and \(T_g\) sees only the
\(Q\)-visible part of the packaged resulting history.  The audit and
no-smuggling records compose for the same reason as above: the wrapper
uses only the two declared lawful computations and the replay package
allowed by the carrier.  Thus \(g\star f\) is lawful and belongs to
\(\EQbang\).  This proves all three closure clauses.\footnote{This
theorem is verified in Lean by the closure-properties theorem
declaration listed in \Cref{app:formalization}; the Lean theorem covers
the admission-level closure under typed operation-record contracts for
post-processing, finite products, and composition.  The
narrowing-disclosure discipline is recorded in
\Cref{app:formalization}.}
\end{proof}
 

\section{Quotients and comparison}
\label{sec:quotients}

The observable families now determine two quotients of the history
carrier \(\Hsat\).  The saturated current-observable family
\(\EQbang\) gives the current quotient \(\Qsat\), while the predictive
observable family \(\EMbang\) gives the predictive quotient \(\Msat\).
The comparison map \(\ComparisonMap\) then forgets predictive
distinctions and lands in the current quotient.  These are the
SAT-specific analogs of the current and predictive quotients in the
general framework of \citet{TsiokosHiddenness2026}; here we give the
explicit setoid definitions on \(\Hsat\).

A \emph{setoid} is simply a named equivalence relation
\citep{MartinLof1984IntuitionisticTypeTheory}, used so that the
quotient set can be formed explicitly.  The equivalence relations
below are observation equivalences.  Two histories are current-
equivalent exactly when no current observable distinguishes them.  Two
histories are predictive-equivalent exactly when no predictive
observable, including the admissible continuation-side tests,
distinguishes them.  This is the standard quotient-by-observation
construction.

The point of using quotients rather than raw histories is that the
history carrier contains more data than a current observer is allowed
to read.  It contains machine identity, replay information, lift
state, and audit fields, while a current observable sees only the
projection made available by \(\Ipoly\).  Quotienting by all current
observables therefore records the observational content of the current
layer, rather than the full internal trace carried by \(\Hsat\).

\begin{definition}[Saturated current quotient]
\label{def:pvnp:q-saturated}
Define the current setoid \(\sim_{Q}\) on \(\Hsat\) by
\[
  h_{1}\sim_{Q}h_{2}
  \quad\Longleftrightarrow\quad
  f(h_{1})=f(h_{2})
  \text{ for every } f\in\EQbang .
\]
The \emph{saturated current quotient} is
\[
  \Qsat=\Hsat/{\sim_{Q}} .
\]
\end{definition}

Thus \(\Qsat\) remembers exactly what the saturated current
observables can see.  If two histories have the same value under
every lawful current observable, they determine the same current
class.  If some observable in \(\EQbang\) separates them, they lie in
different current classes.  \Cref{thm:pvnp:e-q-closure-properties}
ensures that
\(\EQbang\) is stable under the standard observable operations used
later.  The equivalence relation itself is inherited from equality of
observable values: reflexivity comes from equality with oneself,
symmetry from symmetry of equality, and transitivity from transitivity
of equality.

This quotient is deliberately coarse.  It does not say that two
histories are the same computation, or that they have the same
predictive continuations.  It says only that every lawful
polynomial-time current test sees them identically.  In later
sections, decision currentization is phrased in this current layer:
the question is whether the SAT decision bit is already represented by
some member of \(\EQbang\).

\begin{definition}[Saturated predictive quotient]
\label{def:pvnp:m-saturated}
Let \(\EMbang\) be the predictive-side analog of \(\EQbang\): the
family of lawful polynomial-time observables admitted under the
predictive visibility lens, which may read the continuation-side or
witness-side projections of SAT histories.  Define the predictive
setoid \(\sim_{M}\) on \(\Hsat\) by
\[
  h_{1}\sim_{M}h_{2}
  \quad\Longleftrightarrow\quad
  F(h_{1})=F(h_{2})
  \text{ for every } F\in\EMbang .
\]
The \emph{saturated predictive quotient} is
\[
  \Msat=\Hsat/{\sim_{M}} .
\]
\end{definition}

The family \(\EMbang\) is not unpacked again here.  It is the
predictive version of the lawfulness discipline already given for
\(\EQbang\): the same kind of polynomial-time transducer record, time
certificate, replay convention, and audit interface are present, but
the visibility lens is widened to include predictive continuation and
witness-side data.  Consequently \(\Msat\) can distinguish histories
that \(\Qsat\) identifies, exactly when the distinction is invisible
to all current observables but visible to a lawful predictive
observable.

The predictive quotient is correspondingly finer whenever witness or
continuation data matter.  It still uses lawful polynomial-time
observables; the difference is not computational power but visibility.
The same history carrier is being observed through a richer lens.
This common carrier is what makes a comparison map possible.

\begin{definition}[Comparison map]
\label{def:pvnp:comparison-map}
Every current observable is, in particular, a predictive observable
that ignores the continuation-side and witness-side fields.  Hence
the predictive setoid refines the current setoid:
\[
  h_{1}\sim_{M}h_{2}
  \quad\Longrightarrow\quad
  h_{1}\sim_{Q}h_{2}.
\]
The identity map on \(\Hsat\) therefore descends to a well-defined
function
\[
  \ComparisonMap:\Msat\longrightarrow\Qsat,
  \qquad
  \ComparisonMap([h]_{M})=[h]_{Q}.
\]
\end{definition}

The direction is important.  The map goes from predictive classes to
current classes, \(\Msat\to\Qsat\), because predictive equivalence is
finer: forgetting predictive distinctions gives a current equivalence
class.  This is the same direction as the forget-witness map in the
general quotient framework.  The comparison map is the technical
object through which current and predictive information are compared.
When it fails to be bijective, or when it has no polynomial-time
section compatible with the declared instrument, the saturated layer
contains predictive distinctions not recovered by current
observation.  The translation theorems analyze this situation on
SAT-encoded histories.

Thus \(\ComparisonMap\) is a bookkeeping map, but not a cosmetic one.
It records exactly which predictive distinctions survive after current
observation forgets continuation and witness-side structure.  The next
sections use the currentization predicates to ask whether decision
data and witness data become visible on the current side of this map.
 

\section{Autonomous lift on the SAT layer}
\label{sec:autonomous-lift}

We now define the two records that connect the SAT history carrier to
the current and predictive observable layers.  The autonomous lift
\(\OmegaSAT\) is the deterministic policy that generates SAT histories
and lift states.  The instrument-relative visibility map \(\VisBang\)
filters the lift-and-history data down to the projection exposed by
the chosen instrument.  \citet{TsiokosHiddenness2026} develops the
general autonomous-lift framework; here we use its SAT-specific
instance.

\begin{definition}[Autonomous lift \(\OmegaSAT\)]
\label{def:pvnp:omega-sat}
The \emph{autonomous lift} \(\OmegaSAT\) on the saturated SAT layer is
the deterministic-policy record consisting of:
\begin{enumerate}
  \item a SAT-specific lift-state space, whose elements encode the
  internal state of the autonomous policy as it generates SAT
  histories;
  \item an update rule sending a current lift state, active encoded
  SAT instance, and audited-shell bookkeeping data to the next lift
  state;
  \item a history-generation link sending lift-state trajectories to
  the corresponding histories in \(\Hsat\);
  \item a continuation catalog assigning to each lift state the
  admissible witness-side and predictive continuations available from
  that state.
\end{enumerate}
\footnote{The autonomous lift is realized in Lean by the autonomous
lift declaration listed in \Cref{app:formalization}, with the policy
components \((\OmegaSAT\)-state, update rule, history-generation map,
continuation catalog) exposed as independent fields together with
Prop-level no-oracle / no-hidden-witness commitments.  The
deterministic-policy linkage between the update rule and the
history-generation map is recorded paper-side.  See
\Cref{app:formalization} for the narrowing-disclosure discipline.}
\end{definition}

In standard terms, \(\OmegaSAT\) is a deterministic policy.  Its next
state is determined by its current state together with the SAT
instance and the shell data supplied by \(\Saud\).  It does not
consult an external oracle or a hidden witness.  The point of the lift
is to supply the predictive structure on which \(\EMbang\) can later
act: without a continuation catalog, there is no typed account of what
predictive observations are allowed to distinguish.

The lift is therefore not a solver and not a new source of
computational power.  It is the history-generation mechanism for the
saturated layer.  The histories it generates are still the audited,
machine-indexed histories of \(\Hsat\); the lift record says how those
histories are produced and which continuation data are attached to
their states.

\begin{definition}[Instrument-relative visibility map]
\label{def:pvnp:vis-bang}
The \emph{instrument-relative visibility map} \(\VisBang\) is the
typed map
\[
  \VisBang:
  (\OmegaSAT\text{-states})\times\Hsat
  \longrightarrow \mathrm{VisData},
\]
which sends a lift state and a SAT history to the visibility data
exposed by the instrument's selected lens.
\footnote{The instrument-relative visibility map is realized in Lean by
the instrument-relative visibility map declaration listed in
\Cref{app:formalization}.}
\end{definition}

Equivalently, \(\VisBang\) is the partial projection from the combined
lift-and-history record to the fields that the active instrument
allows an observable to read.  It is instrument-relative because the
projection depends on the lens selected by the instrument: current
only, current plus witness, or current plus continuation.  For the
saturated current quotient \(\Qsat\), the current-only projection of
\(\VisBang\) is used.  For the predictive quotient \(\Msat\), the
wider predictive projection is used.

Thus \(\OmegaSAT\) is the predictive engine, and \(\VisBang\) is the
visibility bridge.  The lift determines which continuations are
admissible from each state, and hence what \(\EMbang\) can ever
distinguish.  The visibility map determines what part of the same
lift-and-history data is exposed to current observation.  Together
they specify what the predictive engine generates and what current
observation can see of it.
 

\section{Currentization predicates}
\label{sec:currentize}

We now define the two currentization predicates used by the
translation theorems.  Decision currentization asks whether the SAT
decision bit is already represented by a lawful current observable in
\(\EQbang\).  Witness currentization asks whether a satisfying-witness
producer is already represented by such a current observable.  These
are the framework-internal counterparts of, respectively, standard
polynomial-time decision for SAT and the polynomial-time search
version of SAT.  The next two sections relate these predicates to the
standard complexity-theoretic statements.

Let
\[
  \SATbit:\{\text{encoded SAT instances}\}\longrightarrow\{0,1\}
\]
be the standard SAT decision function
\citep{Cook1971TheoremProving}: \(\SATbit(\varphi)=1\) exactly when
the encoded formula \(\varphi\) is satisfiable.  For satisfiable encoded instances, let
\[
  \omegaSATsel:
  \{\text{satisfiable encoded SAT instances}\}
  \longrightarrow
  \{\text{satisfying assignments}\}
\]
be the lexicographically first satisfying assignment under the
standard variable ordering
\citep{HemaspaandraOgihara2002SelfReducibility}.  The lex-first convention is arbitrary but
fixed.  Any explicit deterministic selector producing a satisfying
assignment whenever one exists would serve the same mathematical role;
the convention makes the witness target unambiguous.

The selector \(\omegaSATsel\) is defined only on satisfiable
instances.  To make a single observable with a total output convention,
the witness predicate below uses a disjoint symbol \(\bot\) on
unsatisfiable instances.  Thus the witness-output alphabet is the
disjoint union of satisfying assignments and this distinguished
failure value.  This is standard bookkeeping: it lets one current
observable be evaluated on every encoded SAT instance, while still
distinguishing a genuine satisfying assignment from the assertion that
none exists.

\begin{definition}[Decision currentization]
\label{def:pvnp:currentize-dec}
Decision currentization, written \(\CurrentizeDec(\mathrm{SAT})\),
holds when there exists a current observable \(f\in\EQbang\) such that
for every encoded SAT instance \(\varphi\),
\[
  f(h_{\varphi})=\SATbit(\varphi).
\]
Here \(h_{\varphi}\in\Hsat\) denotes the history of \(\varphi\) under
the fixed all-input replay convention for the lawful instrument
\(\Ipoly\).
\end{definition}

The choice of \(h_{\varphi}\) is only a bookkeeping choice.  The
all-input replay convention supplies a canonical way to view each
encoded instance as a history in the SAT carrier, for any fixed
machine index admitted by the replay convention.  The predicate above
requires agreement with \(\SATbit\) on every encoded SAT instance, so
the decision content is independent of that bookkeeping choice.  In
plain terms, \(\CurrentizeDec(\mathrm{SAT})\) says that the SAT
decision data are already available at the current layer.

Because \(f\) is required to belong to \(\EQbang\), this is not merely
the existence of an arbitrary Boolean-valued function.  The observable
must be realized by a lawful deterministic polynomial-time transducer,
with the machine index, time certificate, replay convention,
visibility record, and audit data required in
\Cref{sec:observable-saturation}.

\begin{definition}[Witness currentization]
\label{def:pvnp:currentize-wit}
Witness currentization, written \(\CurrentizeWit(\OmegaSAT)\), holds
when there exists a current observable \(g\in\EQbang\) such that for
every encoded SAT instance \(\varphi\),
\[
  g(h_{\varphi})=
  \begin{cases}
    \omegaSATsel(\varphi),
      & \text{if } \varphi \text{ is satisfiable},\\
    \bot,
      & \text{if } \varphi \text{ is unsatisfiable}.
  \end{cases}
\]
The symbol \(\bot\) is a designated output value meaning that no
satisfying assignment exists.
\end{definition}

Thus \(\CurrentizeWit(\OmegaSAT)\) says that the current layer already
contains a lawful polynomial-time observable returning a witness on
satisfiable instances and \(\bot\) on unsatisfiable instances.  In
this formulation the witness is the canonical value
\(\omegaSATsel(\varphi)\) when \(\varphi\) is satisfiable.  This makes
the predicate a functional statement, rather than the ambiguous claim
that the observable returns some satisfying assignment.

The two predicates differ in what they claim is current.  Decision
currentization claims that the SAT decision bit is current.  Witness
currentization claims that a canonical witness producer is current.
The first translation theorem brings standard SAT-in-\(P\) into this
framework language by identifying it with decision currentization.  The
second translation theorem connects decision and witness
currentization using the closure properties of \(\EQbang\).

Both predicates quantify over a single observable in \(\EQbang\).  The
all-input replay convention is what makes this meaningful: the same
lawful observable is evaluated on every encoded SAT instance, not on a
chosen subfamily.  The role of \(\omegaSATsel\) is only to make the
witness target deterministic.  The later self-reducibility argument
uses the standard polynomial relation between SAT instances and
satisfying assignments; it does not depend on any special property of
the lexicographic order beyond determinism.
 

\section{Translation Theorem A}
\label{sec:translation-a}

The first translation theorem identifies the framework-internal
decision-currentization predicate with the standard membership
predicate \(\mathrm{SAT}\in P\).  The forward direction uses observable
saturation: a standard polynomial-time SAT decider, once equipped with
its machine index, time certificate, replay convention, visibility
bookkeeping, and audit record, is a lawful current observable and
therefore belongs to \(\EQbang\).  The reverse direction uses
lawfulness extraction: a current observable computing the SAT decision
bit carries enough typed lawfulness data to recover an ordinary
polynomial-time SAT decider.  Both directions are proved below.

\begin{theorem}[Translation Theorem A]
\label{thm:translation-A}
For the saturated SAT layer and the current instrument \(\Ipoly\),
\[
  \mathrm{SAT}\in P
  \quad\Longleftrightarrow\quad
  \CurrentizeDec(\mathrm{SAT}).
\]
Equivalently, SAT has a standard polynomial-time decision procedure if
and only if there exists a lawful polynomial-time current observable in
\(\EQbang\) computing the SAT decision bit on every encoded SAT
instance.
\end{theorem}

\begin{proof}
Assume first that \(\mathrm{SAT}\in P\).  By the standard definition,
there is a deterministic polynomial-time Turing transducer \(T\) which,
on input an encoded SAT instance \(\varphi\), outputs the decision bit
\(\SATbit(\varphi)\).  Realize \(T\) as a candidate current observable
on the SAT history carrier \(\Hsat\) by evaluating it on the encoded
instance field of the canonical replay history \(h_\varphi\).

We check the five lawfulness clauses of
\Cref{def:pvnp:lawfulness-discipline}.  For machine indexing, assign
\(T\) a free machine index in the all-input replay convention.  For
the time certificate, use the polynomial budget witnessing
\(\mathrm{SAT}\in P\).  For the all-input replay convention, use the
fact that \(T\) is total on encoded SAT instances.  For visibility,
\(T\) reads only the encoded-instance field of \(h_\varphi\), which is
part of the \(Q\)-visible projection selected by \(\Ipoly\).  For the
no-smuggling record, \(T\) uses no SAT oracle, witness oracle, external
verifier oracle, hidden SAT bit, or other forbidden source; its audit
record is therefore the vacuous no-oracle audit.  Hence the induced
observable is lawful.  By the saturation definition of \(\EQbang\),
\Cref{def:pvnp:e-q-saturation}, it belongs to \(\EQbang\).  Its value
on \(h_\varphi\) is \(\SATbit(\varphi)\) for every encoded SAT
instance \(\varphi\), so \(\CurrentizeDec(\mathrm{SAT})\) holds.

Conversely, assume \(\CurrentizeDec(\mathrm{SAT})\).  By
\Cref{def:pvnp:currentize-dec}, there is an observable
\(f\in\EQbang\) such that
\[
  f(h_{\varphi})=\SATbit(\varphi)
\]
for every encoded SAT instance \(\varphi\).  By
\Cref{def:pvnp:e-q-saturation}, membership in \(\EQbang\) is exactly
lawfulness, so unpack the lawfulness record of \(f\).  The
machine-indexing clause gives a declared deterministic transducer
\(T_f\).  The time-certificate clause gives an explicit polynomial
budget \(p_f\).  The all-input replay clause gives totality on encoded
SAT instances.  The visibility-bookkeeping clause says that \(T_f\)
reads only the \(Q\)-visible projection, in particular the encoded
instance field available in the canonical history.  The audit record
supplies the corresponding no-smuggling commitments.

On every encoded instance \(\varphi\), the transducer \(T_f\) outputs
\(\SATbit(\varphi)\), because the observable \(f\) does so on
\(h_\varphi\).  Thus \(T_f\), with its polynomial certificate \(p_f\),
is a deterministic polynomial-time SAT decider.  Hence
\(\mathrm{SAT}\in P\).  This reverse implication is where the
lawfulness record is load-bearing: without that record, an abstract
current observable computing the right bit would not by itself expose
the underlying polynomial-time machine.  Here the record supplies
exactly the data needed to extract the standard decider.  The two
directions prove the equivalence.\footnote{This theorem is verified in
Lean by the Translation Theorem A declaration listed in
\Cref{app:formalization}; the Lean theorem proves the iff packaging
from the standard
\(\mathrm{SAT}\in P\)-to-current-observable bridge and the \(\EQbang\)
admission/extraction predicates taken as parametric hypotheses.  The
narrowing-disclosure discipline is recorded in
\Cref{app:formalization}.}
\end{proof}

The theorem is the structural translation between typed current
observables and the standard polynomial-time membership predicate.
 

\section{Translation Theorem B}
\label{sec:translation-b}

Translation Theorem B bridges decision currentization and witness
currentization, showing that the two predicates are equivalent: a
lawful current observable computing the SAT decision bit exists if
and only if a lawful current observable computing the canonical
satisfying-witness producer exists.  The
forward direction uses SAT self-reducibility
\citep{Schoening1986ComplexityStructure,HemaspaandraOgihara2002SelfReducibility,Selman1988SelfReducible},
namely the standard variable-by-variable witness-search-from-decider
construction, together with the closure properties of \(\EQbang\).  The reverse direction
post-processes the witness producer's output to recover the decision
bit.  We prove both directions below; the forward direction uses the
standard SAT self-reduction as a classical input.

\begin{theorem}[Translation Theorem B]
\label{thm:translation-B}
For the saturated SAT layer and autonomous lift \(\OmegaSAT\),
\[
  \CurrentizeDec(\mathrm{SAT})
  \quad\Longleftrightarrow\quad
  \CurrentizeWit(\OmegaSAT).
\]
Equivalently, the SAT decision bit is current if and only if the
canonical satisfying-witness producer is current.
\end{theorem}

\begin{proof}
We begin with the forward implication.  Assume
\(\CurrentizeDec(\mathrm{SAT})\).  By
\Cref{def:pvnp:currentize-dec}, there is an observable
\(f\in\EQbang\) such that
\[
  f(h_{\varphi})=\SATbit(\varphi)
\]
for every encoded SAT instance \(\varphi\).  Let \(T_{f}\) be the
lawful polynomial-time transducer realizing \(f\), as supplied by the
lawfulness record of \(\EQbang\).

We construct a witness-producing transducer \(T_{g}\).  Let
\(\varphi\) have variables \(x_{1},\ldots,x_{n}\).  For a partial
assignment \(\alpha\), write
\[
  \varphi_{\alpha\cup\{x_i\mapsto b\}}
\]
for the encoded formula obtained from \(\varphi\) by substituting the
assignments already chosen in \(\alpha\), setting \(x_i=b\), and then
re-encoding the resulting formula under the same SAT encoding
convention.  Simplification by propagation may be performed, but is
not essential; the important point is that restriction and
re-encoding are polynomial-time and produce admitted encoded SAT
instances.

On input \(\varphi\), the wrapper first evaluates \(f\) on the
canonical history \(h_{\varphi}\).  If the value is \(0\), it outputs
\(\bot\).  If the value is \(1\), it builds an assignment
variable-by-variable.  Start with the empty partial assignment
\(\alpha=\varnothing\).  At stage \(i\), form
\[
  \psi_i=\varphi_{\alpha\cup\{x_i\mapsto 0\}}.
\]
Evaluate \(f\) on the canonical replay history \(h_{\psi_i}\).  If
\(f(h_{\psi_i})=1\), set \(\alpha(x_i)=0\).  Otherwise set
\(\alpha(x_i)=1\).  This is the lexicographic convention: choose
\(0\) whenever it preserves satisfiability, and choose \(1\) only when
the \(0\)-branch is unsatisfiable.  After \(n\) stages, output the
assembled assignment \(\alpha\).

This is the standard SAT self-reduction
\citep{Schoening1986ComplexityStructure,HemaspaandraOgihara2002SelfReducibility,Selman1988SelfReducible}.
The invariant is that after each stage, the partial assignment
\(\alpha\) can be extended to a satisfying assignment of the original
formula whenever the original formula is satisfiable.  If the
\(0\)-branch remains satisfiable, the wrapper chooses it; otherwise
satisfiability of the current prefix forces the \(1\)-branch.  Thus,
when \(\varphi\) is satisfiable, the output assignment satisfies
\(\varphi\), and by the choice rule it is the lexicographically first
satisfying assignment \(\omegaSATsel(\varphi)\).  When \(\varphi\) is
unsatisfiable, the initial call to \(f\) returns \(0\), and the wrapper
outputs \(\bot\).

The runtime is polynomial.  The wrapper makes at most \(n+1\) calls to
the transducer \(T_{f}\), each on an encoded formula of size bounded
by a polynomial in \(|\varphi|\).  Formula restriction and
re-encoding are polynomial-time operations.  Since \(T_{f}\) carries a
polynomial time certificate, the combined runtime is bounded by a
polynomial in \(|\varphi|\).

It remains to check that the resulting observable belongs to
\(\EQbang\).  The wrapper \(T_g\) is a single deterministic
polynomial-time transducer.  On input \(\varphi\), it runs through
\(n+1\) sequential stages.  Stage \(0\) is the initial decision query.
At stage \(i\in\{1,\ldots,n\}\), the wrapper computes the restricted
formula
\[
  \psi_i=\varphi_{\alpha\cup\{x_i\mapsto 0\}}
\]
from the original formula and the current partial assignment
\(\alpha\), forms the canonical replay history \(h_{\psi_i}\), invokes
the lawful transducer \(T_f\) on that history, and updates \(\alpha\)
according to the returned decision bit.  The final stage outputs the
assembled assignment \(\alpha\), or \(\bot\) if the initial query
returned \(0\).

Each restriction-and-reencoding step is a polynomial-time input
transformation.  Each call to \(T_f\) is a polynomial-time composition
through the lawful transducer interface of \(f\), not an oracle call.
Each update of \(\alpha\), and the final assembly of the witness or
\(\bot\), is polynomial-time post-processing of the current transcript.
Thus the stage maps are obtained using the polynomial-time composition
and post-processing clauses of
\Cref{thm:pvnp:e-q-closure-properties}.  Iterating these stages
preserves polynomial-time lawfulness because the number of stages is
\(n+1\), polynomial in \(|\varphi|\), and the total time is bounded by
\[
  O\bigl((n+1)\cdot p_f(\operatorname{poly}(|\varphi|))\bigr)
  =
  \operatorname{poly}(|\varphi|).
\]

We now check the five lawfulness clauses for \(T_g\).  For machine
indexing, \(T_g\) is a single uniform deterministic transducer with a
declared machine index; the loop is part of the program, not an
external choice.  For the time certificate, the polynomial bound just
displayed supplies the declared certificate.  For all-input replay,
\(T_g\) is defined on every encoded SAT instance \(\varphi\), since
\(T_f\) is defined on every admitted encoded instance and each
\(\psi_i\) is again an admitted encoding.  For visibility bookkeeping,
\(T_g\) reads only the encoded-instance field of \(h_{\varphi}\); the
intermediate formulas \(\psi_i\) are computed internally by
polynomial-time restriction and re-encoding.  For no-smuggling,
\(T_g\) uses no oracle calls, hidden witness source, or external
verifier; the internal calls to \(T_f\) pass through \(T_f\)'s lawful
transducer interface, which already carries the no-oracle audit
record.  The audit record for \(T_g\) is the transcript-level union of
the \(T_f\) audit entries with the wrapper's own no-oracle commitment.
By saturation of \(\EQbang\) in \Cref{def:pvnp:e-q-saturation}, the
observable \(g\) realized by \(T_g\) belongs to \(\EQbang\).

Consequently \(g(h_{\varphi})=\omegaSATsel(\varphi)\) whenever
\(\varphi\) is satisfiable, and \(g(h_{\varphi})=\bot\) whenever
\(\varphi\) is unsatisfiable.  This is exactly the witness
currentization predicate \(\CurrentizeWit(\OmegaSAT)\).  The forward
direction uses only standard SAT self-reducibility and the closure
properties of \(\EQbang\) under post-processing and composition.

Conversely, assume \(\CurrentizeWit(\OmegaSAT)\).  By
\Cref{def:pvnp:currentize-wit}, there is an observable
\(g\in\EQbang\) such that for every encoded SAT instance
\(\varphi\),
\[
  g(h_{\varphi})=
  \begin{cases}
    \omegaSATsel(\varphi),
      & \text{if } \varphi \text{ is satisfiable},\\
    \bot,
      & \text{otherwise}.
  \end{cases}
\]
We construct \(f\in\EQbang\) with
\(f(h_{\varphi})=\SATbit(\varphi)\) for every encoded SAT instance
\(\varphi\).

Let
\[
  \pi:\{\text{witness outputs}\}\cup\{\bot\}\longrightarrow\{0,1\}
\]
be the deterministic post-processing function defined by
\(\pi(w)=1\) when \(w\neq\bot\), and \(\pi(\bot)=0\).  This function
is polynomial-time: it only checks whether the output token is the
distinguished failure value \(\bot\).  It need not verify the witness
against the formula, because the hypothesis
\(\CurrentizeWit(\OmegaSAT)\) already supplies the stronger promise
that \(g\) returns \(\omegaSATsel(\varphi)\) on satisfiable instances
and \(\bot\) on unsatisfiable instances.

Define \(f=\pi\circ g\).  If \(\varphi\) is satisfiable, then
\[
  g(h_{\varphi})=\omegaSATsel(\varphi)\neq\bot,
\]
so
\[
  f(h_{\varphi})
  =
  \pi(\omegaSATsel(\varphi))
  =
  1
  =
  \SATbit(\varphi).
\]
If \(\varphi\) is unsatisfiable, then \(g(h_{\varphi})=\bot\), so
\[
  f(h_{\varphi})
  =
  \pi(\bot)
  =
  0
  =
  \SATbit(\varphi).
\]
Thus \(f(h_{\varphi})=\SATbit(\varphi)\) for every encoded SAT
instance \(\varphi\).

By closure of \(\EQbang\) under polynomial-time post-processing,
\Cref{thm:pvnp:e-q-closure-properties}, the composite
\(f=\pi\circ g\) belongs to \(\EQbang\).  Therefore
\(\CurrentizeDec(\mathrm{SAT})\) holds.  The two implications prove
the equivalence.\footnote{This theorem is verified in Lean by the
Translation Theorem B
declaration listed in \Cref{app:formalization}; the Lean theorem proves
the iff packaging from the self-reduction wrapper, canonical-witness
correctness, closure admission, post-processor admission, and SAT-bit
correctness taken as parametric hypotheses.  The narrowing-disclosure
discipline is recorded in
\Cref{app:formalization}.}
\end{proof}
 

\section{Recognition source}
\label{sec:recognition-source}

We now identify the single structural input imported from the
hiddenness paper of \citet{TsiokosHiddenness2026}.  The recognition
source \(\CSLSATHidden\) is the SAT-specific instance of the \CSLcore{}
developed there: it is a structural content statement about what the
formed SAT layer exposes to current observation.  The present section
gives the conceptual content of that source and explains its
relationship to that paper.  The formal carrier schema, including the
typed fields used to attach the source to the saturated SAT layer, is
deferred to \Cref{app:recognition-source-carrier}.

\begin{definition}[Recognition source $\Gamma_{\mathrm{CSL\text{-}SAT\text{-}hidden}}$]
\label{def:pvnp:recognition-source-Gamma}
On the saturated SAT layer \(\Tsat\), under the lawful current
instrument \(\Ipoly\), the recognition source \(\CSLSATHidden\) is the
predicate
\[
  \CSLSATHidden
  \;:=\;
  \neg\,\CurrentizeWit(\OmegaSAT).
\]
Here \(\CurrentizeWit(\OmegaSAT)\) is the witness-currentization
predicate of \Cref{def:pvnp:currentize-wit}: it asserts that the
canonical witness producer induced by the autonomous lift \(\OmegaSAT\)
is realized by a lawful current observable in \(\EQbang\).  Thus
\(\CSLSATHidden\) asserts that this witness producer is not current on
\(\Tsat\) through \(\Ipoly\).  On the fixed-presentation carrier, the
same source is equivalently the localized residual
\[
  \mathrm{Br}^{\mathrm{lex}}_{\mathrm{SAT}}\notin\EQbang,
\]
where \(\mathrm{Br}^{\mathrm{lex}}_{\mathrm{SAT}}\) is the canonical
lexicographic SAT-branching readout.\footnote{The recognition source
is encoded in Lean as the typed recognition-source carrier listed in
\Cref{app:formalization}; see \Cref{app:formalization} for the
narrowing-disclosure discipline.}
\end{definition}

The localized form
\(\mathrm{Br}^{\mathrm{lex}}_{\mathrm{SAT}}\notin\EQbang\) reads the
canonical lexicographic SAT-branching readout --- the standard
variable-by-variable SAT-branching policy of
\citep{DavisLogemannLoveland1962MachineProgram,NieuwenhuisOliverasTinelli2006DPLL}
--- as an observable, and asserts that it is not lawful in the
saturated family.  This residual form gives \Cref{sec:class-b-barrier}
a concrete target: the barrier theorem acts on the fixed-branch
residual rather than on the abstract negation.  The localization does
not change the conditional theorem's content; it makes precise which
residual the recognition source contributes.

The word ``recognition'' is used here in the framework's audit sense.
A structural recognition source is not a theorem derived inside the
present construction; it is a typed structural input whose content is
tracked and whose uses are explicit.  \citet{TsiokosHiddenness2026}
documents why this source is supplied rather than derived: that paper
records a three-option derivation sweep for the SAT-specific instance
of the \CSLcore{} — a search for a framework-internal derivation, a
search for a reduction from a more abstract structural result, and a
search for a finite-carrier instance witness — and reports that all
three routes fail to produce the SAT-specific witness-currentization
negation.  The present paper therefore takes \(\CSLSATHidden\) as a
structural recognition source supplied by that paper, not as a result
proved from the construction of \(\Tsat\) itself.

This is consistent with the \VDiff{} placement of PvNP in
\citet{TsiokosHiddenness2026}.  That classification separates
substrates whose current observables expose witness-side content from
substrates whose current observables expose only trace state.  Witness-content exposure means that a
current observable reads data such as an assignment or solution
coordinate.  Trace-state exposure means that current observation sees
replay, audit, step, or other history-state information without
revealing the witness content itself.  PvNP is placed on the
trace-only-exposing side.  If PvNP were witness-content-exposing in
the relevant sense, the recognition source would assert the opposite
of what is needed here: that the witness producer is currentizable.
The \VDiff{} placement does not derive \(\CSLSATHidden\), but it gives
the structural context in which the source is plausible.

The carrier fields, their intended semantics, and the way the source is
attached to the saturated SAT layer are recorded in
\Cref{app:recognition-source-carrier}.  This carrier discipline makes
the conditional reading of the contradiction chain precise: every use
of \(\CSLSATHidden\) is a structured invocation of the source, not an
oracle call.

The next section uses \(\CSLSATHidden\) as the single non-derived
input in the contradiction chain.  Every other component of that
chain---Translation A, Translation B, and the standard verifier
package for \(\mathrm{SAT}\in NP\)---is construction-grade or standard
mathematics.  The recognition-source structure is exactly what
separates the framework-internal conditional theorem from an
unconditional theorem in ordinary mathematics.  The barrier theorem
acts on the localized form
\(\mathrm{Br}^{\mathrm{lex}}_{\mathrm{SAT}}\notin \EQbang\), not on the
opaque source name \(\CSLSATHidden\).  The discussion section returns
to this boundary when comparing the present result with the classical
barriers.
 

\section{Four-step contradiction chain}
\label{sec:contradiction-chain}

We now combine the two translation theorems with the recognition
source.  Translation A converts a standard polynomial-time SAT decider
into decision currentization.  Translation B converts decision
currentization into witness currentization.  The recognition source
\(\CSLSATHidden\) says that witness currentization fails on the
formed SAT layer under \(\Ipoly\).  The resulting contradiction gives
\(\mathrm{SAT}\notin P\), conditional on the structural recognition
source.  We then record the standard verifier-package fact
\(\mathrm{SAT}\in NP\).  Together these are the construction-side
inputs for the framework-relative judgment stated in the next section.

\begin{theorem}[Four-step contradiction]
\label{thm:four-step-contradiction-conditional}
Conditional on the structural recognition source \(\CSLSATHidden\),
\[
  \mathrm{SAT}\notin P.
\]
Equivalently, in standard conditional form,
\[
  \CSLSATHidden
  \quad\Longrightarrow\quad
  \mathrm{SAT}\notin P.
\]
Thus, assuming \(\CSLSATHidden\) is supplied as input to the
construction, the SAT decision problem is not in standard polynomial
time.
\end{theorem}

\begin{proof}
Suppose, for contradiction, that \(\mathrm{SAT}\in P\).  By
Translation Theorem A, \Cref{thm:translation-A}, this implies
\[
  \CurrentizeDec(\mathrm{SAT}).
\]
By Translation Theorem B, \Cref{thm:translation-B}, decision
currentization implies witness currentization:
\[
  \CurrentizeWit(\OmegaSAT).
\]
On the other hand, the structural recognition source
\(\CSLSATHidden\), as formalized in
\Cref{def:pvnp:recognition-source-Gamma}, asserts
\[
  \neg\,\CurrentizeWit(\OmegaSAT).
\]
This contradicts the conclusion obtained from Translation B.  Hence
the assumption \(\mathrm{SAT}\in P\) is impossible under
\(\CSLSATHidden\), and therefore \(\mathrm{SAT}\notin P\).\footnote{This
theorem is verified in Lean by the four-step contradiction theorem
declaration listed in \Cref{app:formalization}.}
\end{proof}

\begin{theorem}[SAT is in NP]
\label{thm:pvnp:sat-in-np}
\[
  \mathrm{SAT}\in NP.
\]
Equivalently, SAT has the standard verifier package: a polynomial-size
satisfying assignment serves as a certificate checkable in polynomial
time \citep{Cook1971TheoremProving,AroraBarak2009Complexity}.
\end{theorem}

\begin{proof}
This is the textbook verifier characterization of SAT.  Given a
formula \(\varphi\) and a candidate assignment \(\alpha\), the standard
verifier checks in polynomial time whether \(\alpha\) satisfies every
clause of \(\varphi\); this verifier is the polynomial-time
witness-checker required by the definition of \(NP\).  The assignment
length is bounded by the number of variables, hence by the input length
up to the standard encoding convention, so the witness-size bound is
met.\footnote{This theorem is verified in Lean by the SAT-in-NP
declaration listed in \Cref{app:formalization}; the Lean theorem
transports the verifier-package conclusion from supplied
textbook hypotheses (polynomial-time verifier, witness-size bound,
standard verifier-package combinator) without building a concrete SAT
encoding inside Lean.  The narrowing-disclosure discipline is recorded
in \Cref{app:formalization}.}
\end{proof}

Combining the conditional conclusion \(\mathrm{SAT}\notin P\) with
\(\mathrm{SAT}\in NP\) gives the framework-relative separation that
the next section packages:
\[
  \PSB\neq\NPSB
\]
under the standard closure assumption and the supplied structural
recognition source.  The next section states this in the BirdInt
judgment shape and records the recognition-grade status of the
conclusion.
 \section{BirdInt judgment}
\label{sec:judgment}

The previous section established the construction-side conditional:
under \(\CSLSATHidden\), SAT is not in \(P\), while the standard
verifier package keeps SAT in \(NP\).  This section packages that
conditional theorem in the framework's judgment form.

\begin{definition}[BirdInt judgment for the PvNP closure]
\label{def:pvnp:birdint-judgment}
The \(\BirdInt\) judgment for the PvNP closure is the four-slot
framework-relative assertion
\[
  \BIJudgment{\CSLSATHidden}{\Tsat}{\Ipoly}{\PSB \neq \NPSB}.
\]
The source, layer, instrument, and proposition slots record, in order,
the supplied structural input \(\CSLSATHidden\), the saturated SAT
layer \(\Tsat\), the lawful polynomial-time current instrument
\(\Ipoly\), and the framework-relative separation \(\PSB\neq\NPSB\).
The judgment is qualified by \RecognitionGrade{}.\footnote{The BirdInt
judgment is recorded in the formalization harness as a manifest
obligation rather than a Lean theorem; see \Cref{app:formalization}.}
\end{definition}

Recognition grade means that acceptance uses the supplied recognition
source, not a derivation from framework primitives alone.  Foundations~I
supplies stability of the audit-and-judgment shape
\citep{TsiokosFoundationsI2026}; \(\CSLSATHidden\) supplies the
SAT-specific hiddenness content of \citet{TsiokosHiddenness2026}.  With
that source, the judgment is the bookkeeping outcome of
\Cref{sec:contradiction-chain}.

Read outside the framework, the same content is the conditional
theorem
\[
  \text{formed-SAT-closure-under-\(\Ipoly\)}
  \quad\Longrightarrow\quad
  P\neq NP,
\]
where the antecedent supplies the recognition source and the formed
SAT closure.  The conclusion combines \(\mathrm{SAT}\notin P\) from
the contradiction chain with \(\mathrm{SAT}\in NP\) from the standard
verifier package.

\Cref{sec:pvnp-aor-instance} reads the same conditional content in the
audit-closure register of \citet{Tsiokos2026AOR}, where the same
closure records exhibit \(\Tsat\) as an audit-closed carrier under
\(\CSLSATHidden\).
 

\section{The Class-B Irremovability Barrier}
\label{sec:class-b-barrier}

The conditional theorem of \Cref{sec:contradiction-chain} lands at
recognition grade because it depends on the structural source
\(\CSLSATHidden\).  One can ask whether the source is removable: whether
some framework-internal method can replace it with a derivation.  This
section shows that one natural class of such methods cannot.  No
basis-generated Class-B law on the standard fixed-presentation
polynomial current carrier removes the recognition source.

The section is organized as follows.  Two theorem-grade facts about the
carrier — orbit collapse and polynomial pairing closure — do most of
the work.  These are stated in \Cref{subsec:class-b:meta-results}.
\Cref{subsec:class-b:shape-and-basis} fixes the Class-B law shape and
the native feature basis \(\mathcal B_{\mathrm{fix}}\).  Basis
Combination Closure (\Cref{subsec:class-b:combination-closure}) reduces
every basis-generated feature to one of five normal classes, and Basis
Exhaustiveness (\Cref{subsec:class-b:basis-exhaustiveness}) closes the
basis under intrinsic features.  The Class-B Irremovability theorems
follow in \Cref{subsec:class-b:irremovability}, and the comparison
with the three classical barriers is in
\Cref{subsec:class-b:comparison}.

\subsection{Theorem-grade meta-results}
\label{subsec:class-b:meta-results}

The Class-B analysis rests on two theorem-grade facts about the
standard fixed-presentation polynomial current carrier.  The first says
that lawful current observables separate encoded inputs.  The second
says that separately lawful current observations can be paired without
leaving the saturated current-observable family.

\begin{theorem}[Fixed-Presentation Orbit Collapse]
\label{thm:pvnp:orbit-collapse}
For the standard fixed-presentation polynomial current carrier, the
common kernel of all lawful saturated current observables is the
diagonal:
\[
  \bigcap_{F\in\EQbang}\ker(F)=\Delta_{X}.
\]
Equivalently, if \(x\neq y\) are distinct encoded inputs in the carrier,
then there is a lawful current observable \(F\in\EQbang\) such that
\[
  F(x)\neq F(y).
\]
\end{theorem}

\begin{proof}
The diagonal inclusion
\[
  \Delta_X\subseteq \bigcap_{F\in\EQbang}\ker(F)
\]
is immediate: every observable \(F\) is a function, so \(F(x)=F(x)\)
for all \(x\in X\), and therefore \((x,x)\in\ker(F)\) for every
\(F\in\EQbang\).

For the reverse inclusion, let \(x,y\in X\) with \(x\neq y\).  Distinct
encoded SAT instances differ at some bit, position, or field of the
fixed-presentation encoding.  The identity readout
\(\operatorname{id}_{X}:X\to X\) is realized by a deterministic
polynomial-time transducer that copies its input verbatim, carries the
vacuous no-oracle audit record, and reads only the encoded-instance
field of every history.  Thus it satisfies the five lawfulness clauses
of \Cref{def:pvnp:lawfulness-discipline}, so
\(\operatorname{id}_{X}\in\EQbang\).  Since \(x\neq y\),
\[
  \operatorname{id}_{X}(x)=x\neq y=\operatorname{id}_{X}(y),
\]
and hence \((x,y)\notin\ker(\operatorname{id}_{X})\).  Therefore
\((x,y)\) cannot lie in \(\bigcap_{F\in\EQbang}\ker(F)\).  The common
kernel contains only diagonal pairs.

Equivalently, for distinct encoded inputs \(x\neq y\), the bit
projection at any differing coordinate is a lawful current observable
that separates \(x\) and \(y\).  Therefore
\[
  \bigcap_{F\in\EQbang}\ker(F)=\Delta_X .
\]\footnote{This theorem is verified in Lean by
the orbit-collapse theorem declaration listed in
\Cref{app:formalization}.}
\end{proof}

The second meta-result is the polynomial pairing closure already proved
in \Cref{thm:pvnp:e-q-closure-properties}.  In the present barrier
context we use only the following consequence: if \(F,G\in\EQbang\),
then the paired observable
\[
  x\longmapsto (F(x),G(x))
\]
also belongs to \(\EQbang\).  More generally, finite products of
fixed-stage lawful current observables remain lawful current
observables.  This is not a new theorem in this section; it is the
finite-product clause of the observable-saturation closure theorem,
restated here because the Class-B barrier uses it directly.

Together, orbit collapse and pairing closure eliminate two natural
Class-B mechanisms on the standard polynomial carrier.  Orbit collapse
kills quotient-mediated obstructions: there is no nontrivial quotient
respected by all lawful current observables.  Pairing closure kills
joint-access-mediated obstructions: if the separate accesses are
current, so is their finite joint access.  The rest of the section
makes this reduction precise and shows that no other Class-B mechanism
generated by the native basis survives either.

\subsection{Class-B law shape and feature basis}
\label{subsec:class-b:shape-and-basis}

A \emph{Class-B carrier-property law} is a theorem whose obstruction
comes from an intrinsic feature of the carrier.  Its schematic form is:
if the carrier has feature \(\Phi\), then no admissible observable in
\(\EQbang\) can attain a target \(T\).  The point is that the
impossibility is carried by the carrier feature itself, not by rewriting
what the target means and not by restricting what a claim is allowed to
assert.  This distinguishes Class-B laws from the framework's Class-A
claim-discipline laws, such as no-overread, adequacy, and quotientality,
which deny overclaims by the structure of the claim being made.  Both
classes are legitimate Six Birds methods, but the present barrier is
about Class-B laws because those are the methods that would try to
remove \(\CSLSATHidden\) by exploiting carrier features.  Structural
treatments of complexity classes have a long methodological precedent:
descriptive complexity gives a machine-independent characterization of
\(NP\) via first-order spectra
\citep{Fagin1974GeneralizedSpectra,Immerman1999DescriptiveComplexity},
and more recent categorical approaches characterize complexity through
the structure of admissible morphisms and functors
\citep{BasuIsik2020CategoricalComplexity}.  The Class-B framework is in
the same broad family of structural approaches; the formal Class-B law
record specific to the present setting is given in
\Cref{app:class-b-definitions}.

For the standard fixed-presentation polynomial current carrier, the
native feature basis is denoted \(\mathcal B_{\mathrm{fix}}\).  It has
four primitives.  First, \(\Delta_{X}\) is the diagonal, or equality
relation, on the encoded-input carrier.  Second,
\(\mathsf{PolyClosure}(\EQbang)\) is the closure structure of the
saturated current-observable family: pairing, finite products,
polynomial-time post-processing, and polynomial-time composition, as
proved in \Cref{thm:pvnp:e-q-closure-properties}.  Third,
\(\mathsf{TimeFilt}\) is the polynomial resource ladder
\(E^{[1]}\subseteq E^{[2]}\subseteq\cdots\), with
\(\EQbang=\bigcup_{k}E^{[k]}\).  Fourth,
\(\mathsf{PresData}\) is the fixed-presentation data: bit positions,
variable names, symbol order, clause order, delimiter conventions, and
similar structural features of the chosen encoding.  These are the
carrier-native ingredients from which the basis-generated Class-B
features are allowed to be built; their formal definitions are recorded
in \Cref{app:class-b-definitions}.

A feature is \emph{basis-generated} when it is built from
\(\mathcal B_{\mathrm{fix}}\) by finitely many applications of the
allowed closure operations: quotienting, orbiting, transporting,
holonomy or gluing-class formation, pairing, products, polynomial
post-processing, resource filtration, and presentation refinement.  Each
basis-generated feature is assigned a normal class.  The quotient class
\(Q\) records an equivalence relation that the proposed Class-B law
requires every lawful current observable to respect.  The joint-access
class \(J\) records a tuple of separately current observables whose
joint access is alleged to be inadmissible.  The positive-closure class
\(C\) records a construction yielding new current observables from
current ones.  The presentation-refinement class \(P\) records a
refinement of the carrier by presentation data; it supplies breaker
data, not an obstruction.  The time-filtration class \(T\) records a
feature indexed by polynomial resource strata.  The Basis Combination
Closure theorem in the next subsection proves that these five classes
are exhaustive for basis-generated features.  The formal inductive
definition of basis generation, together with the normal-class
assignment, is deferred to \Cref{app:class-b-definitions}.

\subsection{Basis Combination Closure}
\label{subsec:class-b:combination-closure}

\Cref{subsec:class-b:shape-and-basis} introduced basis-generated
features and the five normal classes.  The first structural reduction
is that every finite basis-generated feature lands in one of those
classes, and in each class the resulting feature is either trivial,
eliminated by closure, positive rather than obstructive, presentation
data rather than invariance, or equivalent to the target claim at the
full polynomial union.

\begin{theorem}[Basis Combination Closure]
\label{thm:pvnp:basis-combination-closure}
Let \(\Phi\) be a finite basis-generated feature of the standard
fixed-presentation polynomial current carrier.  Then \(\Phi\) lies in
one of the five normal classes \(Q,J,C,P,T\).  Moreover, the following
class consequence holds.
\begin{enumerate}
\item[\(Q\).] If \(\Phi\) is a quotient feature, then any equivalence
  relation respected by every \(F\in\EQbang\) is the diagonal; orbit
  collapse forces the feature to be trivial.
\item[\(J\).] If \(\Phi\) is a joint-access feature, then any finite
  tuple of separately current observables admits a jointly current
  product observable; pairing closure kills the alleged obstruction.
\item[\(C\).] If \(\Phi\) is a positive-closure feature, then it
  constructs new current observables from current observables and does
  not exclude a target.
\item[\(P\).] If \(\Phi\) is a presentation-refinement feature, then
  it supplies breaker or refinement data rather than an invariance
  obstruction.
\item[\(T\).] If \(\Phi\) is a time-filtration feature, then it is
  either fixed-stratum only, with no obstruction at the full carrier
  \(\EQbang=\bigcup_k E^{[k]}\), or it is target-equivalent at the
  full polynomial union.
\end{enumerate}
\end{theorem}

\begin{proof}
The proof is by induction on the constructor depth of the
basis-generated feature expression.  The four primitive base cases are
as follows: the diagonal \(\Delta_X\) lies in \(Q\); the polynomial
closure structure \(\mathsf{PolyClosure}(\EQbang)\) lies in \(C\); the
resource ladder \(\mathsf{TimeFilt}\) lies in \(T\); and the
presentation data \(\mathsf{PresData}\) lies in \(P\).

For the inductive step, each allowed closure operation preserves the
normal-class assignment.
\begin{itemize}
\item Quotienting, orbiting, and holonomy or gluing-class formation
  reduce to \(Q\).  The quotient closure lemmas
  \Cref{lem:pvnp:quotient-meet-join,lem:pvnp:orbit-quotient,lem:pvnp:holonomy-quotient},
  together with orbit collapse \Cref{thm:pvnp:orbit-collapse}, force
  the quotient feature to be diagonal.
\item Transport splits into route, record-carrying, and
  resource-indexed forms, reducing respectively to \(Q\), \(P\), and
  \(T\).
\item Pairing and finite products are constructive closure operations:
  as positive closure they reduce to \(C\), while as a proposed
  joint-access obstruction the finite-product clause of
  \Cref{thm:pvnp:e-q-closure-properties} kills the obstruction after
  the directed-union lift to a common stratum: any finite family in
  \(\EQbang=\bigcup_k E^{[k]}\) lies in some \(E^{[k^*]}\), so the
  fixed-stage hypothesis of the finite-product clause is met.
\item Polynomial post-processing reduces to \(C\) by
  \Cref{lem:pvnp:polynomial-postprocessing-closure}.
\item Presentation refinement reduces to \(P\) by
  \Cref{lem:pvnp:presentation-refinement}.
\item Resource filtration reduces to \(T\) by
  \Cref{lem:pvnp:time-stratum}.
\end{itemize}

Therefore a feature built from \(n+1\) operations applied to features
of depth at most \(n\) remains in one of \(Q,J,C,P,T\), with the
corresponding class consequence.  Mixed-class combinations are handled
by the same closure lemmas.  The fifteen pairs from \(Q\times Q\)
through \(T\times T\) reduce to the same five normal classes; the case
table is given in \Cref{app:class-b-closure-lemmas}.\footnote{This
theorem is verified in Lean by the basis combination closure
declaration listed in \Cref{app:formalization}; the Lean proof performs
a per-constructor case analysis over the feature-expression type with
the closure lemmas invoked at the corresponding cases.}
\end{proof}

The theorem reduces any basis-generated Class-B law to a five-way case
analysis on its normal class.  This reduction is used in
\Cref{thm:pvnp:irremovability-scoped}.

\subsection{Basis Exhaustiveness}
\label{subsec:class-b:basis-exhaustiveness}

\Cref{thm:pvnp:basis-combination-closure} reduces any
basis-generated feature to one of five normal classes.  The remaining
question is whether a Class-B law on the standard fixed-presentation
polynomial current carrier could use an intrinsic feature that is not
basis-generated.  Basis Exhaustiveness closes that gap.

An \emph{intrinsic feature} of the fixed-presentation carrier is a
feature definable from the carrier package
\((X,\EQbang,\Ipoly,\mathsf{Pres})\) without importing extra structure:
no external predicate, measure, topology, oracle, group action,
algebraic structure, advice string, probability law, or metric that is
not already part of the carrier package.  Equivalently, it is a feature
invariant under isomorphisms of the carrier package.  The formal
definition is given in \Cref{app:class-b-definitions}.

\begin{theorem}[Basis Exhaustiveness]
\label{thm:pvnp:basis-exhaustiveness}
Every intrinsic feature of the standard fixed-presentation polynomial
current carrier is generated from
\[
  \mathcal B_{\mathrm{fix}}
  =
  \{\Delta_X,\mathsf{PolyClosure}(\EQbang),
    \mathsf{TimeFilt},\mathsf{PresData}\}
\]
by the closure operations of \Cref{subsec:class-b:shape-and-basis}.
Consequently, together with
\Cref{thm:pvnp:basis-combination-closure}, every intrinsic feature
falls into one of the five normal classes \(Q,J,C,P,T\) with the
corresponding class consequence.
\end{theorem}

\begin{proof}
The proof is by source-slot decomposition.  An intrinsic feature of
the carrier package \((X,\EQbang,\Ipoly,\mathsf{Pres})\) draws its
content from one of four slots.

In the carrier-set slot \(X\), the set is a bare encoded-input set.
The only intrinsic relation available at this level is the diagonal
\(\Delta_X\).  Any nontrivial quotient required to be respected by
every \(F\in\EQbang\) collapses to \(\Delta_X\) by
\Cref{thm:pvnp:orbit-collapse}.

In the observable-family slot \(\EQbang\), the intrinsic structure is
the closure algebra \(\mathsf{PolyClosure}(\EQbang)\), recorded in
\Cref{thm:pvnp:e-q-closure-properties}.  Common-kernel features again
collapse to \(\Delta_X\) by orbit collapse.

In the polynomial-time interface slot \(\Ipoly\), the nontrivial
structural parameter is the resource ladder \(\mathsf{TimeFilt}\).
The remaining interface conditions, such as determinism, uniformity,
all-input totality, and no-oracle bookkeeping, are typing requirements
for membership in \(\EQbang\), not independent intrinsic features.

In the fixed-presentation slot \(\mathsf{Pres}\), the encoding
contributes \(\mathsf{PresData}\): bit positions, variable names,
symbol order, clause order, delimiter conventions, and related
presentation features.

Combinations of these four primitives under the allowed closure
operations are basis-generated by construction, with their normal-class
consequences supplied by
\Cref{thm:pvnp:basis-combination-closure}.  Candidate fifth sources of
structure — categorical or two-categorical structure, diagonalization,
topology or metric data, probability or measure data, additional
resources such as space or randomness, algebraic structure, order or
prefix structure, and semantic SAT structure — each either reduce to
one of the four basis primitives or leave the standard
fixed-presentation carrier for an enriched carrier.  These case checks
are recorded in \Cref{app:class-b-definitions}.\footnote{The
formalization harness tracks this theorem through the basis
exhaustiveness declaration listed in \Cref{app:formalization}, as a
conditional schema over a typed decomposition record; the substantive
content is the source-slot argument given here.}
\end{proof}

Combined with Basis Combination Closure, Basis Exhaustiveness lets the
Irremovability theorem treat any Class-B law on the standard carrier by
a five-way normal-class analysis.

\subsection{Class-B Irremovability}
\label{subsec:class-b:irremovability}

\Cref{thm:pvnp:basis-combination-closure} reduces basis-generated
features to the five normal classes, and
\Cref{thm:pvnp:basis-exhaustiveness} says that intrinsic features of
the standard fixed-presentation carrier are basis-generated.  The
consequence is the Class-B Irremovability barrier: carrier-feature
methods of this kind cannot remove the localized recognition residual.

\begin{theorem}[Class-B Irremovability Core]
\label{thm:pvnp:irremovability-core}
No quotient-mediated or joint-access-mediated Class-B law on the
standard fixed-presentation polynomial current carrier can prove
\[
  \mathrm{Br}^{\mathrm{lex}}_{\mathrm{SAT}}\notin\EQbang .
\]
Equivalently, no Class-B law of either of these two mechanisms can
remove \textnormal{\(\CSLSATHidden\)}.
\end{theorem}

\begin{proof}
Let \(\mathcal L_B\) be a Class-B law of one of the two mechanisms.
We analyze the two cases.

In the quotient-mediated case, the proposed law asserts that some
nontrivial equivalence relation \(R\) on the encoded-input carrier
\(X\) is respected by every \(F\in\EQbang\), and that this respect
obstructs
\(\mathrm{Br}^{\mathrm{lex}}_{\mathrm{SAT}}\in\EQbang\).  Respect for
\(R\) by every lawful current observable means
\[
  R\subseteq \bigcap_{F\in\EQbang}\ker(F).
\]
By \Cref{thm:pvnp:orbit-collapse}, the common kernel of all lawful
current observables is the diagonal \(\Delta_X\).  Hence
\(R\subseteq\Delta_X\), so the quotient is the identity relation.  The
identity relation imposes no constraint on any readout, since every
function on \(X\) respects equality.  The quotient mechanism therefore
exerts no force on whether
\(\mathrm{Br}^{\mathrm{lex}}_{\mathrm{SAT}}\) lies in \(\EQbang\), and
cannot prove
\(\mathrm{Br}^{\mathrm{lex}}_{\mathrm{SAT}}\notin\EQbang\).

In the joint-access-mediated case, the proposed law asserts that some
finite family \(F_1,\ldots,F_n\in\EQbang\) of separately current
observables has inadmissible joint access, and that this joint-access
failure obstructs
\(\mathrm{Br}^{\mathrm{lex}}_{\mathrm{SAT}}\in\EQbang\).  The alleged
joint observable is
\[
  x\longmapsto (F_1(x),\ldots,F_n(x)).
\]
By the finite-product clause of
\Cref{thm:pvnp:e-q-closure-properties}, finite products of lawful
current observables at a common audit stage remain in \(\EQbang\).
Since \(\EQbang=\bigcup_k E^{[k]}\) is directed over the polynomial
resource ladder, each \(F_i\) lies in some stratum \(E^{[k_i]}\), and
all \(F_i\) lie in the common stratum \(E^{[k^*]}\) for
\(k^*=\max(k_1,\ldots,k_n)\).  The fixed-stage hypothesis is therefore
satisfied.  Hence \((F_1,\ldots,F_n)\in\EQbang\): the alleged
joint-access obstruction is not an obstruction.  The joint-access
mechanism therefore exerts no force on
\(\mathrm{Br}^{\mathrm{lex}}_{\mathrm{SAT}}\), and cannot prove
\(\mathrm{Br}^{\mathrm{lex}}_{\mathrm{SAT}}\notin\EQbang\).

In each case the mechanism collapses on the standard
fixed-presentation carrier.  Hence no Class-B law of either mechanism
can derive
\(\mathrm{Br}^{\mathrm{lex}}_{\mathrm{SAT}}\notin\EQbang\).\footnote{This
theorem is verified in Lean by
the Class-B irremovability core declaration listed in
\Cref{app:formalization}; the Lean proof performs the same two-case
disjunction, applying orbit collapse and pairing closure at the
corresponding cases.}
\end{proof}

\begin{theorem}[Six Birds Class-B Irremovability, scoped form]
\label{thm:pvnp:irremovability-scoped}
Let \(\mathcal L_B\) be a Class-B carrier-feature law of the standard
fixed-presentation polynomial current carrier, analyzed against the
basis \(\mathcal B_{\mathrm{fix}}\).  Then at least one of the
following holds:
\begin{enumerate}
\item \(\mathcal L_B\) cannot derive
  \(\mathrm{Br}^{\mathrm{lex}}_{\mathrm{SAT}}\notin\EQbang\);
\item \(\mathcal L_B\) introduces a feature not generated by
  \(\mathcal B_{\mathrm{fix}}\);
\item \(\mathcal L_B\) assumes a target-equivalent lower bound, namely
  a full-polynomial-union obstruction whose content is already
  equivalent to
  \(\mathrm{Br}^{\mathrm{lex}}_{\mathrm{SAT}}\notin\EQbang\).
\end{enumerate}
\end{theorem}

\begin{proof}
Let \(\Phi\) be the carrier feature used by \(\mathcal L_B\).  First
suppose that \(\Phi\) is intrinsic to the standard fixed-presentation
polynomial current carrier.  By \Cref{thm:pvnp:basis-exhaustiveness},
\(\Phi\) is basis-generated.  By
\Cref{thm:pvnp:basis-combination-closure}, \(\Phi\) falls into one of
the five normal classes \(Q,J,C,P,T\).  We treat the five cases.

In class \(Q\), \(\Phi\) is quotient-mediated: a nontrivial
equivalence relation \(R\) on \(X\) is required to be respected by
every \(F\in\EQbang\).  By the quotient part of
\Cref{thm:pvnp:irremovability-core}, \(R\) collapses to the diagonal
\(\Delta_X\), and the quotient mechanism cannot derive
\[
  \mathrm{Br}^{\mathrm{lex}}_{\mathrm{SAT}}\notin\EQbang .
\]
Thus the first alternative holds.

In class \(J\), \(\Phi\) is joint-access-mediated: a finite tuple of
separately current observables is alleged to fail joint admission.  By
the pairing-closure part of \Cref{thm:pvnp:irremovability-core},
including the directed-union lift to a common audit stratum before the
finite-product clause is applied, the alleged obstruction is killed by
closure.  The joint-access mechanism cannot derive
\(\mathrm{Br}^{\mathrm{lex}}_{\mathrm{SAT}}\notin\EQbang\), so the
first alternative again holds.

In class \(C\), \(\Phi\) is a positive-closure feature: it constructs
new current observables from current observables.  Such a feature is
constructive rather than exclusionary, so it does not exclude any
target readout, in particular
\(\mathrm{Br}^{\mathrm{lex}}_{\mathrm{SAT}}\).  The first alternative
holds.

In class \(P\), \(\Phi\) is a presentation-refinement feature, such as
bit positions, variable names, or order conventions.  These features
supply breaker or refinement data rather than an invariance
obstruction.  They therefore cannot derive
\(\mathrm{Br}^{\mathrm{lex}}_{\mathrm{SAT}}\notin\EQbang\), and the
first alternative holds.

In class \(T\), \(\Phi\) is indexed by polynomial resource strata
\(E^{[k]}\).  If \(\Phi\) is fixed-stratum, then an exclusion at one
stratum does not imply exclusion at the full carrier
\[
  \EQbang=\bigcup_k E^{[k]},
\]
so the first alternative holds.  If \(\Phi\) is a full-polynomial-union
statement, then it is already equivalent to
\(\mathrm{Br}^{\mathrm{lex}}_{\mathrm{SAT}}\notin\EQbang\).  In that
case the law has not removed the recognition source; it has restated
the target-equivalent lower bound, so the third alternative holds.

It remains to consider the case in which \(\Phi\) is not intrinsic to
the standard carrier, or otherwise imports carrier structure not
generated by \(\mathcal B_{\mathrm{fix}}\).  Then
\(\mathcal L_B\) introduces a non-basis feature, and the second
alternative holds by definition.  The intrinsic five-way analysis and
the non-basis case exhaust the theorem, so at least one of the three
alternatives holds for \(\mathcal L_B\).\footnote{This theorem is
verified in Lean by the scoped Class-B irremovability declaration listed in
\Cref{app:formalization}; the Lean theorem operates at the
feature-classification layer, ranging over a \texttt{classBLawFeature}
value rather than directly over the \texttt{classBLaw} record, with the
connection from a
\texttt{classBLaw} record to its feature classification recorded
paper-side.  The narrowing-disclosure discipline is recorded in
\Cref{app:formalization}.}
\end{proof}

Combining this theorem with the localization in
\Cref{sec:recognition-source} gives the intended barrier reading: no
basis-generated Class-B law on the standard fixed-presentation
polynomial current carrier can remove \(\CSLSATHidden\).  The theorem
acts on the localized residual
\(\mathrm{Br}^{\mathrm{lex}}_{\mathrm{SAT}}\notin\EQbang\); the
recognition-source section explains why that residual is the concrete
form of the supplied source.

The scope is deliberately narrow: the theorem rules out one specified
class of removals on one specified carrier.  The corresponding scope
fences — no claim of unconditional \(P\neq NP\), no claim of absolute
irremovability of \(\CSLSATHidden\), no statement about enriched
carriers or Class-A methods — are recorded in \Cref{sec:scope}.

\subsection{Comparison to classical barriers}
\label{subsec:class-b:comparison}

Baker, Gill, and Solovay's relativization barrier
\citep{BakerGillSolovay1975Relativization} rules out proof methods
whose reasoning remains valid after arbitrary oracle access is added.
The carrier-level fact is that there are oracles \(A\) and \(B\) with
\(P^A=NP^A\) and \(P^B\neq NP^B\).  A proof technique that goes
through unchanged relative to both kinds of oracle therefore cannot
settle the unrelativized \(P\)-versus-\(NP\) question.  Relativization
is a methodological barrier, not a resolution of the question itself.

The Razborov--Rudich natural-proofs barrier
\citep{RazborovRudich1997NaturalProofs} rules out a broad class of
combinatorial lower-bound methods.  In that setting a natural property
of Boolean functions is constructive, large, and useful against the
circuit class under attack.  The carrier-level fact is
pseudorandomness: a large constructive property useful against all easy
functions would distinguish pseudorandom functions from truly random
functions.  Conditional on pseudorandom function generators of the
required strength, no such natural property proves the strong circuit
lower bounds needed for \(P\) versus \(NP\).  Subsequent work has
refined the largeness and constructivity conditions
\citep{Chow2011AlmostNaturalProofs} without removing the barrier in its
intended scope.

Aaronson and Wigderson's algebrization barrier
\citep{AaronsonWigderson2009Algebrization} extends relativization
beyond ordinary oracle access.  It rules out proof methods whose
reasoning survives arithmetization, where oracle access is extended to
low-degree extensions over finite fields.  The carrier-level fact is
that algebraic oracle access is broad enough to model many
arithmetization techniques and still admits worlds where the target
separation behaves differently.  Algebrization is therefore a
methodological barrier for proof strategies that remain valid under
that algebraic extension of the oracle setting.  An axiomatic
reformulation that treats algebrization as a closed proof-method class
appears in
\citep{ImpagliazzoKabanetsKolokolova2009AxiomaticAlgebrization}.

The Six Birds Class-B Irremovability theorem rules out a different
methodological class: basis-generated Class-B carrier-feature laws on
the standard fixed-presentation polynomial current carrier.  The
carrier-level facts doing the work are
\Cref{thm:pvnp:orbit-collapse} and
\Cref{thm:pvnp:e-q-closure-properties}, together with
\Cref{thm:pvnp:basis-exhaustiveness} and
\Cref{thm:pvnp:basis-combination-closure}.  Its scope is internal to
Six Birds.  It is not a resolution of \(P\) versus \(NP\) itself, not
absolute irremovability of \(\CSLSATHidden\), and not a barrier on
enriched carriers.

The analogy with the classical barriers is an analogy of role.  Each
barrier rules out a methodological class on a standard computational
object: relativization uses oracle separation, natural proofs use
pseudorandomness, algebrization uses low-degree extension separation,
and Class-B Irremovability uses the intrinsic feature basis of a typed
carrier together with closure over that basis.  The four barriers are
therefore complementary rather than mutually subsuming.  Successful
lower-bound techniques that pass the classical barriers in specific
settings do exist; Williams's nonuniform \(\mathsf{ACC}\) lower bound
\citep{Williams2014ACCLowerBounds} is the canonical modern example.  A
future proof of \(P\) versus \(NP\) would have to avoid relativization,
naturalness, algebrization, and this Class-B carrier-feature
obstruction on the standard fixed-presentation polynomial carrier; a
broader perspective treats such barriers as independence-like phenomena
relative to particular proof methods \citep{Kolokolova2016BarriersIndependence}.

The limitation is also part of the statement.  Unlike the classical
barriers, the Six Birds Class-B barrier is internal to a particular
framework.  It does not rule out methods outside Six Birds, Class-B
laws on objective, branchwise, presentation-quotiented, or otherwise
enriched carriers, or Class-A claim-discipline laws.  It should be read
neither as equivalent to the classical barriers nor as weaker
bookkeeping: it is a framework-internal barrier theorem identifying a
principled limit of Six Birds Class-B reasoning on the standard
polynomial current layer.  \Cref{sec:scope} records the corresponding
scope fences explicitly.
 

\section{PvNP as an AOR instance}
\label{sec:pvnp-aor-instance}

The conditional separation theorem may also be read in the
Audited Operational Realisability vocabulary of
\citet{Tsiokos2026AOR}.  This reading does not introduce a second
proof of \(\PSB\neq\NPSB\).  It repackages the same typed records
already used in the construction: the saturated SAT layer, the
lawful polynomial-time instrument, the current and predictive
quotients, the autonomous lift, the currentization predicates, the
translation theorems, the recognition source, the Class-B scope
fence, and the \(\BirdInt\) judgment.  In the AOR register, the
sentence ``the PvNP closure holds on \(\Tsat\)'' means that these
records form an audit-closed carrier for the declared SAT scope.

The correct target class is \(\mathbf{AOR}^{\mathrm{casc}}\), not
the full refinement-stable class.  The history carrier is presented
as a directed colimit of finite-stage carriers, and the observable
families are saturated; but the proof does not use a nontrivial
asymptotic discharge analogous to a scale or frequency-limit
argument.  The live work is a finite discharge cascade:
construct the saturated carrier, identify current and predictive
quotients, translate standard \(P\)-membership into currentization,
translate decision currentization into witness currentization, supply
the recognition-source negation, and package the contradiction in
\(\BirdInt\).  Thus the membership statement is
\[
  \Tsat \in
  \mathbf{AOR}^{\mathrm{casc}}_{\mathfrak S_{\mathrm{PvNP}}},
\]
where \(\mathfrak S_{\mathrm{PvNP}}\) is the declared SAT scope under
\(\Ipoly\).  A stronger refinement-stable statement would require a
separate audit showing that every directed refinement of the SAT
history carrier and observable families remains covered by the same
payload library.  That stronger statement is not needed for the
conditional theorem proved here.

\subsection{AOR-membership reading of the SAT carrier}

The AOR reading is carrier-relative.  The carrier is not the set of
all mathematical descriptions of SAT, nor the standard complexity
class \(P\) considered in isolation.  It is the saturated SAT layer
\[
  \Tsat
  =
  (\Saud,\Hsat,\EQbang,\EMbang,\Qsat,\Msat,\ComparisonMap,
    \operatorname{Audit}^{!})
\]
of \Cref{def:pvnp:saturated-closure}, equipped with the lawful
polynomial-time current instrument \(\Ipoly\), the current-observable
family \(\EQbang\), the autonomous lift \(\OmegaSAT\), and the
visibility relation \(\VisBang\).

The record-level reading is this: within the formed SAT layer, every
native residual used by the proof of the framework-relative judgment
\[
  \BIJudgment{\CSLSATHidden}{\Tsat}{\Ipoly}{\PSB\neq\NPSB}
\]
has a typed status.  The mechanical SAT-layer records are
by-construction or imported as ordinary polynomial-time computation
facts.  The structural residual is the recognition-source residual:
\(\CSLSATHidden\) supplies the failure of witness currentization on
\(\Tsat\), localized to the fixed-branch readout
\[
  \mathrm{Br}^{\mathrm{lex}}_{\mathrm{SAT}}\notin\EQbang .
\]
The Class-B barrier does not remove this source.  It records that one
specified removal route---basis-generated Class-B carrier-feature
laws on the standard fixed-presentation polynomial carrier---is
blocked.

Thus the AOR-instance statement has the same content as the
recognition-grade theorem.  It says that the conditional proof stack
is an audit-closed carrier package.  It does not say that
\(\CSLSATHidden\) has been derived from the other fields of \(\Tsat\).

\subsection{The saturated SAT layer as an AOR-instance object}

\begin{definition}[The AOR instance object attached to \(\Tsat\)]
\label{def:pvnp:aor-instance:object}
Let \(\mathfrak S_{\mathrm{PvNP}}\) be the SAT scope consisting of the
saturated SAT layer \(\Tsat\), the lawful polynomial-time current
instrument \(\Ipoly\), the saturated current and predictive observable
families \(\EQbang\) and \(\EMbang\), the quotients \(\Qsat\) and
\(\Msat\), the comparison map \(\ComparisonMap\), the autonomous lift
\(\OmegaSAT\), the visibility relation \(\VisBang\), the
currentization predicates \(\CurrentizeDec\) and \(\CurrentizeWit\),
the recognition source \(\CSLSATHidden\), the standard SAT verifier
package, the Class-B basis \(\mathcal B_{\mathrm{fix}}\), and the
\(\BirdInt\) judgment package.

The AOR instance object
\[
  C_{\mathrm{PvNP}}(\Tsat,\Ipoly,\CSLSATHidden)
\]
has the following typed fields.
\begin{enumerate}
  \item The carrier is \(\Hsat\), read through the finite-stage
        carriers described in \Cref{def:pvnp:history-carrier} and the
        saturated aggregate of \Cref{def:pvnp:saturated-closure}.
  \item The current observations are the lawful polynomial-time
        observables \(\EQbang\) of \Cref{def:pvnp:e-q-saturation}
        under the instrument \(\Ipoly\) of \Cref{def:pvnp:i-poly}.
        The predictive observations are \(\EMbang\).
  \item The current and predictive quotients are \(\Qsat\) and
        \(\Msat\), with comparison map \(\ComparisonMap\), as in
        \Cref{def:pvnp:q-saturated,def:pvnp:m-saturated,def:pvnp:comparison-map}.
  \item The routes are the polynomial-time post-processing, finite
        product, and composition operations of
        \Cref{thm:pvnp:e-q-closure-properties}, together with the
        self-reduction route used in Translation Theorem~B.
  \item The source records are the standard \(\mathrm{SAT}\in P\)
        bridge used in Translation Theorem~A, the self-reduction
        package used in Translation Theorem~B, the standard verifier
        package for \(\mathrm{SAT}\in NP\), and the recognition source
        \(\CSLSATHidden\) supplied by \citet{TsiokosHiddenness2026}.
  \item The hidden-state record is the autonomous lift \(\OmegaSAT\)
        of \Cref{def:pvnp:omega-sat}, together with the visibility
        relation \(\VisBang\) of \Cref{def:pvnp:vis-bang}.
  \item The target records are decision currentization
        \(\CurrentizeDec(\mathrm{SAT})\), witness currentization
        \(\CurrentizeWit(\OmegaSAT)\), the localized fixed-branch
        residual
        \(\mathrm{Br}^{\mathrm{lex}}_{\mathrm{SAT}}\notin\EQbang\),
        and the judgment proposition \(\PSB\neq\NPSB\).
  \item The residual statuses are the discharges recorded in
        \Cref{thm:pvnp:aor:mechanical-records,thm:pvnp:aor:recognition-source-discharge,thm:pvnp:aor:translation-discharge,thm:pvnp:aor:class-b-scope-fence,thm:pvnp:aor:judgment-package}.
  \item The nonclaim register is the register of \Cref{sec:scope},
        including the absence of an unconditional standard-ZFC
        theorem, the absence of a claim to circumvent classical
        barriers in their native formulations, the non-derivation of
        \(\CSLSATHidden\), and the scoped nature of Class-B
        Irremovability.
\end{enumerate}
\footnote{Realized in Lean by the AOR-instance object declaration
listed in \Cref{app:formalization}; carrier-attached to the concrete
PvNP types.}
\end{definition}

The definition is a bookkeeping object.  It is not a new carrier and
not a new complexity-theoretic hypothesis.  It places the existing
records of the paper into the AOR field order.

\subsection{Mechanical records as membership components}

The definition-side rows of the PvNP construction are mechanical in
the AOR sense.  They introduce typed carriers, instruments, targets,
quotients, and operation records.  They do not supply the hiddenness
content of \(\CSLSATHidden\).

\begin{theorem}[Mechanical AOR records for the saturated SAT layer]
\label{thm:pvnp:aor:mechanical-records}
For the scope \(\mathfrak S_{\mathrm{PvNP}}\), the following records
are AOR-mechanical:
\[
\begin{gathered}
  \text{\Cref{def:pvnp:audited-shell}, }
  \text{\Cref{def:pvnp:history-carrier}, }
  \text{\Cref{def:pvnp:saturated-closure},}\\
  \text{\Cref{def:pvnp:i-poly}, }
  \text{\Cref{def:pvnp:lawfulness-discipline}, }
  \text{\Cref{def:pvnp:e-q-saturation},}\\
  \text{\Cref{def:pvnp:q-saturated}, }
  \text{\Cref{def:pvnp:m-saturated}, }
  \text{\Cref{def:pvnp:comparison-map},}\\
  \text{\Cref{def:pvnp:omega-sat}, }
  \text{\Cref{def:pvnp:vis-bang}, }
  \text{\Cref{def:pvnp:currentize-dec}, }
  \text{\Cref{def:pvnp:currentize-wit}.}
\end{gathered}
\]
The operation-closure theorem \(\Cref{thm:pvnp:e-q-closure-properties}\)
is mechanical at the admission level and \(\texttt{approved\_other}\)
for the ordinary polynomial-time computation facts carried by its
operation records.  The standard verifier package
\(\Cref{thm:pvnp:sat-in-np}\) is \(\texttt{approved\_other}\) from
textbook SAT verification.
\end{theorem}

\begin{proof}
The listed definitions supply the typed carrier, instrument,
observable-family, quotient, comparison, autonomous-lift, visibility,
and currentization records used by the proof.  Their AOR status is
therefore \(\texttt{by\_construction}\).  The closure properties of
\(\EQbang\) add typed operation records for polynomial-time
post-processing, finite products, and polynomial-time composition.
By-construction and approved-other discharges name these typed-shell
definitions and operation-record contracts; the ordinary semantic
computation clauses are carried as \(\texttt{approved\_other}\)
polynomial-time computation facts.  The verifier package for SAT is
likewise an \(\texttt{approved\_other}\) standard verifier record.
None of these rows supplies or derives the recognition-source
negation.\footnote{This theorem is verified in Lean by the
AOR-mechanical-records declaration listed in \Cref{app:formalization};
the Lean theorem inherits the narrowed-surrogate status of
\texttt{eQClosureProperties} (admission-level closure via
operation-record contracts) and \texttt{satInNP} (transport from
textbook verifier hypotheses).  The narrowing-disclosure discipline is
recorded in \Cref{app:formalization}.}
\end{proof}

\subsection{Substantive discharges}

The substantive AOR work is concentrated in the recognition source,
the two translation interfaces, the Class-B scope fence, and the
\(\BirdInt\) packaging of the judgment.

\begin{theorem}[Recognition-source discharge]
\label{thm:pvnp:aor:recognition-source-discharge}
The recognition source \(\CSLSATHidden\) has primary residual type
\(\Delta_{\mathrm{source}}\).  Its forced secondary residual types are
\(\Delta_{\mathrm{target}}\), \(\Delta_{\mathrm{role}}\), and
\(\Delta_{\mathrm{horizon}}\).  Its AOR-8 status is
\(\texttt{hidden}\): the source is read here as the typed assertion
that the autonomous-lift witness content is hidden from current
observation on the formed SAT carrier.  It supplies
\[
  \neg\,\CurrentizeWit(\OmegaSAT),
\]
localized in the present paper to
\[
  \mathrm{Br}^{\mathrm{lex}}_{\mathrm{SAT}}\notin\EQbang .
\]
\end{theorem}

\begin{proof}
The source is not generated by the construction of \(\Tsat\).  It is
the provenance record for the witness-currentization negation used in
\Cref{thm:four-step-contradiction-conditional}.  The target component
is the predicate \(\neg\,\CurrentizeWit(\OmegaSAT)\), equivalently the
fixed-branch residual
\(\mathrm{Br}^{\mathrm{lex}}_{\mathrm{SAT}}\notin\EQbang\).  The role
component records that this is a recognition source rather than a
framework-derived theorem.  The horizon component is the boundary
between current observation through \(\EQbang\) and the autonomous
lift's witness-side content.  The citation provenance is a source
record; the residual being statused is the hiddenness of witness data,
so the AOR-8 status is \(\texttt{hidden}\), not \(\texttt{bridged}\).
No transport, scale, limit, presentation, or local/global obligation
is generated: no route, asymptotic refinement, presentation change, or
gluing claim is being made by the source record.  The nonclaim that the
source is not derived inside this paper is recorded in
\Cref{nonclaim:pvnp:hiddenness-derivable}.\footnote{Verified in Lean
by the recognition-source discharge declaration listed in
\Cref{app:formalization}; the Lean theorem consumes the
\texttt{recognitionSource} carrier and rewrites through its
\texttt{not\_currentize\_wit} field.  The
narrowing-disclosure discipline is recorded in
\Cref{app:formalization}.}
\end{proof}

\begin{theorem}[Translation-interface discharge]
\label{thm:pvnp:aor:translation-discharge}
Translation Theorems A and B have primary residual type
\(\Delta_{\mathrm{transport}}\).  Their forced secondary types are
\(\Delta_{\mathrm{role}}\), \(\Delta_{\mathrm{target}}\), and
\(\Delta_{\mathrm{source}}\).  Translation A is discharged by the
standard \(P\)-membership-to-current-observable bridge; Translation B
is discharged by the SAT self-reduction and witness-currentization
bridge using the closure properties of \(\EQbang\).
\end{theorem}

\begin{proof}
Translation A transports the standard assertion \(\mathrm{SAT}\in P\)
to the internal predicate \(\CurrentizeDec(\mathrm{SAT})\).  The
source is the standard polynomial-time decider package, the target is
decision currentization, and the role is a translation bridge rather
than an independent structural source.  Translation B transports
\(\CurrentizeDec(\mathrm{SAT})\) to
\(\CurrentizeWit(\OmegaSAT)\).  The source is the decision
currentizer, together with the SAT self-reduction wrapper; the target
is witness currentization for the autonomous lift; the transport is
lawful because \(\EQbang\) is closed under polynomial-time
post-processing and polynomial-time composition
(\Cref{thm:pvnp:e-q-closure-properties}, clauses~(1) and~(3)).  The
self-reduction wrapper is an adaptive sequential composition of
polynomial-time input-restriction transformations with the decision
transducer, followed by polynomial-time post-processing to assemble
the canonical assignment; the lawfulness record for the wrapper is
verified explicitly in the proof of \Cref{thm:translation-B}.  Thus the
two translation residuals are discharged by \(\texttt{bridged}\)
transport records, with their target and role components typed by
\Cref{def:pvnp:currentize-dec,def:pvnp:currentize-wit}.\footnote{Verified
in Lean by the translation-discharge declaration listed in
\Cref{app:formalization}; the Lean theorem inherits the
narrowed-surrogate iff-packaging of
\texttt{translationA} (parametric
\(\mathrm{SAT}\in P\)-to-current-observable bridge) and
\texttt{translationB} (parametric
self-reduction wrapper and closure-correctness contracts).  The
narrowing-disclosure discipline is recorded in
\Cref{app:formalization}.}
\end{proof}

\begin{theorem}[Class-B scope-fence discharge]
\label{thm:pvnp:aor:class-b-scope-fence}
The Class-B Irremovability theorem
\(\Cref{thm:pvnp:irremovability-scoped}\) has primary residual type
\(\Delta_{\mathrm{horizon}}\).  Its forced secondary types are
\(\Delta_{\mathrm{role}}\), \(\Delta_{\mathrm{target}}\), and
\(\Delta_{\mathrm{local/global}}\).  Its status is
\(\texttt{blocked}\) for basis-generated Class-B laws on the standard
fixed-presentation carrier and
\(\texttt{out\_of\_scope}/\texttt{nonclaim}\) for source-removal
methods outside that class.
\end{theorem}

\begin{proof}
The residual being fenced is not \(\CSLSATHidden\) itself, but the
attempt to remove it by intrinsic carrier-feature reasoning over the
standard fixed-presentation polynomial current carrier.  The Class-B
basis \(\mathcal B_{\mathrm{fix}}\), the normal-class analysis, and
the irremovability core identify the carrier-side horizon of this
method class: quotient features collapse, joint-access obstructions
are killed by pairing closure, and the remaining basis features do
not exclude the fixed-branch residual.  Inside that scoped method
class, the removal route is blocked.  Outside that class, the paper
makes no absolute claim; this is recorded in
\Cref{nonclaim:pvnp:absolute-irremovability}.  A positive
horizon-sufficiency residual would carry interface/descent ancestry,
but this theorem is not such a claim.  It is a method-class
scope-fence.  Therefore the forced secondary package is exhausted by
the method-role record, the removal-target record, and the
local/global basis-exhaustiveness record.  No further
\(\Delta_{\mathrm{interface}}\), \(\Delta_{\mathrm{descent}}\),
\(\Delta_{\mathrm{scale}}\), or \(\Delta_{\mathrm{limit}}\) payload is
generated.\footnote{Verified in Lean as
the Class-B scope-fence discharge declaration listed in
\Cref{app:formalization}; the Lean theorem invokes
\texttt{irremovabilityScoped} and inherits the feature-classification
narrowing.}
\end{proof}

\begin{theorem}[BirdInt package discharge]
\label{thm:pvnp:aor:judgment-package}
The four-step contradiction chain and the \(\BirdInt\) judgment of
\Cref{sec:contradiction-chain,sec:judgment} have primary residual type
\(\Delta_{\mathrm{role}}\).  Their forced secondary types are
\(\Delta_{\mathrm{source}}\), \(\Delta_{\mathrm{target}}\), and
\(\Delta_{\mathrm{meta}}\).  After the discharges in
\Cref{thm:pvnp:aor:recognition-source-discharge,thm:pvnp:aor:translation-discharge}, the judgment package is
\(\texttt{zero}\) as a bookkeeping transport from the conditional
contradiction to
\[
  \BIJudgment{\CSLSATHidden}{\Tsat}{\Ipoly}{\PSB\neq\NPSB}.
\]
\end{theorem}

\begin{proof}
The contradiction chain consumes three typed inputs: Translation A,
Translation B, and the recognition-source negation.  It produces
\(\mathrm{SAT}\notin P\) under \(\CSLSATHidden\).  The verifier
package supplies \(\mathrm{SAT}\in NP\).  The \(\BirdInt\) judgment
does not change these claims; it assigns them their framework roles:
source \(\CSLSATHidden\), layer \(\Tsat\), instrument \(\Ipoly\), and
proposition \(\PSB\neq\NPSB\).  Recognition grade records the
meta-status of the judgment: the closure apparatus is framework
hygiene, while the SAT-specific hiddenness content is supplied by the
recognition source.  No further residual remains after the roles are
assigned.\footnote{Verified in Lean by the BirdInt judgment-package
discharge declaration listed in \Cref{app:formalization}; the Lean
theorem inherits the narrowed-surrogate status of \texttt{satInNP}
(verifier-package transfer from textbook hypotheses) and consumes the
parametric contradiction-chain conclusion supplied by
\texttt{fourStepContradictionConditional}.  The narrowing-disclosure
discipline is recorded in \Cref{app:formalization}.}
\end{proof}

\subsection{AOR-instance theorem}

\begin{theorem}[PvNP AOR-instance theorem]
\label{thm:pvnp:aor:in-aor-casc}
Under the records and discharges of
\Cref{def:pvnp:aor-instance:object,thm:pvnp:aor:mechanical-records,thm:pvnp:aor:recognition-source-discharge,thm:pvnp:aor:translation-discharge,thm:pvnp:aor:class-b-scope-fence,thm:pvnp:aor:judgment-package}, the saturated SAT carrier is a
cascade-stable AOR instance on the PvNP scope:
\[
  \Tsat \in
  \mathbf{AOR}^{\mathrm{casc}}_{\mathfrak S_{\mathrm{PvNP}}}.
\]
The record-level reading is the same as the theorem of
\Cref{sec:judgment}: the framework-relative judgment
\[
  \BIJudgment{\CSLSATHidden}{\Tsat}{\Ipoly}{\PSB\neq\NPSB}
\]
is accepted at \RecognitionGrade{}.
\end{theorem}

\begin{proof}
The saturated carrier, instrument, observable families, quotients,
comparison map, autonomous lift, visibility relation, and
currentization predicates are supplied by construction.  The
operation closure of \(\EQbang\), the SAT verifier package, and the
standard polynomial-time computation facts are recorded at their
declared admitted status.

We now spell out the finite discharge walk.  First, the recognition
source starts as a \(\Delta_{\mathrm{source}}\) residual.  The
\(\texttt{hidden}\) status generates the target, role, and horizon
costs described in
\Cref{thm:pvnp:aor:recognition-source-discharge}.  The target is
\(\neg\,\CurrentizeWit(\OmegaSAT)\), already typed by
\Cref{def:pvnp:currentize-wit}; the role is recognition-source, fixed
by the source record; the horizon is the current/autonomous boundary
\((\EQbang,\OmegaSAT)\), fixed by \(\VisBang\) and closed by the
same hiddenness record.  The only remaining cost is the nonclaim of
internal derivation, recorded in
\Cref{nonclaim:pvnp:hiddenness-derivable}; the branch then reaches the
empty record.

Second, Translation A and Translation B start as
\(\Delta_{\mathrm{transport}}\) residuals.  The
\(\texttt{bridged}\) status generates source, role, and target costs.
For Translation A, the source is the assumed standard
polynomial-time decider under the contradiction hypothesis, the role
is translation, and the target is \(\CurrentizeDec(\mathrm{SAT})\).
For Translation B, the source is the decision currentizer plus the
SAT self-reduction package, the role is witness-currentization
transport, and the target is \(\CurrentizeWit(\OmegaSAT)\).  The
closure properties of \(\EQbang\) discharge the transport costs by
post-processing, finite product, and composition records.  The
standard SAT self-reduction and verifier components are covered by
the approved computational records of
\Cref{thm:pvnp:aor:mechanical-records}.  The translation branch then
reaches the empty record.

Third, the Class-B scope fence starts as a \(\Delta_{\mathrm{horizon}}\)
residual for one removal route.  Inside the Class-B carrier-feature
method class, \(\Cref{thm:pvnp:irremovability-scoped}\) gives a
\(\texttt{blocked}\) discharge.  This generates method-role,
removal-target, and local/global basis-exhaustiveness costs.  These
are discharged by the Class-B basis, orbit-collapse, basis
exhaustiveness, basis-combination closure, and irremovability-core
records.  Outside that method class the cost-table row is
\(\texttt{out\_of\_scope}/\texttt{nonclaim}\), closed by
\Cref{nonclaim:pvnp:absolute-irremovability}.  No other secondary
obligation remains.

Fourth, the \(\BirdInt\) package starts as a \(\Delta_{\mathrm{role}}\)
residual: the source, layer, instrument, target proposition, and
judgment status must be assigned.  The preceding three branches
supply the source-side and translation-side content.  The layer is
\(\Tsat\), the instrument is \(\Ipoly\), and the target proposition is
\(\PSB\neq\NPSB\).  The contradiction chain supplies the proposition
under the recognition source, and the verifier package supplies
\(\mathrm{SAT}\in NP\).  The \(\BirdInt\) record therefore discharges
at \(\texttt{zero}\), with recognition grade recorded as the
meta-status of the judgment.

Every residual native to the proof stack is thus
\(\texttt{by\_construction}\), \(\texttt{approved\_other}\),
\(\texttt{zero}\), \(\texttt{hidden}\), \(\texttt{bridged}\),
\(\texttt{blocked}\), or
\(\texttt{out\_of\_scope}/\texttt{nonclaim}\).  Each branch terminates
after finitely many cost-table moves, and every cross-cascade
obligation is served by one of the discharges named above.  No
asymptotic refinement payload is used.  Hence the carrier belongs to
\(\mathbf{AOR}^{\mathrm{casc}}\) on the declared PvNP scope.\footnote{Tracked
in the formalization harness; see \Cref{app:formalization}.  The
\(\mathbf{AOR}^{\mathrm{casc}}\) tag is carried as a local field on the
\texttt{aorInstanceObject} rather than as membership in an imported
predicate.  The projection-packaging discipline is recorded in
\Cref{app:formalization}.}
\end{proof}
 

\section{Scope and non-claims}
\label{sec:scope}

The preceding sections state a framework-relative judgment at
\RecognitionGrade{} and the conditional theorem it reads as outside the
framework.  This section makes the boundary explicit, so that
\(\PSB\neq\NPSB\) is not mistaken for an unconditional standard-ZFC
theorem about \(P\) and \(NP\).  The scope of the Class-B Irremovability
barrier proved in \Cref{sec:class-b-barrier} is recorded alongside the
recognition-grade scope of the first two results.

Positively, the paper claims the following.  The saturated SAT layer
construction, including \(\Tsat\), \(\EQbang\), \(\Msat\),
\(\Qsat\), and \(\OmegaSAT\), is standard mathematics expressed in
the framework's typed vocabulary.  The closure properties of
\(\EQbang\) are the current-observable closure rules needed by the
translation theorems.  Translation A identifies standard
\(\mathrm{SAT}\in P\) with \(\CurrentizeDec(\mathrm{SAT})\), and
Translation B identifies \(\CurrentizeDec(\mathrm{SAT})\) with
\(\CurrentizeWit(\OmegaSAT)\).  Conditional on the structural
recognition source \(\CSLSATHidden\), the four-step contradiction
gives \(\mathrm{SAT}\notin P\).  The BirdInt judgment records
\(\PSB\neq\NPSB\) at \RecognitionGrade{} under the standard Six Birds
closure assumption.  Read as ordinary mathematics, the result is the
conditional theorem
\[
  \text{formed-SAT-closure-under-\(\Ipoly\)}
  \quad\Longrightarrow\quad
  P\neq NP.
\]
The paper also proves the Six Birds Class-B Irremovability barrier on
the standard fixed-presentation polynomial current carrier
(\Cref{thm:pvnp:irremovability-scoped}).  The localization in
\Cref{sec:recognition-source} identifies the concrete residual on
which the barrier acts:
\(\mathrm{Br}^{\mathrm{lex}}_{\mathrm{SAT}}\notin\EQbang\).

\paragraph{No standard-ZFC unconditional separation.}
\label{nonclaim:pvnp:standard-zfc-pneqnp}
This paper does not claim a standard-ZFC unconditional separation of
\(P\) from \(NP\) \citep{Cook2000PNPClay,Fortnow2009StatusPNP,Fortnow2022FiftyYears}.  The headline judgment \(\PSB\neq\NPSB\) is a
framework-relative judgment at \RecognitionGrade{} under the supplied
structural recognition source.  Outside the framework, what is
established is the conditional theorem just displayed.  Without
\(\CSLSATHidden\), the contradiction chain has no source for
\(\neg\CurrentizeWit(\OmegaSAT)\), and no separation of standard
\(P\) from \(NP\) is proved.

\paragraph{No classical-barrier circumvention.}
\label{nonclaim:pvnp:circumvent-classical-barriers}
This paper does not claim to circumvent natural proofs
(Razborov--Rudich, \citealp{RazborovRudich1997NaturalProofs}),
algebrization (Aaronson--Wigderson,
\citealp{AaronsonWigderson2009Algebrization}), or relativization
(Baker--Gill--Solovay, \citealp{BakerGillSolovay1975Relativization}) in
their native senses.  The construction operates at a different layer:
typed observable saturation over a SAT-specific history carrier,
together with a structural recognition source.  The Class-B barrier of
\Cref{sec:class-b-barrier} is analogous to the three classical barriers
in role but methodologically distinct and internal to Six Birds; the
detailed comparison is in \Cref{subsec:class-b:comparison}.

\paragraph{Recognition source not derivable from framework primitives.}
\label{nonclaim:pvnp:hiddenness-derivable}
This paper does not claim that the structural recognition source
\(\CSLSATHidden\) is derivable from the framework's primitive
constructions.  \citet{TsiokosHiddenness2026} documents the relevant
no-go: the SAT-specific source is tested against three derivation
routes — a framework-internal derivation, a reduction from a more
abstract structural result, and a finite-carrier witness construction
— and all three routes fail.  The source is therefore supplied to the
present construction as structural content from that paper, not derived
here.

\paragraph{No absolute irremovability of the recognition source.}
\label{nonclaim:pvnp:absolute-irremovability}
The Class-B Irremovability theorem
(\Cref{thm:pvnp:irremovability-scoped}) is scoped.  It rules out
basis-generated Class-B carrier-feature laws on the standard
fixed-presentation polynomial current carrier; it does not claim
absolute irremovability of \(\CSLSATHidden\) against all future Six
Birds laws.  In particular, it does not preclude laws using new
carrier structure outside \(\mathcal B_{\mathrm{fix}}\), Class-B laws
on enriched carriers (objective SAT, branchwise carriers,
presentation-quotiented carriers), or Class-A claim-discipline laws
acting on the same recognition source.

\paragraph{AOR-instance reading restricted to the formed SAT layer.}
\label{nonclaim:pvnp:aor-instance-formed-layer-only}
The AOR-instance reading of \Cref{sec:pvnp-aor-instance}, including
\Cref{thm:pvnp:aor:in-aor-casc}, is restricted to the saturated SAT
layer \(\Tsat\) on the standard fixed-presentation polynomial current
carrier under the SAT-specific recognition source
\(\CSLSATHidden\).  No AOR-membership claim is made for enriched
carriers, alternative recognition sources, Class-A claim-discipline
methods, or any closure scope outside \(\mathfrak S_{\mathrm{PvNP}}\).
The Class-B \(\Delta_{\mathrm{horizon}}\) residual stratum is
explicitly statused as \texttt{out\_of\_scope/nonclaim} in
\Cref{thm:pvnp:aor:class-b-scope-fence}.

\paragraph{No independent unconditional \(P\neq NP\) claim via AOR instance.}
\label{nonclaim:pvnp:aor-instance-no-unconditional}
\Cref{thm:pvnp:aor:in-aor-casc} establishes that the saturated SAT
carrier \(\Tsat\) is tagged as cascade-stable AOR under the listed
discharges, with the transported propositions \(\mathrm{SAT}\notin P\)
and \(\mathrm{SAT}\in NP\) exposed in the conclusion.  Independently
of the audit records ---
in particular the recognition source \(\CSLSATHidden\) and the
six declared per-discharge structures --- the result is not an
unconditional standard-ZFC \(P\neq NP\) claim; the AOR-instance
reading reframes the closure stack as the audit-closure membership
witness but does not upgrade the conditional separation to an
unconditional one outside the audit-record presentation.

These six nonclaims distinguish what the paper proves from what
would be required to upgrade the result to an unconditional
standard-mathematics theorem, or to a result that addresses the
classical barriers in their native formulations.  The fourth nonclaim
is the scope fence for the Class-B Irremovability barrier theorem of
\Cref{sec:class-b-barrier}; the fifth and sixth are the scope fences
for the AOR-instance recasting of
\Cref{sec:pvnp-aor-instance}.  The discussion section spells out the
technical layer at which the construction operates and what additional
bridge would be required to strengthen the result beyond the
conditional theorem stated here.
 

\section{Discussion}
\label{sec:discussion}

The preceding sections establish a conditional theorem and package it
as a framework-relative judgment.  This discussion positions that
result against the broader landscape: the distinction between
construction-grade and recognition-grade content, the granularity
mismatch inherited from \citet{TsiokosHiddenness2026}, the
methodological situation after the Class-B Irremovability barrier of
\Cref{sec:class-b-barrier}, and what the conditional theorem plus that
barrier leave open.  The detailed comparison with the classical
barriers is now in \Cref{subsec:class-b:comparison}.  The scope fences
of \Cref{sec:scope} remain in force throughout.

\paragraph{Construction grade vs recognition grade.}
The construction-grade content of the paper is the part built directly:
the saturated SAT layer, observable saturation and its closure
properties, Translation Theorem A, Translation Theorem B, the
conditional four-step contradiction, and the standard verifier package
for \(\mathrm{SAT}\in NP\).  These are ordinary mathematical
constructions expressed in the framework's typed vocabulary.  The
recognition-grade content is different.  It is the \(\BirdInt\)
judgment that packages the conditional theorem into the
framework-relative assertion \(\PSB\neq\NPSB\) at
\RecognitionGrade{}, using the supplied recognition source
\(\CSLSATHidden\).  The finite-audit judgment apparatus is the
Foundations-III setting for such records \citep{TsiokosFoundationsIII2026}.
This distinction is the reason the paper is not presented as a
candidate unconditional proof: the contradiction chain is theorem-grade
conditional on the source, while the final judgment is recognition
grade relative to that source.

\paragraph{Granularity-mismatch bridge.}
\citet{TsiokosHiddenness2026} establishes a granularity-mismatch
theorem: framework current and predictive quotients are finer than the
standard language classes \(P\) and \(NP\) at the history level, and
the strictness can be lost under standard language-level projection.  This is the structural reason the
present paper's conditional theorem does not automatically project to
an unconditional standard theorem \(P\neq NP\).  The framework
operates over typed quotients of histories and observables.  Standard
complexity theory operates over languages and verifier predicates.  The
projection from the former to the latter discards trace identity,
machine bookkeeping, and witness-identity information.  The
conditional theorem is therefore the readout available before
crossing that projection.

\paragraph{Classical-barrier positioning.}
Relativization, natural proofs, and algebrization live at a different
methodological level from the present construction: they concern
oracle extension, pseudorandomness and natural circuit properties, and
arithmetized oracle access, respectively.  The present paper instead
works with typed observable saturation, lawfulness records, and a
structural recognition source on the formed SAT layer.  The detailed
comparison is given in \Cref{subsec:class-b:comparison}.  The headline
there is that Six Birds Class-B Irremovability is analogous in role to
the classical barriers, but methodologically distinct and internal to
the Six Birds framework.

\paragraph{Why the result lands at recognition grade.}
The conditional theorem lands at \RecognitionGrade{} because it depends
on the structural source \(\CSLSATHidden\).  One can therefore ask
whether the source is removable: whether a framework-internal
derivation can replace it.  \Cref{sec:class-b-barrier} answers this
question for one methodological class.  Intrinsic carrier features of
the standard fixed-presentation polynomial current carrier are
basis-generated (\Cref{thm:pvnp:basis-exhaustiveness}), basis-generated
features reduce to the five normal classes
(\Cref{thm:pvnp:basis-combination-closure}), and no basis-generated
Class-B carrier-feature law removes \(\CSLSATHidden\)
(\Cref{thm:pvnp:irremovability-scoped}).  The Class-B route to removing
the source is therefore structurally closed.  The explanation is
scoped: Class-A claim-discipline methods, Class-B laws on enriched
carriers, and methods using new carrier structure outside
\(\mathcal B_{\mathrm{fix}}\) remain open.

\paragraph{Standard constructions used.}
Apart from the recognition source, the mathematical ingredients are
standard.  Translation B uses the usual SAT self-reducibility argument
\citep{Schoening1986ComplexityStructure,HemaspaandraOgihara2002SelfReducibility}:
a decider is queried on restricted formulas to recover a witness
variable by variable.  The observable closure theorem uses standard
polynomial-time composition, finite products, and post-processing.
The quotients are ordinary quotients by observation equivalence
relations.  The proof that \(\mathrm{SAT}\in NP\) is the standard
verifier package.  The paper's contribution is not a new computational
primitive; it is the typed saturated-SAT-layer organization, the
translation theorems, and the recognition-grade conditional packaging.

\paragraph{Why the conditional reading is primary.}
The source \(\CSLSATHidden\) has structural content: it says that
witness data are hidden from the current layer on the formed SAT
carrier under \(\Ipoly\).  \citet{TsiokosHiddenness2026} records that
the SAT-specific instance is not derivable from the framework's
primitive constructions.  A reader who accepts the framework's closure assumption
and finds the supplied source structurally compelling may read the
result as the \(\BirdInt\) judgment \(\PSB\neq\NPSB\) at
\RecognitionGrade{}; a reader who does not make that commitment reads
the same mathematics as a conditional theorem with explicit antecedent
\(\CSLSATHidden\).  These are two readings of the same conditional
structure, not two different mathematical claims.  The Class-B
Irremovability barrier adds a further reason for the conditional
reading: the carrier-feature route for removing the antecedent is itself
blocked within the scoped methodological class.

\paragraph{Open directions.}
Several questions remain open.  One may ask whether the
recognition source extends cleanly to other NP-complete carriers, or
whether SAT is the simplest instance of a more abstract structural
hypothesis.  One may study additional closure properties of
\(\EQbang\), beyond post-processing, finite products, and composition,
and ask how they further constrain the current-versus-predictive
distinction.  One may also investigate whether weakening the standard
closure assumption changes the \(\BirdInt\) judgment while leaving the
ordinary conditional theorem intact.  The Class-B barrier suggests
three further directions: Class-B laws on enriched carriers such as
objective SAT, branchwise SAT, or presentation-quotiented SAT;
Class-A claim-discipline approaches to the recognition source, which
operate on claim shape rather than carrier features; and possible
extensions of \(\mathcal B_{\mathrm{fix}}\) by new intrinsic features
native to the standard polynomial-time carrier, together with the
corresponding barrier theorem on the extended basis.  These are
framework-extension questions, not changes to the conditional theorem
proved here.
 

\section{Conclusion}
\label{sec:conclusion}

The paper has proved three linked results.  First, it constructs the
saturated SAT layer \(\Tsat\), the lawful polynomial-time current
instrument \(\Ipoly\), and the saturated current-observable family
\(\EQbang\), then proves Translation Theorem A and Translation Theorem
B.  Conditional on the structural recognition source
\(\CSLSATHidden\), the four-step contradiction gives
\(\mathrm{SAT}\notin P\); combined with the standard verifier-package
fact \(\mathrm{SAT}\in NP\) \citep{Cook1971TheoremProving}, this yields
the framework-relative \(\BirdInt\) judgment \(\PSB\neq\NPSB\) at
\RecognitionGrade{}.  Outside the framework, this is the conditional
theorem that formed SAT closure under \(\Ipoly\) implies \(P\neq NP\).
Second, \Cref{sec:recognition-source} makes the structural source
concrete by localizing \(\CSLSATHidden\) to the fixed-branch residual
\[
  \mathrm{Br}^{\mathrm{lex}}_{\mathrm{SAT}}\notin\EQbang,
\]
the assertion that the canonical lexicographic SAT-branching readout is
not a lawful current observable on the fixed-presentation carrier.
Third, \Cref{sec:class-b-barrier} proves the Six Birds Class-B
Irremovability barrier: no basis-generated Class-B carrier-feature law
on the standard fixed-presentation polynomial current carrier can
remove \(\CSLSATHidden\).  That barrier rests on Fixed-Presentation
Orbit Collapse (\Cref{thm:pvnp:orbit-collapse}) and polynomial
pairing closure, organized through Basis Exhaustiveness
(\Cref{thm:pvnp:basis-exhaustiveness}), Basis Combination Closure
(\Cref{thm:pvnp:basis-combination-closure}), and the scoped
Irremovability theorem (\Cref{thm:pvnp:irremovability-scoped}).  It
is analogous in role to Baker--Gill--Solovay relativization,
Razborov--Rudich natural proofs, and Aaronson--Wigderson algebrization
\citep{BakerGillSolovay1975Relativization,RazborovRudich1997NaturalProofs,AaronsonWigderson2009Algebrization},
but methodologically distinct and internal to Six Birds.

The reason for three results rather than one is methodological.  The
conditional theorem establishes the recognition-grade judgment.  The
localization identifies what the recognition source contributes — the
fixed-branch residual
\(\mathrm{Br}^{\mathrm{lex}}_{\mathrm{SAT}}\notin\EQbang\) — and gives
the Class-B barrier a concrete target.  The Class-B barrier then
explains why the carrier-feature route for removing that source is
closed on the standard polynomial current carrier.  Read as a single
statement: \(P\) versus \(NP\) at recognition grade, plus a
framework-internal barrier against one natural class of attempts to
upgrade beyond it.

The load-bearing structural research problem remains the derivation of
the recognition source.  The Class-B barrier provides a partial negative
answer: the source cannot be removed by basis-generated Class-B
carrier-feature methods on the standard fixed-presentation polynomial
current carrier.  \citet{TsiokosHiddenness2026} documents that three
further routes fail for the SAT-specific source — a framework-internal
derivation, a reduction from a more abstract structural result, and a
finite-carrier witness construction.  The broader question stays open:
Class-A claim-discipline methods, Class-B laws on enriched carriers,
and methods using carrier structure outside
\(\mathcal B_{\mathrm{fix}}\) are not addressed here.

The contribution is therefore a typed structural calculation under an
explicit hypothesis, a precise localization of that hypothesis, and a
framework-internal barrier theorem against one class of attempts to
remove it.  Whether the hypothesis can be made unconditional is a
separate structural question; for basis-generated Class-B laws on the
standard carrier the answer is negative, while for Class-A methods,
enriched carriers, and non-basis structures it remains open.

The conditional theorem of \Cref{sec:judgment} packages the
construction as the framework-relative \(\BirdInt\) judgment
\(\PSB\neq\NPSB\) at \RecognitionGrade{}; equivalently, by
\Cref{sec:pvnp-aor-instance}, the same closure records exhibit
\(\Tsat\) as an audit-closed carrier under \(\CSLSATHidden\) in the
sense of \citet{Tsiokos2026AOR}.  Both readings share the same
audit-record requirement.
 
\appendix


\section{Primitive-role schema: six mechanisms}
\label{app:primitive-roles}

This appendix records the six-mechanism schema used by the saturated
SAT layer.  In \Cref{sec:saturated-closure}, the audited shell
\(\Saud\) is introduced as the finite-stage state record on which the
SAT construction runs; here we spell out how that shell instantiates
the primitive-role schema \(P_{1}\)--\(P_{6}\).  The appendix is
therefore a bookkeeping appendix for the construction, not a second
definition of the layer \(\Tsat\).

The Foundations-II primitive-role schema provides six typed roles that
organize a closure construction into named components rather than
implicit ambient choices: machine update, admissibility, timescale,
lens selection, packaging, and audit \citep{TsiokosFoundationsII2026}.
The SAT instance of the schema names these roles on the history
carrier \(\Hsat\), under the lawful polynomial-time instrument
\(\Ipoly\) and the saturated current-observable family \(\EQbang\).

\begin{definition}[Six mechanisms (P1--P6)]
\label{def:pvnp:six-mechanisms}
The saturated SAT shell carries the following six mechanisms.
\begin{enumerate}
  \item \(P_{1}\), \emph{machine update}.  This is the deterministic
  polynomial-time transducer class admitted on the SAT layer, namely
  the class \(K_{\mathrm{poly}}\) declared in the audited shell.
  Update is by lawful polynomial-time machine action on encoded SAT
  instances.

  \item \(P_{2}\), \emph{admissibility}.  This is the set of encoded
  SAT instances admitted at the current finite stage.  Admissibility is
  determined by the finite-stage restriction of the audited shell's
  instance database, together with the type-correctness and
  no-smuggling checks imposed by the lawfulness discipline.

  \item \(P_{3}\), \emph{timescale}.  This is the polynomial budget
  declared by the lawful instrument \(\Ipoly\), including the local
  clock and the time certificate recorded by each lawful current
  observable.

  \item \(P_{4}\), \emph{lens selection}.  This is the visibility lens
  of the audited shell, determining which projection of histories is
  exposed to current observation, and which wider projection is used
  for predictive observation.

  \item \(P_{5}\), \emph{packaging}.  This is the \(Q/M\) quotient and
  continuation packaging: the SAT-saturated current and predictive
  quotients \(\Qsat\) and \(\Msat\), together with the comparison map
  \(\ComparisonMap\) on the history carrier \(\Hsat\).

  \item \(P_{6}\), \emph{audit}.  This is the audit ledger of the
  audited shell, including the machine-index record, time certificate,
  replay log, and gate predicates declared in the audit interface.  The
  closed audit record is part of the saturated SAT layer \(\Tsat\), and
  the seven gates \(G_{1}\)--\(G_{7}\) are listed in
  \Cref{app:gate-audit}.
\end{enumerate}
\end{definition}

This schema instance is mathematically content-free as a labeling: it
names the six components without adding material beyond the body
definitions.  Its purpose is to make the typed organization explicit.
A reader who wants to check how the SAT construction maps onto the
Foundations-II primitive roles can read \(P_{1}\)--\(P_{6}\) directly
against the shell \(\Saud\), the carrier \(\Hsat\), the instrument
\(\Ipoly\), and the observable family \(\EQbang\).  The same labeling
is used uniformly across Foundations-II closure constructions; the SAT
instance is the natural one once the substrate, lift, instrument,
observable families, quotient packaging, and audit ledger have been
fixed.

The formalization records this schema instance as a typed-record
declaration.  The coverage table in \Cref{app:formalization} lists the
corresponding row and the per-row formalization status.
 

\section{Recognition-source carrier schema}
\label{app:recognition-source-carrier}

This appendix records the carrier schema for the recognition source
\(\CSLSATHidden\).  \Cref{sec:recognition-source} gave the conceptual
content: witness currentization fails on the
formed saturated SAT layer under the lawful current instrument.  The
present appendix gives the typed-record carrier by which that content
is tracked in the formalization.  The source is encoded as a structural
recognition-source carrier with explicit fields; it is not introduced
as a primitive assertion of the proof assistant.

\begin{center}
\small
\begin{tabularx}{\linewidth}{@{}>{\raggedright\arraybackslash}p{0.32\linewidth}>{\raggedright\arraybackslash}X@{}}
\toprule
Field & Intended semantics \\
\midrule
Recognition-content predicate &
A typed predicate over the saturated SAT layer asserting that, on
\(\Tsat\) under \(\Ipoly\), no lawful current observable in
\(\EQbang\) realizes the canonical witness producer generated by
\(\OmegaSAT\).  Equivalently, this is the negation of
\(\CurrentizeWit(\OmegaSAT)\). \\

Source identity &
A typed name slot recording the source as \(\CSLSATHidden\).  This
identifier ties the carrier to its origin in
\citet{TsiokosHiddenness2026} and to the V-Differential placement of
PvNP as trace-only-exposing. \\

Closure context &
A typed context slot recording that the source is invoked inside the
framework's audited closure shape, with layer \(\Tsat\) and instrument
\(\Ipoly\) fixed.  It is the place where the standard closure
assumption applies to the carrier use. \\

Provenance record &
A typed provenance slot recording the structural lineage: the
\CSLcore{} of \citet{TsiokosHiddenness2026}, the V-Differential
placement of PvNP, and the documented failure of the natural derivation
routes.  Provenance is bookkeeping; it does not make the source
derivable from the present construction. \\

Audit boundary &
A typed audit slot recording the no-smuggling commitments attached to
the carrier use: no SAT oracle, no witness oracle, no hidden SAT bit,
and no invocation outside the saturated SAT layer. \\
\bottomrule
\end{tabularx}
\end{center}

The carrier discipline is more explicit than an informal appeal to
hiddenness.  The recognition source is realized as a typed record:
every conditional use in the four-step contradiction chain is a typed
call on the carrier predicate, together with the surrounding source,
closure-context, provenance, and audit bookkeeping.  The formal
development uses no primitive assertion to introduce this source.
Its axiom audit remains under the same trust base used elsewhere in
the project---Classical choice, propositional extensionality,
quotient soundness, choice, and functional extensionality
\citep{MartinLof1984IntuitionisticTypeTheory}.

The schema is consistent with the content of
\citet{TsiokosHiddenness2026}.  The recognition-content predicate is
the SAT-specific instantiation of the \CSLcore{} developed there: on
the formed SAT layer, witness content generated by \(\OmegaSAT\) is
not exposed to lawful current observation under \(\Ipoly\).  This
appendix is the disclosure point for the paper's source posture.  The
source is supplied, named, audited, and carried through the
conditional theorem; it is not derived in the present construction.
 

\section{Gate audit (six gates plus Gate 7)}
\label{app:gate-audit}

This appendix records the seven-gate audit used by the lawfulness
discipline in \Cref{sec:observable-saturation}.  A lawful current
observable in \(\EQbang\) is not merely a polynomial-time transducer;
it carries audit data certifying that the transducer has not imported
the SAT bit, a witness, an external verifier, or other hidden
information.  The seven gates below are the typed check predicates
attached to that audit data.

Each gate is represented in the audit record by the same four-part
schema.  First there is a check predicate, a typed assertion about the
candidate observable or the saturated layer on which it acts.  Second
there is a status saying whether the check passes.  Third there is an
evidence record witnessing the status.  Fourth there is an implication
from the passed check to the corresponding no-smuggling or
no-overreading property.  The same schema is used for
\(\GateOne\)--\(\GateSix\) and for the seventh gate.

\begin{proposition}[Gate audit]
\label{prop:pvnp:gate-audit}
Every lawful current observable in \(\EQbang\), acting on the
saturated SAT layer \(\Tsat\), passes the following seven gate checks
with the typed evidence records carried in its lawfulness data.
\begin{enumerate}
  \item \(\GateOne\), \emph{no SAT oracle smuggling}.  The check
  predicate says that the observable does not invoke an oracle
  \citep{BakerGillSolovay1975Relativization} for the SAT decision
  predicate.  The status passes by construction: the audited shell
  admits ordinary deterministic polynomial-time transducers, not
  SAT-oracle machines.
  The evidence is the audit ledger recording the absence of oracle
  commitments.  The implication is that the observable's computation
  does not implicitly use the SAT bit as input.

  \item \(\GateTwo\), \emph{no witness oracle smuggling}.  The check
  predicate says that the observable does not invoke an oracle for a
  satisfying-witness producer \citep{FeigeShamir1990WitnessHiding,GoldwasserMicaliRackoff1989KnowledgeComplexity}.  The status passes by construction, and
  the evidence is again the audit ledger of the admitted transducer.
  The implication is that satisfying assignments are not pre-loaded
  into the observable or made available as hidden input.

  \item \(\GateThree\), \emph{no external verifier oracle smuggling}.
  The check predicate says that the observable does not invoke an
  external NP-verifier oracle.  The status passes by the same audit
  evidence: verifier behavior must be realized by the declared
  polynomial-time computation itself.  The implication is that the
  observable cannot test whether an assignment satisfies a formula by
  an undeclared oracle call.

  \item \(\GateFour\), \emph{no hidden \(B\)-bit smuggling}.  The
  check predicate says that the visibility lens does not encode a
  hidden bit, such as a disguised SAT bit, into the \(Q\)-visible
  package data
  \citep{GoguenMeseguer1982SecurityPolicies,SabelfeldMyers2003InformationFlow}.  The status passes by visibility bookkeeping.  The
  evidence is the visibility-lens record attached to the current
  instrument.  The implication is that the \(Q\)-visible projection
  does not encode predictive or witness-side data.

  \item \(\GateFive\), \emph{no time-certificate smuggling}.  The
  check predicate says that the polynomial time certificate is honest:
  the observable does not exceed its declared polynomial budget by
  hiding deferred evaluation or unrecorded subcomputation.  The status
  passes by the time-certificate clause of the lawfulness discipline.
  The evidence is the certificate itself, together with the replay
  record.  The implication is that the observable's polynomial budget
  is the budget actually audited.

  \item \(\GateSix\), \emph{no scope violation}.  The check predicate
  says that the observable's domain is the all-input replay
  convention's instance set, not a favorable subfamily of SAT
  instances.  The status passes by the all-input replay convention.
  The evidence is the replay record.  The implication is that the
  observable cannot be a partial decider presented as a total current
  observable.

  \item \(\GateSeven\), \emph{no overreading suppression}.  The check
  predicate says that the observable does not suppress overreading of
  its own audit data: the lawfulness record is faithful to the actual
  computation being audited.  The status passes by the Foundations-III
  no-overreading-suppression theorem
  \citep{TsiokosFoundationsIII2026}.  The evidence is the audit
  witness field.  The implication is that the audit ledger's records
  cannot be falsified or selectively hidden at audit time.
\end{enumerate}
\end{proposition}

The gate audit is verified in the formalization as a typed projection
over the audit-record schema.  Lean proves the gate-audit proposition
over the typed audit record by exhibiting, for each gate, the check
predicate, the passing status, the evidence record, and the implication
from the passed check to the corresponding no-smuggling or
no-overreading property.  Each gate rules out a specific computational
shortcut, and the lawfulness records make those rule-outs explicit and
audit-friendly.
 

\section{Formalization (Lean coverage)}
\label{app:formalization}

This appendix records how the paper-side definitions, propositions,
theorems, recognition source, obligation, and nonclaims are tracked in
the Lean development.  It is a per-row map from the statements of
record in the paper to their Lean declarations or formalization-harness
status, together with the disclosure wording used in the body.  It is
not a restatement of the body theorems, not a substitute for the
mathematical proof sketches, and not an external-bibliography list.

Lean module paths and identifier strings appear in this appendix.  The
body prose, including per-theorem disclosure footnotes, uses paper-prose
names such as ``saturated current observable family,'' ``Translation
Theorem B,'' or ``recognition source,'' and points to this appendix for
the corresponding formalization names.  The present
paper has 49 statements-of-record rows: 41 mechanized rows, one
recognition-source carrier, one BirdInt-judgment obligation, and six
nonclaim rows.  Of the 41 mechanized rows, 15 are new with the Class-B
Irremovability barrier of
\Cref{sec:class-b-barrier}: one orbit-collapse theorem, four Class-B
definitions, five closure lemmas, the Basis Combination Closure
theorem, the Basis Exhaustiveness theorem, and the two Irremovability
theorems; seven are the AOR-instance rows of
\Cref{sec:pvnp-aor-instance}.  The coverage table below lists all 49
rows.

\paragraph{Lean-disclosure wording.}
The body and appendix use the following wording discipline.
\begin{center}
\small
\begin{tabularx}{\textwidth}{@{}>{\raggedright\arraybackslash}p{0.22\textwidth}>{\raggedright\arraybackslash}p{0.30\textwidth}>{\raggedright\arraybackslash}X@{}}
\toprule
Coverage / alignment & Permitted wording & Meaning in this paper \\
\midrule
theorem / faithful &
``Lean proves'' or ``verified in Lean'' &
The Lean theorem matches the paper statement.  This applies to
Translation A, Translation B, the conditional four-step contradiction,
the SAT-in-\(NP\) theorem, the gate audit, orbit collapse, the Class-B
closure lemmas, Basis Combination Closure, and both Class-B
Irremovability theorems. \\

theorem / narrowed surrogate &
``Lean proves the stated scope; the broader claim is paper-side with
narrowing disclosed here'' &
The structural shape is mechanized, but the paper discloses the
narrowing.  In this paper the closure theorem for \(\EQbang\) is
realized through typed operation records for post-processing, finite
products, and composition. \\

theorem / projection-packaged &
``tracked by the formalization harness as a conditional schema over
explicit hypotheses'' &
The Lean theorem repackages a typed input record that already supplies
the equivalence rather than deriving it.  In this paper Basis
Exhaustiveness is parameterized by an intrinsic-feature decomposition
record whose fields encode the bi-implication between intrinsic
features and basis-generated representatives. \\

definition / faithful &
``realized in Lean as \texttt{Module.name}'' &
The Lean declaration is the paper definition's typed realization.  This
is the default for the saturated-layer, observable, quotient, lift,
currentization, appendix-definition, and Class-B definition rows. \\

definition / narrowed surrogate &
``realized in Lean with the parameter or scope exposed'' &
The Lean definition exposes a parameter or finite-stage scope that the
body treats in its intended mathematical form.  In this paper the
history carrier is realized with the finite-stage form exposed as a
free argument. \\

recognition source &
``encoded as a structural recognition-source carrier'' or ``supplied by
\citet{TsiokosHiddenness2026} and tracked as a typed structure'' &
The source \(\CSLSATHidden\) is not proved by the formalization.  It is
tracked by a typed carrier whose content is the negation of
\(\CurrentizeWit(\OmegaSAT)\) on the formed SAT layer. \\

obligation &
``recorded in the formalization harness as a manifest obligation'' &
The framework-relative \(\BirdInt\) judgment is bookkeeping over the
conditional theorem.  It is not a Lean theorem. \\

not mechanized &
No Lean citation &
The row is a scope-fence nonclaim.  The body uses the nonclaim wording
from \Cref{sec:scope} and does not cite Lean for it. \\
\bottomrule
\end{tabularx}
\end{center}

All Lean declarations \citep{DeMouraUllrich2021Lean4} listed in the
coverage table below live under the \nolinkurl{SixBirdsHiddenness.PvNP.*}
namespace.

\begin{landscape}
\footnotesize
\setlength{\tabcolsep}{1pt}
\renewcommand{\arraystretch}{0.82}
\noindent\textbf{Coverage table.}\par\smallskip
\begin{longtable}{@{}>{\raggedright\arraybackslash}p{0.33\linewidth}>{\raggedright\arraybackslash}p{0.10\linewidth}>{\raggedright\arraybackslash}p{0.30\linewidth}>{\raggedright\arraybackslash}p{0.08\linewidth}>{\raggedright\arraybackslash}p{0.17\linewidth}@{}}
\toprule
Paper label & Paper-prose name & Lean decl & Coverage / alignment & Disclosure wording (body location) \\
\midrule
\endfirsthead
\toprule
Paper label & Paper-prose name & Lean decl & Coverage / alignment & Disclosure wording (body location) \\
\midrule
\endhead
\midrule
\multicolumn{5}{r}{\footnotesize\itshape continued on next page} \\
\endfoot
\bottomrule
\endlastfoot
\path|def:pvnp:audited-shell| & Audited SAT shell & \mbox{\scriptsize\texttt{PvNP.SaturatedClosure.auditedShell}} & definition / faithful & Definition entry; realized in Lean as listed here (\Cref{sec:saturated-closure}). \\

\path|def:pvnp:six-mechanisms| & Six mechanisms (\(P_{1}\)--\(P_{6}\)) & \mbox{\scriptsize\texttt{PvNP.SaturatedClosure.sixMechanisms}} & definition / faithful & Appendix definition; realized in Lean as listed here (\Cref{app:primitive-roles}). \\

\path|def:pvnp:history-carrier| & SAT history carrier & \mbox{\scriptsize\texttt{PvNP.SaturatedClosure.historyCarrier}} & definition / narrowed surrogate & Realized with finite-stage form exposed (\Cref{sec:saturated-closure}). \\

\path|def:pvnp:saturated-closure| & Saturated SAT layer & \mbox{\scriptsize\texttt{PvNP.SaturatedClosure.saturatedClosure}} & definition / narrowed surrogate & Lean record carries the typed core (\(\Hsat,\EQbang,\EMbang,\ComparisonMap,\operatorname{Audit}^{!}\)) plus the closure operator and admissibility/audit fields; the audited shell \(\Saud\) and the saturated quotients \(\Qsat,\Msat\) live in companion records and are recorded paper-side as the aggregate object (\Cref{sec:saturated-closure}). \\

\path|def:pvnp:i-poly| & Lawful current instrument & \mbox{\scriptsize\texttt{PvNP.ObservableSaturation.iPoly}} & definition / faithful & Definition entry; realized in Lean as listed here (\Cref{sec:observable-saturation}). \\

\path|def:pvnp:e-q-saturation| & Saturated current-observable family & \mbox{\scriptsize\texttt{PvNP.ObservableSaturation.eQSaturation}} & definition / faithful & Definition entry; realized in Lean as listed here (\Cref{sec:observable-saturation}). \\

\path|def:pvnp:lawfulness-discipline| & Lawfulness discipline & \mbox{\scriptsize\texttt{PvNP.ObservableSaturation.lawfulnessDiscipline}} & definition / faithful & Definition entry; realized in Lean as listed here (\Cref{sec:observable-saturation}). \\

\path|thm:pvnp:e-q-closure-properties| & Closure properties of \(\EQbang\) & \mbox{\scriptsize\texttt{PvNP.ObservableSaturation.eQClosureProperties}} & theorem / narrowed surrogate & Lean proves the admission-level closure (the wrapped transducer preserves certificate, conventions, replay/audit, no-oracle commitments, and admitted-because-lawful status); the semantic computation clauses are supplied by the operation-record contracts rather than derived from a primitive eval clause (\Cref{sec:observable-saturation}). \\

\path|def:pvnp:q-saturated| & Saturated current quotient & \mbox{\scriptsize\texttt{PvNP.Quotients.qSaturated}} & definition / faithful & Definition entry; realized in Lean as listed here (\Cref{sec:quotients}). \\

\path|def:pvnp:m-saturated| & Saturated predictive quotient & \mbox{\scriptsize\texttt{PvNP.Quotients.mSaturated}} & definition / faithful & Definition entry; realized in Lean as listed here (\Cref{sec:quotients}). \\

\path|def:pvnp:comparison-map| & Comparison map & \mbox{\scriptsize\texttt{PvNP.Quotients.comparisonMap}} & definition / faithful & Definition entry; realized in Lean as listed here (\Cref{sec:quotients}). \\

\path|def:pvnp:omega-sat| & Autonomous lift & \mbox{\scriptsize\texttt{PvNP.AutonomousLift.omegaSat}} & definition / narrowed surrogate & Record exposes \texttt{generateHistory}, \texttt{sigma\_L}, \texttt{Cont\_M}, an \texttt{updateRule}, and no-oracle Prop fields as independent components; the deterministic-policy linkage between the update rule and history generation is paper-side (\Cref{sec:autonomous-lift}). \\

\path|def:pvnp:vis-bang| & Instrument-relative visibility map & \mbox{\scriptsize\texttt{PvNP.AutonomousLift.visBang}} & definition / faithful & Definition entry; realized in Lean as listed here (\Cref{sec:autonomous-lift}). \\

\path|def:pvnp:currentize-dec| & Decision currentization & \mbox{\scriptsize\texttt{PvNP.Currentize.currentizeDec}} & definition / faithful & Definition entry; realized in Lean as listed here (\Cref{sec:currentize}). \\

\path|def:pvnp:currentize-wit| & Witness currentization & \mbox{\scriptsize\texttt{PvNP.Currentize.currentizeWit}} & definition / faithful & Definition entry; realized in Lean as listed here (\Cref{sec:currentize}). \\

\path|thm:translation-A| & Translation Theorem A & \mbox{\scriptsize\texttt{PvNP.TranslationA.translationA}} & theorem / narrowed surrogate & Lean proves the iff packaging of Translation Theorem A from the standard \(\mathrm{SAT}\in P\) bridge taken as a parametric hypothesis; the standard membership-to-observable bridge is not mechanized internally (\Cref{sec:translation-a}). \\

\path|thm:translation-B| & Translation Theorem B & \mbox{\scriptsize\texttt{PvNP.TranslationB.translationB}} & theorem / narrowed surrogate & Lean proves the iff packaging of Translation Theorem B from the self-reduction wrapper, canonical-witness correctness, closure admission, post-processor admission, and SAT-bit correctness taken as parametric hypotheses; the wrapper construction itself is not built inside Lean (\Cref{sec:translation-b}). \\

\path|axiom:pvnp:csl-sat-hidden| & Recognition source \(\CSLSATHidden\) & \mbox{\scriptsize\texttt{PvNP.RecognitionSource.recognitionSource}} & recognition source / not applicable & Encoded as typed structure carrier; supplied by \citet{TsiokosHiddenness2026} (\Cref{sec:recognition-source,app:recognition-source-carrier}). \\

\path|thm:four-step-contradiction-conditional| & Four-step contradiction (conditional) & \begin{tabular}[t]{@{}l@{}}{\scriptsize\texttt{PvNP.ContradictionChain.}}\\{\scriptsize\texttt{fourStepContradictionConditional}}\end{tabular} & theorem / faithful & Lean proves the conditional four-step contradiction (\Cref{sec:contradiction-chain}). \\

\path|thm:pvnp:sat-in-np| & SAT in \(NP\) via standard verifier & \mbox{\scriptsize\texttt{PvNP.ContradictionChain.satInNP}} & theorem / narrowed surrogate & Lean transports the verifier-package conclusion from supplied textbook hypotheses (polynomial-time verifier, witness-size bound, standard verifier-package combinator); no concrete SAT encoding or verifier is constructed inside Lean (\Cref{sec:contradiction-chain}). \\

\path|oblig:pvnp:final-judgment| & BirdInt judgment & --- & obligation / not applicable & Recorded as manifest obligation (\Cref{sec:judgment}). \\

\path|prop:pvnp:gate-audit| & Gate audit (\(\GateOne\)--\(\GateSeven\)) & \mbox{\scriptsize\texttt{PvNP.GateAudit.gateAudit}} & theorem / faithful & Lean proves the gate-audit proposition over the typed audit record (\Cref{app:gate-audit}). \\

\path|nonclaim:pvnp:standard-zfc-pneqnp| & Standard-ZFC unconditional \(P\neq NP\) (nonclaim) & --- & not mechanized & Scope-fence wording only (\Cref{sec:scope}). \\

\path|nonclaim:pvnp:circumvent-classical-barriers| & Classical-barrier circumvention (nonclaim) & --- & not mechanized & Scope-fence wording only (\Cref{sec:scope}). \\

\path|nonclaim:pvnp:hiddenness-derivable| & Hiddenness derivable from primitives (nonclaim) & --- & not mechanized & Scope-fence wording only (\Cref{sec:scope}). \\

\path|thm:pvnp:orbit-collapse| & Fixed-Presentation Orbit Collapse & \mbox{\scriptsize\texttt{PvNP.ClassB.OrbitCollapse.orbitCollapse}} & theorem / faithful & Lean proves the common-kernel collapse (\Cref{subsec:class-b:meta-results}). \\

\path|def:pvnp:class-b-law| & Class-B carrier-property law & \mbox{\scriptsize\texttt{PvNP.ClassB.Definitions.classBLaw}} & definition / faithful & Definition entry; realized in Lean as listed here (\Cref{app:class-b-definitions}). \\

\path|def:pvnp:basis-fix| & Fixed-presentation feature basis & \mbox{\scriptsize\texttt{PvNP.ClassB.Definitions.basisFix}} & definition / faithful & Definition entry; realized in Lean as listed here (\Cref{app:class-b-definitions}). \\

\path|def:pvnp:normal-classes| & Normal feature classes & \mbox{\scriptsize\texttt{PvNP.ClassB.Definitions.normalClasses}} & definition / faithful & Definition entry; realized in Lean as listed here (\Cref{app:class-b-definitions}). \\

\path|def:pvnp:basis-generated| & Basis-generated feature expressions & \mbox{\scriptsize\texttt{PvNP.ClassB.Definitions.basisGenerated}} & definition / faithful & Definition entry; realized in Lean as listed here (\Cref{app:class-b-definitions}). \\

\path|lem:pvnp:quotient-meet-join| & Quotient meet/join closure & \mbox{\scriptsize\texttt{PvNP.ClassB.Lemmas.quotientMeetJoin}} & theorem / faithful & Lean proves the closure lemma directly (\Cref{app:class-b-closure-lemmas}). \\

\path|lem:pvnp:orbit-quotient| & Orbit quotient closure & \mbox{\scriptsize\texttt{PvNP.ClassB.Lemmas.orbitQuotient}} & theorem / faithful & Lean proves the closure lemma directly (\Cref{app:class-b-closure-lemmas}). \\

\path|lem:pvnp:holonomy-quotient| & Holonomy quotient closure & \mbox{\scriptsize\texttt{PvNP.ClassB.Lemmas.holonomyQuotient}} & theorem / faithful & Lean proves the closure lemma directly (\Cref{app:class-b-closure-lemmas}). \\

\path|lem:pvnp:presentation-refinement| & Presentation refinement closure & \mbox{\scriptsize\texttt{PvNP.ClassB.Lemmas.presentationRefinement}} & theorem / faithful & Lean proves the closure lemma directly (\Cref{app:class-b-closure-lemmas}). \\

\path|lem:pvnp:time-stratum| & Time-stratum refinement & \mbox{\scriptsize\texttt{PvNP.ClassB.Lemmas.timeStratum}} & theorem / faithful & Lean proves the closure lemma directly (\Cref{app:class-b-closure-lemmas}). \\

\path|thm:pvnp:basis-combination-closure| & Basis Combination Closure & \begin{tabular}[t]{@{}l@{}}{\scriptsize\texttt{PvNP.ClassB.CombinationClosure.}}\\{\scriptsize\texttt{basisCombinationClosure}}\end{tabular} & theorem / faithful & Lean proves a per-constructor case analysis over the feature-expression type, with the normal-class assignment defined by outer constructor pattern; the polynomial pairing and post-processing facts are supplied through the narrowed admission-level \texttt{eQClosureProperties} record clauses (\Cref{subsec:class-b:combination-closure}). \\

\path|thm:pvnp:basis-exhaustiveness| & Basis Exhaustiveness & \begin{tabular}[t]{@{}l@{}}{\scriptsize\texttt{PvNP.ClassB.BasisExhaustiveness.}}\\{\scriptsize\texttt{basisExhaustiveness}}\end{tabular} & theorem / projection-packaged & Tracked by the formalization harness as a conditional schema over a typed decomposition record; the bi-implication is supplied by the record fields rather than derived (\Cref{subsec:class-b:basis-exhaustiveness}). \\

\path|thm:pvnp:irremovability-core| & Class-B Irremovability Core & \mbox{\scriptsize\texttt{PvNP.ClassB.Irremovability.irremovabilityCore}} & theorem / faithful & Lean proves the quotient/joint-access case analysis (\Cref{subsec:class-b:irremovability}). \\

\path|thm:pvnp:irremovability-scoped| & Six Birds Class-B Irremovability & \begin{tabular}[t]{@{}l@{}}{\scriptsize\texttt{PvNP.ClassB.Irremovability.}}\\{\scriptsize\texttt{irremovabilityScoped}}\end{tabular} & theorem / narrowed surrogate & Lean proves the scoped five-way case analysis at the feature-classification layer; the theorem ranges over a \texttt{classBLawFeature} value (basis-generated or non-basis), and the connection from a \texttt{classBLaw} record to a \texttt{classBLawFeature} value is paper-side (\Cref{subsec:class-b:irremovability}). \\

\path|nonclaim:pvnp:absolute-irremovability| & Absolute irremovability (nonclaim) & --- & not mechanized & Scope-fence wording only (\Cref{sec:scope}). \\

\path|def:pvnp:\allowbreak aor-instance:\allowbreak object| & AOR instance object attached to \(\Tsat\) & \begin{tabular}[t]{@{}l@{}}{\scriptsize\texttt{PvNP.AORInstance.}}\\{\scriptsize\texttt{aorInstanceObject}}\end{tabular} & definition / narrowed surrogate & Carrier-attached bookkeeping record; fields bind to the concrete PvNP types (\(\Saud\), \(\Hsat\), \(\Ipoly\), \(\EQbang\), quotients, \(\OmegaSAT\), currentization, recognition source, Class-B basis). Narrowed: AOR field order is bookkeeping, not new mathematical content (\Cref{sec:pvnp-aor-instance}). \\

\path|thm:pvnp:\allowbreak aor:\allowbreak mechanical-records| & Mechanical AOR records for \(\Tsat\) & \begin{tabular}[t]{@{}l@{}}{\scriptsize\path|PvNP.AORInstance.|}\\{\scriptsize\texttt{aorMechanicalRecords}}\end{tabular} & theorem / narrowed surrogate & By-construction / approved-other discharges for the typed-shell definitions, \(\EQbang\)-closure admission, and SAT verifier package (\Cref{sec:pvnp-aor-instance}). \\

\path|thm:pvnp:\allowbreak aor:\allowbreak recognition-source-discharge| & Recognition-source discharge & \begin{tabular}[t]{@{}l@{}}{\scriptsize\path|PvNP.AORInstance.|}\\{\scriptsize\texttt{aorRecognitionSourceDischarge}}\end{tabular} & theorem / faithful & \(\Delta_{\mathrm{source}}\) primary, AOR-8 status \texttt{hidden}; Lean rewrites through \(\Gamma\).\texttt{not\_currentize\_wit} to derive \(\neg\,\CurrentizeWit(\OmegaSAT)\); consumes \path|RecognitionSource.recognitionSource| substantively (axiom-free) (\Cref{sec:pvnp-aor-instance}). \\

\path|thm:pvnp:\allowbreak aor:\allowbreak translation-discharge| & Translation-interface discharge & \begin{tabular}[t]{@{}l@{}}{\scriptsize\path|PvNP.AORInstance.|}\\{\scriptsize\texttt{aorTranslationDischarge}}\end{tabular} & theorem / narrowed surrogate & \(\Delta_{\mathrm{transport}}\) primary, AOR-8 status \texttt{bridged}; Lean invokes \path|TranslationA.translationA| and \path|TranslationB.translationB| to supply the two iff bridges. Inherits the upstream narrowed-surrogate iff-packaging of both translations (\Cref{sec:pvnp-aor-instance}). \\

\path|thm:pvnp:\allowbreak aor:\allowbreak class-b-scope-fence| & Class-B scope-fence discharge & \begin{tabular}[t]{@{}l@{}}{\scriptsize\path|PvNP.AORInstance.|}\\{\scriptsize\texttt{aorClassBScopeFence}}\end{tabular} & theorem / narrowed surrogate & \(\Delta_{\mathrm{horizon}}\) primary; status \texttt{blocked} inside the method class, \texttt{out\_of\_scope/nonclaim} outside; Lean invokes \path|Irremovability.irremovabilityScoped|. Inherits the upstream feature-classification-layer narrowing (\Cref{sec:pvnp-aor-instance}). \\

\path|thm:pvnp:\allowbreak aor:\allowbreak judgment-package| & BirdInt judgment-package discharge & \begin{tabular}[t]{@{}l@{}}{\scriptsize\path|PvNP.AORInstance.|}\\{\scriptsize\texttt{aorJudgmentPackage}}\end{tabular} & theorem / narrowed surrogate & \(\Delta_{\mathrm{role}}\) primary, AOR-8 status \texttt{zero}; Lean invokes \path|ContradictionChain.fourStepContradictionConditional| and \path|ContradictionChain.satInNP|. Inherits the upstream verifier-package narrowing of \texttt{satInNP} (\Cref{sec:pvnp-aor-instance}). \\

\path|thm:pvnp:\allowbreak aor:\allowbreak in-aor-casc| & PvNP AOR-instance theorem & \begin{tabular}[t]{@{}l@{}}{\scriptsize\path|PvNP.AORInstance.|}\\{\scriptsize\texttt{aorInAorCasc}}\end{tabular} & theorem / projection-packaged & Headline: \(\Tsat\) tagged as \(\mathbf{AOR}^{\mathrm{casc}}\). Conclusion exposes the five discharge statuses, the transported \(\neg\,C.\texttt{SATInPProp}\) and \(C.\texttt{SATInNPProp}\), and the carrier tag. Projection-packaged as detailed in \Cref{subsec:pvnp-aor-narrowing}; the cascade-stable AOR tag is a local carrier marker, not membership in an imported \(\mathbf{AOR}^{\mathrm{casc}}\) predicate. \\

\path|nonclaim:pvnp:\allowbreak aor-instance-\allowbreak formed-layer-only| & AOR instance restricted to formed SAT layer (nonclaim) & --- & not mechanized & Scope-fence wording only (\Cref{sec:scope}). \\

\path|nonclaim:pvnp:\allowbreak aor-instance-\allowbreak no-unconditional| & No unconditional \(P\neq NP\) via AOR instance (nonclaim) & --- & not mechanized & Scope-fence wording only (\Cref{sec:scope}). \\
\end{longtable}
\end{landscape}

The six nonclaim rows carry no Lean coverage.  They are included in
the table because they are statements of record for the paper's scope,
but the body does not cite Lean for them.  Their content is expressed
through the scope-fence language in \Cref{sec:scope}: no standard-ZFC
unconditional \(P\neq NP\) separation, no native circumvention claim
for the classical barriers, no derivation of the hiddenness source from
framework primitives, no absolute irremovability claim for
\(\CSLSATHidden\), no AOR-instance enlargement beyond the formed SAT
layer, and no unconditional separation via the AOR-instance reading.

The narrowed-surrogate rows are the main mechanization disclosures.
For \texttt{def:pvnp:\allowbreak history-carrier}, the body uses the
intended directed-colimit reading while the Lean realization exposes
the finite-stage carrier as a free argument.  For
\texttt{thm:pvnp:\allowbreak e-q-\allowbreak closure-properties}, Lean proves
the admission-level closure --- the wrapped transducer preserves the
polynomial-time certificate, the input/output conventions, the replay
and audit bookkeeping, the no-oracle commitments, and the
admitted-because-lawful status --- by deriving these clauses from a
typed operation-record schema (post-processing, finite-product, and
composition records); the semantic computation clauses (post-processed
wrapper computes \(g\circ f\), finite-product wrapper computes the
tuple of outputs, composition wrapper computes the composite) are
supplied as typed-record commitments rather than derived from a
primitive \texttt{eval} clause.  For \texttt{def:pvnp:\allowbreak omega-sat}, the Lean
record stores the policy components (\texttt{generateHistory},
\texttt{sigma\_L}, \texttt{Cont\_M}, \texttt{updateRule}) and the
no-oracle Prop fields as independent components; the deterministic
linkage between the update rule and the history-generation map is
paper-side.  For \texttt{thm:translation-A} and
\texttt{thm:translation-B}, Lean proves the iff packaging from
parametric inputs that name the standard membership-to-observable
bridge (Translation A) or the self-reduction wrapper and its
correctness contracts (Translation B); the standard constructions
themselves are not built inside Lean.

The Class-B rows are theorem / faithful, definition / faithful,
theorem / projection-packaged (the Basis Exhaustiveness row, whose
Lean theorem is parameterized by a typed decomposition record rather
than derived from carrier primitives), or theorem / narrowed surrogate
(the scoped Irremovability row, whose Lean theorem ranges over a
\texttt{classBLawFeature} value rather than directly over the
\texttt{classBLaw} record).  The inductive feature-expression type and
the inductive outcome predicate carry the per-class consequences used
by the paper.  The Lean proof of Basis Combination Closure performs
per-constructor case analysis over the feature-expression type and
invokes the closure lemmas at the corresponding cases; five of the
seven closure lemmas are mechanized directly in
\texttt{ClassB.Lemmas} and the two polynomial closure facts (pairing
and post-processing) are supplied through the narrowed admission-level
clauses of \texttt{eQClosureProperties}.  The Basis
Exhaustiveness theorem is tracked by the formalization harness as a
conditional schema over that decomposition record: the substantive
content is the paper-side slot-by-slot argument.  The Irremovability
theorems are parameterized by the closure and orbit-collapse
hypotheses; the scoped theorem operates on the feature-classification
layer, and the connection from a typed \texttt{classBLaw} record to
its feature classification is recorded paper-side (\Cref{app:class-b-definitions}).
A later mechanization extension may instantiate the abstract carrier
presentation directly against the concrete PvNP carrier and thread
the record-to-feature connection through Lean.

The recognition-source row uses the carrier discipline disclosed in
\Cref{app:recognition-source-carrier}: the Lean encoding is a typed
structure, not a primitive proof-assistant assertion.  The BirdInt
judgment row is a manifest obligation, not a Lean theorem, and
\Cref{sec:judgment} gives the framework-relative judgment shape.  The
Lean development is mathlib-free
\citep{MathlibCommunity2020LeanMathematicalLibrary}, and the theorem
rows listed above pass the project's axiom audit under the trust base
\(\{\)\texttt{Classical.choice}, \texttt{propext},
\texttt{Quot.sound}, \texttt{choice}, \texttt{funext}\(\}\)
\citep{MartinLof1984IntuitionisticTypeTheory}.

\subsection*{Projection-packaging narrowing for the AOR-instance theorem}
\label{subsec:pvnp-aor-narrowing}

The seven AOR-instance rows
(\path|def:pvnp:aor-instance:object|,
 \path|thm:pvnp:aor:mechanical-records|,
 \path|thm:pvnp:aor:recognition-source-discharge|,
 \path|thm:pvnp:aor:translation-discharge|,
 \path|thm:pvnp:aor:class-b-scope-fence|,
 \path|thm:pvnp:aor:judgment-package|,
 \path|thm:pvnp:aor:in-aor-casc|)
mechanize the recasting of \Cref{sec:pvnp-aor-instance} under the
audit-closure meta-theory of \citet{Tsiokos2026AOR}.  Three points
deserve disclosure.

First, the AOR-instance object \texttt{aorInstanceObject} is
carrier-attached: its fields are typed to the concrete PvNP records
(\path|SaturatedClosure.auditedShell|, \texttt{historyCarrier},
\texttt{saturatedClosure}, \path|ObservableSaturation.iPoly|,
\texttt{eQSaturation}, \path|Quotients.qSaturated|,
\texttt{mSaturated}, \texttt{comparisonMap},
\path|AutonomousLift.omegaSat|, \texttt{visBang},
\path|Currentize.currentizeDec|, \texttt{currentizeWit},
\path|RecognitionSource.recognitionSource|, and Class-B
\texttt{basisFix}).  The four substantive discharge theorems
(\path|thm:pvnp:aor:recognition-source-discharge|,
 \path|thm:pvnp:aor:translation-discharge|,
 \path|thm:pvnp:aor:class-b-scope-fence|,
 \path|thm:pvnp:aor:judgment-package|) consume the corresponding
existing PvNP theorem constants
(\path|TranslationA.translationA|,
\path|TranslationB.translationB|,
\path|Irremovability.irremovabilityScoped|,
\path|ContradictionChain.fourStepContradictionConditional|,
\path|ContradictionChain.satInNP|, and the
\texttt{recognitionSource} carrier) substantively, not as decorative
imports.

Second, three of the four substantive discharge rows carry a
\texttt{narrowed surrogate} alignment because each inherits the
upstream narrowing of the PvNP theorem it consumes: the
translation discharge from Translation A and Translation B (iff
packaging from supplied contracts); the Class-B scope-fence
discharge from the scoped irremovability theorem (which ranges
over \texttt{classBLawFeature} rather than the \texttt{classBLaw}
record); the BirdInt judgment-package discharge from the standard
verifier package (transfer from supplied textbook hypotheses).  The
recognition-source discharge is \texttt{faithful}: it consumes the
\texttt{recognitionSource} carrier substantively and rewrites through
\(\Gamma\).\texttt{not\_currentize\_wit} to derive
\(\neg\,\CurrentizeWit(\OmegaSAT)\) without inheriting an upstream
narrowing.  The discharge theorems are faithful to the AOR meta-theory
and transmit any upstream narrowings to the AOR-instance layer.

Third, the headline \path|thm:pvnp:aor:in-aor-casc| carries the
distinct alignment \texttt{projection-packaged}: the cascade-stable
AOR tag (\(\Tsat\) tagged as \(\mathbf{AOR}^{\mathrm{casc}}\)) is a
local marker carried by the \texttt{aorInstanceObject} field
\texttt{targetClass}, not membership in an imported
\(\mathbf{AOR}^{\mathrm{casc}}\) predicate.  The AOR meta-theory is
consumed paper-side; this subsection records the projection-packaging
discipline.  The conclusion of
\path|thm:pvnp:aor:in-aor-casc| exposes the five discharge statuses,
the transported \(\neg\,C.\texttt{SATInPProp}\) and
\(C.\texttt{SATInNPProp}\), and the carrier tag itself.

The \texttt{aorCascadePackage} structure also carries the
\texttt{judgment} slot tied to the carrier's own
\texttt{SATInPProp}/\texttt{SATInNPProp} fields, so the cascade
cannot splice in an unrelated BirdInt judgment proposition.  The
seven AOR-instance entries remain within the project trust base
\(\{\)\texttt{propext}, \path|Quot.sound|\(\}\); two of the six discharge theorems
(\texttt{aorRecognitionSourceDischarge} and
\texttt{aorJudgmentPackage}) are axiom-free.  A future strengthening
would promote the proposition-valued discharge fields to
data-bearing AOR records and lift the projection-packaging.
 

\section{Class-B Irremovability apparatus}
\label{app:class-b-apparatus}

This appendix records the formal Class-B definitions and closure
lemmas used in \Cref{sec:class-b-barrier}.  The body states the
Class-B Irremovability barrier as the paper's third result; the
technical apparatus supporting that theorem is collected here so that
the precise definitions do not interrupt the body argument.  The
present subsection gives the definition of a Class-B law, the fixed
feature basis, the five normal classes, basis-generated feature
expressions, and intrinsic features of the fixed-presentation carrier.
The closure lemmas are recorded in \Cref{app:class-b-closure-lemmas}.

\subsection{Class-B definitions}
\label{app:class-b-definitions}

\begin{definition}[Class-B carrier-property law]
\label{def:pvnp:class-b-law}
A \emph{Class-B carrier-property law} is a theorem of the form
\[
  \Phi(\mathcal C) \Longrightarrow
  \neg \exists O\in\mathcal O_{\mathcal C}\,:\,T(O),
\]
where \(\mathcal C\) is a Six Birds carrier, \(\Phi(\mathcal C)\) is
an intrinsic feature of that carrier, and the impossibility conclusion
follows from the feature itself.  The feature is not an external
recognition source, is not a restatement of the target, and is not a
claim-discipline restriction on how the target is allowed to be
asserted.
\end{definition}

Class-B laws are contrasted with Class-A claim-discipline laws.
Class-A laws, such as no-overread, adequacy, and quotientality, act on
the structure of the claim being made.  Class-B laws act on the
carrier: they try to derive an impossibility from a property of the
computational substrate.  The formalization records the intrinsic
feature condition, the non-recognition-source condition, the
non-claim-discipline condition, and the carrier-feature impossibility
conclusion as explicit fields of the corresponding typed record.

\begin{definition}[Fixed-presentation feature basis]
\label{def:pvnp:basis-fix}
For the standard fixed-presentation polynomial current carrier, the
\emph{fixed-presentation feature basis} is
\[
  \mathcal B_{\mathrm{fix}}
  =
  \{\Delta_X,\ \mathsf{PolyClosure}(\EQbang),\
    \mathsf{TimeFilt},\ \mathsf{PresData}\}.
\]
Here \(\Delta_X\) is the diagonal equality relation on the encoded
input carrier; \(\mathsf{PolyClosure}(\EQbang)\) is the closure
structure of the saturated current-observable family under pairing,
finite products, polynomial-time post-processing, and polynomial-time
composition; \(\mathsf{TimeFilt}\) is the polynomial resource ladder
\[
  E^{[1]}\subseteq E^{[2]}\subseteq \cdots,
  \qquad
  \EQbang=\bigcup_k E^{[k]};
\]
and \(\mathsf{PresData}\) is the fixed-presentation structural data:
bit positions, variable names, symbol order, clause order, delimiter
conventions, and the analogous features of the chosen SAT encoding.
\end{definition}

The four primitives correspond to the four source slots of the carrier
package \((X,\EQbang,\Ipoly,\mathsf{Pres})\) used in
\Cref{thm:pvnp:basis-exhaustiveness}: the carrier set, the observable
family, the polynomial-time interface, and the fixed presentation.
The substantive closure facts for \(\mathsf{PolyClosure}(\EQbang)\)
are those of \Cref{thm:pvnp:e-q-closure-properties}.

\begin{definition}[Normal feature classes]
\label{def:pvnp:normal-classes}
The five \emph{normal feature classes} are \(Q,J,C,P,T\).
\begin{enumerate}
  \item \(Q\), the \emph{quotient class}, consists of equivalence
  relations \(R\subseteq X\times X\) such that the proposed Class-B
  law requires \(R\subseteq\ker(F)\) for every \(F\in\EQbang\).
  Equivalently, every lawful current observable is required to be
  constant on \(R\)-classes.

  \item \(J\), the \emph{joint-access class}, consists of finite
  tuples \(F_1,\ldots,F_n\in\EQbang\) of separately current
  observables whose joint access \((F_1,\ldots,F_n)\) is alleged not
  to lie in \(\EQbang\).

  \item \(C\), the \emph{positive-closure class}, consists of
  constructions that yield new current observables from current ones
  by composition, pairing, products, or polynomial-time
  post-processing.

  \item \(P\), the \emph{presentation-refinement class}, consists of
  refinements of the carrier by presentation data.  These refinements
  supply breaker data rather than invariance constraints.

  \item \(T\), the \emph{time-filtration class}, consists of features
  indexed by polynomial resource strata \(E^{[k]}\), either at a fixed
  \(k\) or uniformly across all \(k\).
\end{enumerate}
\end{definition}

These five classes are exhaustive for basis-generated feature
expressions by \Cref{thm:pvnp:basis-combination-closure}.

\begin{definition}[Basis-generated feature expression]
\label{def:pvnp:basis-generated}
A \emph{basis-generated feature expression} is built from
\(\mathcal B_{\mathrm{fix}}\) by finitely many applications of the
following closure operations: quotienting, orbiting, transporting,
holonomy or gluing-class formation, pairing, finite products,
polynomial-time post-processing, resource filtration, and presentation
refinement.  This is the standard closure-operator construction on a
basis of generators \citep{DaveyPriestley2002LatticesOrder}, restricted
to the operations admitted by the carrier package.  Transport has three
flavors: route-equivalence transport, record-carrying transport, and
resource-indexed transport.

The normal-class assignment is a total function
\[
  \mathsf{normalClass}:\{\text{basis-generated expressions}\}
  \longrightarrow \{Q,J,C,P,T\}.
\]
The diagonal primitive, quotienting, orbiting,
transport-by-route, and holonomy or gluing-class formation land in
\(Q\).  The polynomial-closure primitive, pairing, finite products,
and polynomial-time post-processing land in \(C\).  The time-filtration
primitive, transport-by-resource, and resource filtration land in
\(T\).  The presentation primitive, transport-by-record, and
presentation refinement land in \(P\).  Joint-access expressions land
in \(J\) precisely when they assert obstruction to finite joint access
of separately current observables.
\end{definition}

\begin{definition}[Intrinsic feature of the fixed-presentation carrier]
An \emph{intrinsic feature} of
\((X,\EQbang,\Ipoly,\mathsf{Pres})\) is a feature definable from that
carrier package alone.  It does not import an external predicate, a
measure, a topology, an oracle, an advice string, a probability law, a
metric not induced by the encoding, a group action not already present,
or algebraic structure not already present.  Equivalently, it is
invariant under isomorphisms of the carrier package.
\end{definition}

Every intrinsic feature draws its content from one of the four carrier
slots \(X\), \(\EQbang\), \(\Ipoly\), or \(\mathsf{Pres}\).  Features
drawn from \(X\) reduce to \(\Delta_X\); features drawn from
\(\EQbang\) reduce to \(\mathsf{PolyClosure}(\EQbang)\) or to common
kernels collapsed by \Cref{thm:pvnp:orbit-collapse}; features drawn
from \(\Ipoly\) reduce to \(\mathsf{TimeFilt}\); and features drawn
from \(\mathsf{Pres}\) reduce to \(\mathsf{PresData}\).  Combinations
of these four slots are basis-generated by
\Cref{thm:pvnp:basis-combination-closure}.

There are several natural candidate ``fifth features.''  In the
standard fixed-presentation carrier they do not produce additional
Class-B basis primitives.  Categorical or 2-categorical structure
reduces to \(\mathsf{PolyClosure}(\EQbang)\), to
\(\mathsf{PresData}\), or else exits to a richer carrier.  Diagonal or
universal Turing structure reduces to the resource-ladder simulation
recorded by \(\mathsf{TimeFilt}\).  Topological or metric structure
reduces to the discrete diagonal \(\Delta_X\) or to an encoding-induced
presentation feature.  Probabilistic or measure structure exits the
standard deterministic carrier.  Other resources such as space,
randomness, advice, or nonuniformity exit the deterministic
polynomial-time interface.  Algebraic structure either reduces to
computable operations already accounted for by
\(\mathsf{PolyClosure}(\EQbang)\), or exits to an enriched carrier.
Order, length, and prefix structure are encoding-derived and therefore
belong to \(\mathsf{PresData}\).  Semantic SAT structure is either
presentation-readable, reducing to \(\mathsf{PresData}\) together with
\(\mathsf{PolyClosure}(\EQbang)\), or target-equivalent.

\subsection{Closure lemmas}
\label{app:class-b-closure-lemmas}

This subsection records the seven closure lemmas invoked by the
per-constructor proof of \Cref{thm:pvnp:basis-combination-closure}.
The lemmas formalize the structural facts about quotient meets and
joins, group-action orbits, holonomy and gluing-class invariance,
polynomial pairing, polynomial post-processing, presentation
refinement, and time-stratum refinement.  Five of the seven lemmas
(quotient meet/join, orbit, holonomy/gluing, presentation refinement,
time-stratum) are mechanized as axiom-free derivations under
\texttt{SixBirdsHiddenness.PvNP.ClassB.Lemmas}; the remaining two
(polynomial pairing, polynomial post-processing) are supplied through
the narrowed admission-level clauses of
\Cref{thm:pvnp:e-q-closure-properties}, as detailed in the closing
paragraph of this subsection.  The paper statements are included here
for review.

\begin{lemma}[Quotient products, meets, and joins]
\label{lem:pvnp:quotient-meet-join}
Let \(R_1\) and \(R_2\) be equivalence relations on \(X\), and suppose
that both are respected by every lawful current observable
\(F\in\EQbang\).  Define
\[
  R_{\times}=R_1\cap R_2,
  \qquad
  R_{\vee}=\operatorname{EqCl}(R_1\cup R_2).
\]
Then \(R_{\times}\) is an equivalence relation respected by every
\(F\in\EQbang\), and \(R_{\vee}\) is an equivalence relation respected
by every \(F\in\EQbang\).  Moreover, by fixed-presentation orbit
collapse,
\[
  R_{\times}=R_{\vee}=\Delta_X .
\]
\end{lemma}

\begin{proof}[Sketch.]
Intersections of equivalence relations are equivalence relations.
The equivalence closure of a union is an equivalence relation, and
respect by all \(F\in\EQbang\) is preserved along finite chains of
identifications.  Thus both \(R_{\times}\) and \(R_{\vee}\) are
contained in the common kernel of all lawful current observables.
\Cref{thm:pvnp:orbit-collapse} identifies that common kernel with
\(\Delta_X\); reflexivity gives the reverse diagonal inclusion.
\end{proof}

\begin{lemma}[Group-action orbits are quotient features]
\label{lem:pvnp:orbit-quotient}
Let a group or groupoid \(G\) act on \(X\).  Let \(R_G\) be the orbit
relation
\[
  x\,R_G\,y
  \quad\Longleftrightarrow\quad
  \exists g\in G,\ y=g\cdot x .
\]
Then \(R_G\) is an equivalence relation.  If a proposed Class-B law
requires every \(F\in\EQbang\) to be invariant on \(G\)-orbits, then
\[
  R_G=\Delta_X .
\]
\end{lemma}

\begin{proof}[Sketch.]
Reflexivity is given by identity arrows, symmetry by inverses, and
transitivity by composition.  Orbit-invariance says that
\((x,y)\in R_G\) implies \(F(x)=F(y)\) for every \(F\in\EQbang\), so
\(R_G\) is contained in the common kernel.  Orbit collapse identifies
that common kernel with \(\Delta_X\), while reflexivity gives
\(\Delta_X\subseteq R_G\).
\end{proof}

\begin{lemma}[Holonomy and gluing-class invariance are quotient features]
\label{lem:pvnp:holonomy-quotient}
Let \(H\) be a holonomy relation, transport-equivalence relation, or
global gluing-class relation on \(X\).  Suppose a proposed Class-B law
requires every \(F\in\EQbang\) to be invariant under \(H\).  If
\[
  R_H=\operatorname{EqCl}(H),
\]
then \(R_H\) is respected by every \(F\in\EQbang\), and
\[
  R_H=\Delta_X .
\]
\end{lemma}

\begin{proof}[Sketch.]
Equivalence closure preserves invariance: a finite chain of
\(H\)-identifications gives a finite chain of equalities after applying
any \(F\in\EQbang\).  Hence \(R_H\) lies in the common kernel of all
lawful current observables.  The conclusion follows from
\Cref{thm:pvnp:orbit-collapse}, exactly as in
\Cref{lem:pvnp:orbit-quotient}.
\end{proof}

\begin{lemma}[Polynomial pairing closure]
\label{lem:pvnp:polynomial-pairing-closure}
If \(F,G\in\EQbang\), then the paired observable
\[
  x\longmapsto (F(x),G(x))
\]
belongs to \(\EQbang\).  More generally, every finite product of
observables in \(\EQbang\) is again an observable in \(\EQbang\).
\end{lemma}

\begin{proof}[Sketch.]
This is the finite-product clause of
\Cref{thm:pvnp:e-q-closure-properties}, restated in the Class-B
framework as a named closure lemma.
\end{proof}

\begin{lemma}[Polynomial post-processing closure]
\label{lem:pvnp:polynomial-postprocessing-closure}
Let \(F\in\EQbang\), and let \(g\) be a polynomial-time
post-processing function on the output of \(F\).  Then
\[
  g\circ F\in\EQbang .
\]
More generally, polynomial-time post-processing of a finite product of
\(\EQbang\)-members is again in \(\EQbang\).
\end{lemma}

\begin{proof}[Sketch.]
This is the post-processing clause of
\Cref{thm:pvnp:e-q-closure-properties}.  The proof in
\Cref{sec:observable-saturation} records the time-certificate and
audit data needed to keep the post-processed observable lawful.
\end{proof}

\begin{lemma}[Presentation refinement]
\label{lem:pvnp:presentation-refinement}
Let \(P\) be a fixed-presentation readout, such as a bit projection,
first-variable readout, encoding hash, or the identity observable.
Let \(R\) be an equivalence relation that a proposed obstruction
requires to be respected by every \(F\in\EQbang\), and define
\[
  R^P = R\cap \ker(P).
\]
Then \(R^P\) is an equivalence relation finer than \(R\).  If \(R\) is
respected by every \(F\in\EQbang\), then so is \(R^P\).  If \(P\)
separates an \(R\)-class, meaning that \(x\,R\,y\) but
\(P(x)\neq P(y)\) for some \(x,y\), then \(R^P\subsetneq R\).  Under
full fixed-presentation saturation, iterating such refinements yields
\(\Delta_X\).
\end{lemma}

\begin{proof}[Sketch.]
The kernel of a readout is an equivalence relation, so \(R^P\) is an
intersection of equivalence relations.  The inclusion \(R^P\subseteq R\)
gives refinement, and strictness follows when the readout separates a
pair in an \(R\)-class.  Since fixed-presentation saturation includes
the lawful presentation readouts, repeated refinement reaches the
common kernel of \(\EQbang\), which is \(\Delta_X\) by
\Cref{thm:pvnp:orbit-collapse}.
\end{proof}

\begin{lemma}[Time-stratum refinement]
\label{lem:pvnp:time-stratum}
Let \(E^{[k]}\) denote the family of observables with time bound
\(O(n^k)\), with
\[
  \EQbang=\bigcup_k E^{[k]} .
\]
A time-stratum feature has two possible statuses.  At a fixed
stratum, an exclusion \(T\notin E^{[k]}\) does not imply
\(T\notin\EQbang\), since \(T\) may lie in \(E^{[\ell]}\) for some
\(\ell>k\).  Uniformly over all strata,
\[
  \bigl(\forall k,\ T\notin E^{[k]}\bigr)
  \quad\Longleftrightarrow\quad
  T\notin\EQbang .
\]
For \(T=\mathrm{Br}^{\mathrm{lex}}_{\mathrm{SAT}}\), the uniform
statement is exactly the sharpened recognition residual.
\end{lemma}

\begin{proof}[Sketch.]
The fixed-\(k\) case is the observation that membership in the
polynomial union is not determined by membership in one stratum.  The
uniform case is the De Morgan dual of
\(\EQbang=\bigcup_k E^{[k]}\):
\(\bigwedge_k \neg(T\in E^{[k]})\) is equivalent to
\(\neg\bigvee_k(T\in E^{[k]})\).
\end{proof}

The mixed-class combinations \(Q\times Q\) through \(T\times T\)
reduce to the same five normal classes \(Q,J,C,P,T\).  The reduction is
case-by-case: quotient combinations use
\Cref{lem:pvnp:quotient-meet-join}, orbit and holonomy combinations
use \Cref{lem:pvnp:orbit-quotient,lem:pvnp:holonomy-quotient}, joint
and positive-closure combinations use the pairing and post-processing
lemmas, presentation combinations use
\Cref{lem:pvnp:presentation-refinement}, and resource combinations use
\Cref{lem:pvnp:time-stratum}.  The proof of
\Cref{thm:pvnp:basis-combination-closure} invokes these lemmas at the
corresponding per-constructor cases of its case analysis over the
feature-expression type.

Five of the seven lemmas are mechanized in Lean as axiom-free
derivations under the module
\texttt{SixBirdsHiddenness.PvNP.ClassB.Lemmas}:
\Cref{lem:pvnp:quotient-meet-join,lem:pvnp:orbit-quotient,lem:pvnp:holonomy-quotient,lem:pvnp:presentation-refinement,lem:pvnp:time-stratum}.
The remaining two,
\Cref{lem:pvnp:polynomial-pairing-closure,lem:pvnp:polynomial-postprocessing-closure},
are restatements of clauses of
\Cref{thm:pvnp:e-q-closure-properties}; they are mechanized as the
finite-product and post-processing clauses of the Lean theorem
\texttt{SixBirdsHiddenness.PvNP.ObservableSaturation.eQClosureProperties}
rather than as separately named lemmas in the \texttt{ClassB.Lemmas}
module.  The coverage table in \Cref{app:formalization} records the
five \texttt{ClassB.Lemmas} entries and the closure-properties theorem
under the substantive ``Lean proves'' wording.
 
\clearpage
\bibliographystyle{plainnat}
\bibliography{references}

\end{document}
